\documentclass[11pt,twoside,openany]{book}
\newcommand{\FgaLang}{spanish}
\usepackage[a4paper,inner=31mm,outer=27mm,top=28mm,bottom=31mm,marginparwidth=17mm]{geometry}
\usepackage{fontspec}
\defaultfontfeatures{Ligatures=TeX,Scale=MatchLowercase}
\setmainfont{TeX Gyre Pagella}
\setsansfont{TeX Gyre Heros}
\setmonofont{TeX Gyre Cursor}
\usepackage{polyglossia}
\setdefaultlanguage{\FgaLang}
\usepackage{microtype}
\usepackage{graphicx}
\usepackage{mathtools}
\usepackage{amsthm}
\usepackage{unicode-math}
\setmathfont{Latin Modern Math}
\usepackage{bookmark}
\usepackage{hyperref}
\hypersetup{hidelinks,unicode=true}

\setcounter{secnumdepth}{3}
\setcounter{tocdepth}{2}
\numberwithin{equation}{section}

\newcommand{\src}[1]{%
  \phantomsection\hypertarget{src:#1}{}%
  \marginpar{\raggedleft\footnotesize\sffamily #1}%
}
\newcommand{\sid}[1]{\phantomsection\hypertarget{fga:#1}{}}
\newcommand{\Osh}{\underline{O}}
\newcommand{\Zsh}{\underline{\mathbf Z}}
\newcommand{\Ccat}{\mathcal C}
\newcommand{\Homsh}{\underline{\operatorname{Hom}}}
\newcommand{\Extsh}{\underline{\operatorname{Ext}}}
\newcommand{\hsh}{\underline{h}}
\newcommand{\fquote}[1]{\textquotedbl#1\textquotedbl}
\newcommand{\seminar}[1]{%
  \begin{flushleft}\sffamily SEMINARIO BOURBAKI\\#1\end{flushleft}%
}
\newcommand{\seminaren}[1]{%
  \begin{flushleft}\sffamily BOURBAKI SEMINAR\\#1\end{flushleft}%
}
\newcommand{\paperhead}[2]{%
  \begin{center}
    {\Large\bfseries #1\par}
    \vspace{0.6\baselineskip}
    #2
  \end{center}
  \vspace{0.8\baselineskip}
}

\theoremstyle{plain}
\newtheorem{proposition}{Proposición}[section]
\newtheorem{theoreme}[proposition]{Teorema}
\newtheorem{theorem}[proposition]{Theorem}
\newtheorem{corollaire}[proposition]{Corolario}
\newtheorem{corollary}[proposition]{Corollary}
\newtheorem{lemme}[proposition]{Lema}
\newtheorem{lemma}[proposition]{Lemma}

\theoremstyle{definition}
\newtheorem{definition}[proposition]{Definición}
\newtheorem{definitionen}[proposition]{Definition}

\theoremstyle{remark}
\newtheorem*{remarque}{Observación}
\newtheorem*{remark}{Remark}

\newcommand{\unit}[1]{%
  \chapter{#1}%
  \markboth{#1}{#1}%
}
\newcommand{\fsection}[2]{%
  \section*{#1. #2}%
  \addcontentsline{toc}{section}{#1. #2}%
}
\newcommand{\result}[2]{%
  \par\medskip\noindent\textbf{#1%
  \if\relax\detokenize{#2}\relax\else\space#2\fi. ---}\enspace
}

% Printed page numbers recur across different works and errata.
% Namespace only the invisible PDF destinations; preserve all body bytes.
\newcommand{\FgaUnitKey}{frontmatter}
\renewcommand{\src}[1]{%
  \phantomsection\hypertarget{src:\FgaUnitKey:#1}{}%
  \marginpar{\raggedleft\footnotesize\sffamily #1}%
}
\setlength{\emergencystretch}{3em}
% The only inherited boldsymbol use is the Greek mu of the group scheme.
% Select its Unicode bold-italic glyph instead of a nonexistent bold math font.
\renewcommand{\boldsymbol}[1]{\symbfit{#1}}
\newcommand{\seminares}[1]{\begin{flushleft}\sffamily SEMINARIO BOURBAKI\\#1\end{flushleft}}
\theoremstyle{plain}
\newtheorem{teorema}[proposition]{Teorema}
\newtheorem{corolario}[proposition]{Corolario}
\newtheorem{lema}[proposition]{Lema}
\theoremstyle{definition}
\newtheorem{definicion}[proposition]{Definición}
\theoremstyle{remark}
\newtheorem*{observacion}{Observación}
\hypersetup{pdftitle={Fundamentos de la geometría algebraica -- Edición española actual -- tramo B},pdfauthor={Alexandre Grothendieck},pdfsubject={Exposiciones 212, 221, 232 y 236, erratas y comentarios}}
\begin{document}
% BEGIN RIGHTS AND AI TRANSLATION NOTICE
\begingroup
\small
\noindent\textbf{Fuentes y autoría}

La obra original, Fondements de la géométrie algébrique, es de Alexandre Grothendieck. La traducción humana al inglés de \href{https://github.com/thosgood/fga/tree/9976dd5f8e2ddcb388f0a07f3aa334f120d1019f}{Tim Hosgood} constituye una fuente de referencia de esta edición. Su proyecto reconoce las correcciones de Steve Hnizdur y el apoyo técnico de Jon Sterling. La versión 9976dd5f8e2ddcb388f0a07f3aa334f120d1019f identifica el texto inglés utilizado para la comparación de procedencia; no establece la versión histórica exacta incorporada inicialmente. La traducción española y las correcciones y composición posteriores fueron producidas por IA. La autoría de la traducción inglesa permanece con su traductor; no se le atribuye la traducción española. Esta atribución no constituye permiso ni aval.

\noindent\textbf{Derechos, disponibilidad y autoría}

Los derechos sobre las obras originales y sobre las contribuciones protegidas a estas traducciones y transcripciones derivadas corresponden a sus respectivos titulares. Este proyecto no dedica esas obras o contribuciones al dominio público ni concede una licencia sobre ellas. Las menciones de derechos de terceros y las autorizaciones ya concedidas siguen siendo aplicables.

Estas ediciones se ofrecen gratuitamente, con fines no comerciales y exclusivamente educativos, bajo la presunción de que los respectivos titulares de derechos no se opondrían a que estas traducciones y transcripciones estuvieran disponibles gratuitamente junto a los originales franceses en NUMDAM (www.numdam.org/). Esta presunción no significa que se haya obtenido autorización; el acceso gratuito a los originales no autoriza por sí solo su traducción, adaptación o republicación.

Si usted es titular de derechos sobre algún material presentado aquí y se opone a su disponibilidad, puede abrir una incidencia en GitHub (github.com/KokunoYumeto/ega-fga-es/issues). El responsable de la publicación retirará el material afectado a petición del titular de los derechos.

La traducción al español está producida por IA. Las correcciones, la edición y la composición también están producidas por IA. El responsable de la publicación no tradujo el texto. Se conservan la autoría y las atribuciones de las obras originales y de las contribuciones heredadas; no se atribuyen a la IA. Los modelos documentados se identifican en cada edición.
\par
\endgroup
\clearpage
% END RIGHTS AND AI TRANSLATION NOTICE

\frontmatter
\begin{titlepage}
\centering
\vspace*{0.13\textheight}
{\Huge\bfseries Fundamentos de la\par geometría algebraica\par}
\vspace{1.5\baselineskip}
{\Large Alexandre Grothendieck\par}
\vfill
{\Large Edición española actual -- tramo B\par}
\vspace{\baselineskip}
{\large Exposiciones 212, 221, 232 y 236, erratas y comentarios\par}
\vspace{1.5\baselineskip}
{\normalsize Edición española\par}
\end{titlepage}
\chapter*{Nota sobre esta edición}
\addcontentsline{toc}{chapter}{Nota sobre esta edición}
Este volumen contiene la traducción española de las exposiciones 212, 221, 232 y 236, sus erratas y los comentarios de 1962.

La edición actual y la traducción según el impreso se ofrecen como textos distintos.

Las cifras marginales remiten a las páginas de las exposiciones originales.

Traducción, correcciones y composición producidas por IA: OpenAI Codex, GPT-6 Astra, esfuerzo Ultra.
La autoría de la obra original corresponde a Alexandre Grothendieck.

El PDF se acompaña de fuentes LaTeX editables y de los archivos necesarios para su compilación.

\tableofcontents
\mainmatter
\def\FgaUnitKey{212}\unit{Exposición 212 -- Preesquemas cociente}
\sid{212}
\src{99}
\seminares{13.er año, 1960/61, n.º 212\\Febrero de 1961}
\paperhead{TÉCNICAS DE CONSTRUCCIÓN Y TEOREMAS DE EXISTENCIA EN GEOMETRÍA ALGEBRAICA\\
III. PREESQUEMAS COCIENTE}%
  {por Alexander GROTHENDIECK}

\subsection*{Introducción}
\addcontentsline{toc}{subsection}{Introducción}

Los problemas tratados en la presente exposición difieren de los considerados en
los dos anteriores en que se intenta representar ciertos funtores covariantes,
y no ya contravariantes, de esquemas variables. Sin embargo, el procedimiento
de paso al cociente es esencial en muchos problemas de construcción en
geometría algebraica, incluidos los de las Exposiciones I y II ([1], [2]). Así, la
cuestión de la \emph{efectividad de un dato de descenso} sobre un
$T$-preesquema $X$, con respecto a un morfismo fielmente plano y cuasicompacto
$T\longrightarrow S$, equivale a la cuestión de la existencia de un cociente
de $X$ (que satisfaga las propiedades razonables examinadas más adelante) por
la relación de equivalencia plana en $X$ definida por el dato de descenso; las
cuestiones planteadas en [1], A, \S~2 c se resolverán sin duda al mismo tiempo
que las cuestiones formuladas en el párrafo 2 de la presente exposición. Del mismo
modo, el \emph{esquema de Picard} de un $S$-esquema $X$ (para la definición,
véase [2], C, \S~3) puede definirse de diversas maneras como cociente de otros
esquemas (de divisores positivos, o de inmersiones en un esquema proyectivo)
por relaciones de equivalencia planas, mientras que la definición y la
construcción de esos esquemas auxiliares son técnicamente más sencillas: son
esquemas de la forma $\underline{\operatorname{Hom}}_S(X,Y)$ y variantes
definidas en [2], C, \S~2, cuya construcción será objeto de la exposición siguiente
(bajo hipótesis adecuadas de proyectividad). Combinando, pues, los resultados
de la presente exposición y de la siguiente se obtiene la construcción de los esquemas
de Picard bajo hipótesis adecuadas.

El problema del paso al cociente en la categoría de los preesquemas presenta
aún varias cuestiones no resueltas. La más importante se menciona en el
párrafo 8. Actualmente es el único obstáculo para la construcción de
\emph{esquemas de módulos sobre los enteros para curvas de género arbitrario,
variedades abelianas polarizadas}, etc. Esto muestra que su solución merece
los esfuerzos de los especialistas en grupos algebraicos.

\src{100}
\subsection*{1. Relaciones de equivalencia y relaciones de equivalencia efectivas}
\addcontentsline{toc}{subsection}{1. Relaciones de equivalencia y relaciones de equivalencia efectivas}

Sean $\mathcal C$ una categoría y $X$ un objeto de $\mathcal C$. Un par de
morfismos
\[
  p_1,p_2:R\rightrightarrows X
\]
se llama \emph{par de equivalencia} en $\mathcal C$, con \emph{destino} $X$
y \emph{fuente} $R$, si, para todo objeto $T$ de $\mathcal C$, las dos
aplicaciones correspondientes
\[
  p_1(T),p_2(T):R(T)\rightrightarrows X(T)
\]
(donde, para todo objeto $Y$ de $\mathcal C$, se pone
$Y(T)=\operatorname{Hom}(T,Y)$) definen una aplicación
\[
  R(T)\longrightarrow X(T)\times X(T)
\]
que induce una biyección de $R(T)$ sobre el grafo de una relación de
equivalencia en el conjunto $X(T)$. Existe una relación de equivalencia
evidente entre los pares de equivalencia con destino $X$; una clase de
equivalencia para esta última relación se llama una \emph{relación de
equivalencia de $\mathcal C$ en $X$}, o simplemente una relación de
equivalencia cuando no pueda producirse confusión.

Si existe $X\times X$, dar una relación de equivalencia en $X$ equivale a dar
un subobjeto $R$ de $X\times X$ tal que, para todo objeto $T$ de $\mathcal C$,
el subconjunto de
\[
  (X\times X)(T)=X(T)\times X(T)
\]
correspondiente a $R(T)$ sea el grafo de una relación de equivalencia en
$X(T)$. Denotando por $p_1$ y $p_2$ los morfismos de $R$ en $X$ inducidos por
las proyecciones $\operatorname{pr}_1$ y $\operatorname{pr}_2$, la condición
anterior afirma que $(p_1,p_2)$ es un par de equivalencia. También pueden
expresarse diagramáticamente en $\mathcal C$ (siempre que existan $X\times X$
y el producto fibrado $(R,p_2)\times_X(R,p_1)$) los axiomas conjuntistas de
una relación de equivalencia para los $R(T)$ en los $X(T)$, conforme al
principio general explicado en [2], A, \S~1. No lo necesitaremos.

Siempre que se tenga un par de morfismos $(p_1,p_2)$ con la misma fuente $R$
y el mismo destino $X$, puede definirse el \emph{conúcleo} del par como el
objeto $Y$ de $\mathcal C$ que representa el funtor covariante en $Z$
\[
  \operatorname{Hom}_{p_1,p_2}(X,Z),
\]
el conjunto de morfismos $u$ de $X$ en $Z$ tales que $up_1=up_2$. Si existe
$Y$, queda determinado salvo isomorfismo único. Lo denotamos por
$X/(p_1,p_2)$ o, por abuso de notación, por $X/R$; esta última notación se usa
principalmente cuando
\src{101}$(p_1,p_2)$ es un par de equivalencia: entonces es habitual
identificar en la notación la relación de equivalencia definida por el par con
$R$. Obsérvese que, si $Y$ se considera como cociente de $X$, en realidad solo
depende de la relación de equivalencia definida por el par $(p_1,p_2)$.

Partamos ahora de un morfismo
\[
  f:X\longrightarrow Y,
\]
que permite considerar $X$ como un «objeto sobre $Y$», y supongamos que existe
el producto fibrado
\[
  R(f)=X\times_Y X
\]
Sean $p_1$ y $p_2$ sus dos proyecciones. Entonces $(p_1,p_2)$ es un par de
equivalencia, llamado par \emph{asociado} al morfismo $f$. Define, por tanto,
una relación de equivalencia, llamada relación asociada a $f$.

Un par de morfismos $(p_1,p_2)$ con destino $X$ y fuente común $R$ se llama
\emph{par de equivalencia efectivo} si
\begin{enumerate}
  \item[(i)] existe el conúcleo $X/(p_1,p_2)=Y$;
  \item[(ii)] existe el producto fibrado $X\times_Y X$;
  \item[(iii)] el morfismo $R\longrightarrow X\times_Y X$ cuyas componentes son
  $p_1$ y $p_2$ es un isomorfismo.
\end{enumerate}
Entonces $(p_1,p_2)$ es efectivamente un par de equivalencia. La relación de
equivalencia que define se llama también \emph{relación de equivalencia
efectiva}.

Un morfismo $f:X\longrightarrow Y$ se llama \emph{epimorfismo efectivo} si
\begin{enumerate}
  \item[(i)] existe el producto fibrado $X\times_Y X=R$;
  \item[(ii)] existe el cociente $X/(p_1,p_2)$, donde $p_1$ y $p_2$ son las dos
  proyecciones de $R$ en $X$;
  \item[(iii)] el morfismo $X/(p_1,p_2)\longrightarrow Y$ inducido por $f$ es un
  isomorfismo.
\end{enumerate}
Entonces $f$ es efectivamente un epimorfismo, e incluso un epimorfismo
estricto (véase [1], A, \S~2.3); el recíproco vale si existe el producto
fibrado $X\times_Y X$. El objeto cociente de $X$ definido por el epimorfismo
$f$ se llama también \emph{cociente efectivo} de $X$.

Las definiciones anteriores implican la siguiente «correspondencia de
Galois»:

\src{102}
\result{PROPOSICIÓN}{1.1} Existe una correspondencia biunívoca que conserva
los órdenes naturales entre el conjunto de relaciones de equivalencia
efectivas $R$ en $X$ y el conjunto de cocientes efectivos $Y$ de $X$: a $R$
le corresponde el cociente efectivo $X/R$, y a $Y$ le corresponde la relación
de equivalencia efectiva definida por la proyección canónica
$X\longrightarrow Y$ (que está definida por el producto fibrado
$X\times_Y X$ provisto de sus dos proyecciones).

En categorías muy buenas (conjuntos, haces de conjuntos, etc.), todo cociente
es efectivo y toda relación de equivalencia es efectiva. Esto deja de ser
cierto en categorías como la categoría de preesquemas sobre un preesquema fijo
$S$, incluso cuando $S$ es el espectro de un cuerpo, y aun cuando uno se
restringe a esquemas finitos sobre $S$. Las cuestiones de efectividad, e
incluso (para los preesquemas no finitos sobre $S$) las cuestiones de
existencia de cocientes, resultan casi siempre delicadas.

\subsection*{2. Ejemplo: preesquemas finitos sobre $S$}
\addcontentsline{toc}{subsection}{2. Ejemplo: preesquemas finitos sobre S}

Sea $\mathcal C$ la categoría de los preesquemas finitos sobre $S$, donde se
supone que $S$ es localmente noetheriano. Es equivalente a la opuesta de la
categoría de haces coherentes de álgebras conmutativas sobre $S$, o, cuando
$S$ es afín de anillo $A$, a la opuesta de la categoría de $A$-álgebras
finitas sobre $A$ (es decir, finitamente generadas como $A$-módulos). Se sigue
inmediatamente que existen en $\mathcal C$ los límites proyectivos finitos y
los límites inductivos finitos. Esto es bien conocido (sin hipótesis alguna
de finitud\ldots) para los primeros. Así, el producto fibrado de preesquemas
$X,Y$ sobre $S$ corresponde al producto tensorial $B\otimes_A C$ de las
álgebras correspondientes; el núcleo de dos morfismos
$X\rightrightarrows Y$, definidos por dos homomorfismos de $A$-álgebras
$u,v:C\rightrightarrows B$, corresponde al cociente de $B$ por el ideal
generado por los elementos $u(c)-v(c)$, etc. Para los límites inductivos
finitos basta considerar, por una parte, las sumas finitas, que corresponden
al producto ordinario de $A$-álgebras, y, por otra, los conúcleos de pares de
morfismos $X\rightrightarrows Y$, que corresponden (como se ve de inmediato)
al subanillo de $C$ formado por los elementos sobre los que coinciden los
homomorfismos $u,v:C\rightrightarrows B$ (este subanillo es finito sobre $A$
por la hipótesis noetheriana). Además, usando la hipótesis noetheriana, puede
demostrarse que los límites inductivos finitos, y en particular los cocientes,
así construidos en la categoría $\mathcal C$ de preesquemas finitos sobre $S$
son en realidad cocientes en la categoría de \emph{todos} los preesquemas.

\src{103}
Como dijimos en [1], existen epimorfismos no efectivos en $\mathcal C$ (o
epimorfismos no estrictos\footnote{El texto impreso omite «no» antes de
«estrictos». En la categoría considerada existen productos fibrados, y el
párrafo anterior establece precisamente la equivalencia entre efectividad y
epimorfía estricta.}; es lo mismo, puesto que existen los productos fibrados).
\emph{No sé si las relaciones de equivalencia son siempre efectivas} cuando
no se impone ninguna hipótesis de planitud. En esta dirección solo he obtenido
resultados positivos muy parciales, indispensables para la demostración del
teorema fundamental de la teoría formal de módulos (véase [2], B, teo. 1).
Señalemos que en el problema planteado es fácil reducirse al caso en que $S$
es el espectro de un anillo local artiniano $A$ con cuerpo residual
algebraicamente cerrado. Pero incluso si $A$ es un cuerpo, no se conoce la
respuesta.

También puede considerarse un preesquema $X$ sobre $S$ que ya no se suponga
finito sobre $S$, pero provisto de una relación de equivalencia $R$ en $X$ tal
que $p_1:R\longrightarrow X$ sea un morfismo finito. Entonces $R$ se llama
una relación de equivalencia finita. Supongamos, para simplificar, que $S$ y
$X$ son afines (y, por tanto, $R$ es afín, de modo que la situación se reduce
a una cuestión de álgebra conmutativa pura). \emph{Incluso en este caso se
ignora si existe un cociente $X/R=Y$ y si el morfismo canónico
$X\longrightarrow Y$ es finito}. (El caso más sencillo es aquel en que $S$ es
el espectro de un cuerpo $k$ y $X$ es el espectro de $k[t]$, es decir, la
recta afín.) Por supuesto, si los dos problemas anteriores tienen soluciones
afirmativas, puede concluirse en la situación presente que $R$ es efectiva.
Obsérvese que el problema de la \emph{existencia} de un cociente $Y$ y de la
finitud de $f:X\longrightarrow Y$ se plantea exactamente en los mismos
términos si, en lugar del grafo de una relación de equivalencia en $X$, solo
se tiene un pregrafo de equivalencia en $X$ en el sentido del párrafo 4.

La cuestión del paso al cociente por una relación de equivalencia finita más
o menos arbitraria se presenta al construir preesquemas «pegando» preesquemas
dados $X_i$ a lo largo de ciertos subpreesquemas cerrados; la ley de pegado se
expresa precisamente mediante una relación de equivalencia finita en el
preesquema suma $X$ de los $X_i$. Cabe esperar también que la solución de los
problemas aquí planteados, y de diversas variantes, sea una condición previa
para desarrollar una técnica general de construcciones no proyectivas en la
dirección iniciada en [2].

El único hecho positivo general conocido por el autor es el siguiente:

\result{PROPOSICIÓN}{2.1} Sea $S$ un preesquema localmente noetheriano, sea
$s$ un punto de $S$ y sea $\Omega$ una extensión algebraicamente cerrada de
$k(s)$. Considérese el
\src{104}«funtor fibra» correspondiente $F$, que asigna a todo $S$-esquema
$X$ finito sobre $S$ el conjunto de puntos con valores en $\Omega$ de $X_s$.
Este funtor (trivialmente exacto por la izquierda) es \emph{exacto por la
derecha}, es decir, conmuta con los límites inductivos finitos y, en
particular, con los conúcleos de pares de morfismos.

Aplicando este resultado a todos los «puntos geométricos» de $S$, se deduce
que la categoría «cociente» $\mathcal C'$ de $\mathcal C$, obtenida razonando
«módulo morfismos radiciales sobreyectivos» (es decir, adjuntando formalmente
inversos a los morfismos en cuestión), es una categoría «geométrica», esto es,
satisface las mismas propiedades «de naturaleza finita» que la categoría de
conjuntos. En particular, toda relación de equivalencia en ella es efectiva.
Esto implica que, si $R$ es una relación de equivalencia en un $X$ finito
sobre $S$, entonces el morfismo canónico
\[
  R\longrightarrow X\times_Y X\qquad (\text{donde }Y=X/R)
\]
es \emph{radicial y sobreyectivo} (de hecho, una inmersión cerrada
sobreyectiva, puesto que es un monomorfismo).

\subsection*{3. El caso de un grupo de operadores}
\addcontentsline{toc}{subsection}{3. El caso de un grupo de operadores}

Supongamos de nuevo que $\mathcal C$ es una categoría arbitraria. Sean $G,X$
objetos de $\mathcal C$, y supongamos que $G$ es un grupo de operadores de
$\mathcal C$ sobre el objeto $X$ de $\mathcal C$. Esto significa (véase [2],
A, \S~1) que, para todo objeto $T$ de $\mathcal C$, se han dado una estructura
de grupo en $G(T)$ y una estructura de $G(T)$-conjunto en $X(T)$, de tal modo
que, al variar $T$, las estructuras en cuestión «varían funtorialmente» con
$T$. Si existen los productos $G\times G$ y $G\times X$ en $\mathcal C$, una
estructura tal puede definirse también mediante un par de morfismos
\[
  G\times G\longrightarrow G,
  \qquad \pi:G\times X\longrightarrow X,
\]
sometidos a la condición de que, para todo objeto $T$ de $\mathcal C$, las
leyes de composición correspondientes en los conjuntos $G(T)$ y $X(T)$ hagan
de $G(T)$ un grupo que actúa en $X(T)$. Traducir este axioma en la
conmutatividad de ciertos diagramas en $\mathcal C$ es fácil pero tedioso y,
de hecho, completamente inútil en todos los casos que conozco.

Supongamos que existe $G\times X$ y consideremos los dos morfismos
\[
  p_1,p_2:G\times X\rightrightarrows X
\]
dados por\footnote{El texto impreso da \texttt{pr\_1}. Puesto que los dos
morfismos tienen codominio $X$, el primero de ellos debe ser
\texttt{pr\_2}; se corrige aquí la errata tipográfica.}
\src{105}
\[
  p_1=\operatorname{pr}_2,\qquad p_2=\pi.
\]
Se comprueba de inmediato que $(p_1,p_2)$ es un par de equivalencia si y solo
si, para todo objeto $T$ de $\mathcal C$, la aplicación
\[
  G(T)\times X(T)\simeq(G\times X)(T)
  \longrightarrow X(T)\times X(T),
\]
definida por este par es inyectiva, es decir, si el grupo $G(T)$ actúa
\emph{libremente} en el conjunto $X(T)$, esto es, si $g\in G(T)$,
$x\in X(T)$ y $g\cdot x=x$ implican que $g$ es el elemento identidad de
$G(T)$. Se dice entonces que $G$ \emph{actúa libremente en $X$} (o que $X$ es
un \emph{espacio principal de $\mathcal C$ bajo $G$}). La relación de
equivalencia asociada al par $(p_1,p_2)$ se llama entonces \emph{relación de
equivalencia definida por el grupo $G$ que actúa libremente en $X$}. Si existe
también $X\times X$ y se considera el morfismo
\[
  p:G\times X\longrightarrow X\times X
\]
definido por el par $(p_1,p_2)$, la condición de que $G$ actúe libremente
significa que $p$ es un \emph{monomorfismo}.

Por supuesto, incluso cuando $G$ no actúa libremente en $X$, se desean
criterios para la existencia de un cociente de $X$ por $G$, es decir, del
conúcleo del par anterior $(p_1,p_2)$.

El conúcleo en cuestión se denotará a menudo por $X/G$, o preferiblemente por
$G\backslash X$ si $G$ actúa por la izquierda (la primera notación se reserva
para el caso en que $G$ actúa por la derecha). Obsérvese que, incluso cuando
existe la «imagen» de $G\times X$ por $p$ (definida, por ejemplo, como el
menor subobjeto de $X\times X$ a través del cual se factoriza $p$), digamos
$R$, esta no es casi nunca una relación de equivalencia en $X$. Si se intenta
entonces pasar directamente al cociente por $R$ (o, más precisamente, por el
par de morfismos de $R$ en $X$ inducidos por las dos proyecciones
$\operatorname{pr}_i$), se pierden irremediablemente las características
especiales del par original $(p_1,p_2)$. Por ello es importante hallar una
generalización de la noción de relación de equivalencia que se aplique
directamente al par definido por un grupo de operadores de $\mathcal C$.

\subsection*{4. Prerrelaciones de equivalencia}
\addcontentsline{toc}{subsection}{4. Prerrelaciones de equivalencia}

Recordemos que una categoría en la que todo morfismo es un isomorfismo se
llama un \emph{grupoide}. Por otra parte, una categoría debe definirse como
constituida por dos conjuntos subyacentes $(X,R)$, el conjunto de
\emph{objetos} y el conjunto de \emph{flechas}, provistos de las estructuras
siguientes:
\src{106}
\begin{enumerate}
\item[(i)] Un par de aplicaciones
\[
  p_1,p_2:R\rightrightarrows X
\]
llamadas \emph{aplicación fuente} y \emph{aplicación destino};

\item[(ii)] Una aplicación
\[
  \pi:(R,p_2)\times_X(R,p_1)\longrightarrow R
\]
llamada \emph{aplicación de composición}.
\end{enumerate}
Estos datos deben satisfacer axiomas conocidos, que no repetiremos aquí y que
podrían expresarse mediante la conmutatividad de ciertos diagramas, junto con
la existencia de una aplicación $D:X\to R$ (necesariamente única) que haga
conmutativos otros dos diagramas. La aplicación $D$ envía un objeto a la
flecha identidad correspondiente y satisface
\[
  p_1D=p_2D=\operatorname{id}_X.
\]
Decir que la categoría es un grupoide significa entonces que existe una
aplicación
\[
  s:R\longrightarrow R
\]
(necesariamente única), llamada la \emph{inversión} de $R$, que transforma
cada flecha en una flecha inversa. Esto podría expresarse mediante la
conmutatividad de otros cuatro diagramas formados a partir de $s$, $\Delta$ y
los datos anteriores; los dos primeros son
\[
  p_1s=p_2,\qquad p_2s=p_1.
\]

Recordados estos puntos, las definiciones generales de ([2], A, \S~1)
muestran, en particular, qué se entiende por una estructura de
$\mathcal C$-\emph{categoría}, respectivamente de
$\mathcal C$-\emph{grupoide}, sobre un par de objetos $X,R$ de una categoría
arbitraria $\mathcal C$: por definición, es el dato, para todo objeto $T$ de
$\mathcal C$, de una estructura de categoría, respectivamente de grupoide, en
el sentido conjuntista, cuyo conjunto de objetos es $X(T)$ y cuyo conjunto de
flechas es $R(T)$, «variando funtorialmente» estas estructuras con la variable
$T$. Esto proporciona, por tanto, dos morfismos
\[
  p_1,p_2:R\rightrightarrows X
\]
llamados \emph{morfismo fuente} y \emph{morfismo destino}, y, cuando existe el
producto fibrado requerido, un morfismo
\src{107}
\[
  \pi:(R,p_2)\times_X(R,p_1)\longrightarrow R
\]
llamado \emph{morfismo de composición}. Estos tres datos bastan entonces para
determinar la estructura de categoría sobre $X,R$, y el axioma que se les
impone es el siguiente: para todo $T$, los tres datos correspondientes para
$X(T),R(T)$ definen una estructura de categoría (respectivamente de grupoide)
sobre este par de conjuntos. Cuando proceda, esto puede expresarse mediante
la conmutatividad de ciertos diagramas en los que interviene un morfismo
determinado de manera única
\[
  D:X\longrightarrow R
\]
y, en el caso de los grupoides, un morfismo determinado de manera única
\[
  s:R\longrightarrow R;
\]
los diagramas se escriben exactamente como en el caso «conjuntista». Esta
laboriosa interpretación de los axiomas es, afortunadamente, inútil en la
práctica. El único interés teórico de poder expresar los datos y axiomas en
términos de morfismos e igualdades de morfismos entre ciertos productos
fibrados es el siguiente: si $F:\mathcal C\to\mathcal C'$ es un funtor exacto
por la izquierda (es decir, que conmuta con productos finitos y productos
fibrados), entonces transforma una $\mathcal C$-categoría (respectivamente un
$\mathcal C$-grupoide) en una $\mathcal C'$-categoría (respectivamente un
$\mathcal C'$-grupoide), siempre que existan productos finitos y productos
fibrados en $\mathcal C$.

En la práctica, es importante saber interpretar los morfismos
$p_1,p_2,\pi,D,s$ como \emph{operaciones simpliciales} en un objeto
semisimplicial o simplicial adecuado de $\mathcal C$, al menos cuando existen
productos fibrados en $\mathcal C$. Para fijar la terminología, introduzcamos
la \emph{categoría $\mathcal S$ de los símplices estándar} como la categoría
cuyos objetos son los conjuntos finitos de la forma
\[
  \Delta_n=[0,n],\qquad\text{donde }n\geq 0
\]
(el intervalo de enteros de $0$ a $n$), y cuyos morfismos o flechas son
\emph{aplicaciones arbitrarias} entre esos conjuntos finitos. Obsérvese que la
categoría $\mathcal S$ es equivalente a la categoría de conjuntos finitos
\emph{no vacíos}, con las aplicaciones de conjuntos finitos como morfismos.
En $\mathcal S$ existe claramente el coproducto de una familia finita
\emph{no vacía} de objetos, así como la suma amalgamada de dos objetos bajo un tercero
(la operación dual del producto fibrado). Denotemos por $\mathcal S'$ la
subcategoría de $\mathcal S$ que tiene los mismos objetos pero cuyos morfismos
son las \emph{aplicaciones crecientes} entre los $\Delta_n$. Esta categoría es
equivalente a la categoría de conjuntos finitos no vacíos totalmente
ordenados. En esta categoría el
\src{108}
coproducto de dos objetos nunca existe, y la suma amalgamada de dos objetos $A,B$ bajo
un tercero $C$ no existe en general (tómese, por ejemplo, $C=\Delta_0$,
$A=B=\Delta_1$, con las dos aplicaciones estructurales $u:C\to A$ y
$v:C\to B$ iguales). Sin embargo, existe en ciertos casos; por ejemplo,
\[
  C=\Delta_0,\qquad A=\Delta_m,\qquad B=\Delta_n,\qquad
  u(0)=m,\qquad v(0)=0,
\]
y en este caso
\[
  \Delta_m\amalg_{\Delta_0}\Delta_n=\Delta_{m+n}.
\]

Un \emph{objeto simplicial} (respectivamente, un \emph{objeto
semisimplicial}) en una categoría $\mathcal C$ es, por definición, un funtor
contravariante $K$ de $\mathcal S$ (respectivamente de $\mathcal S'$) en
$\mathcal C$. Por tanto, un objeto simplicial define por restricción un objeto
semisimplicial, pero la primera noción difiere esencialmente de la segunda por
la presencia de \emph{operaciones de simetría} en los $K_n=K(\Delta_n)$,
correspondientes a las imágenes por el funtor $K$ de los elementos del grupo
simétrico de $n+1$ letras (considerado como el grupo de automorfismos de
$\Delta_n$ en $\mathcal S$).

Para cada $n$, sea $\Delta'_n$ (respectivamente $\Delta''_n$) la categoría
finita cuyo conjunto de objetos es $\Delta_n$, definida por la relación de
orden caótica (respectivamente por el orden total natural) en $\Delta_n$ (el
conjunto de flechas es el grafo de esa relación de orden). Está claro que
$\Delta'_n$ (respectivamente $\Delta''_n$) depende funtorialmente del objeto
$\Delta_n$ de $\mathcal S$ (respectivamente de $\mathcal S'$). Si $Z$ es una
categoría, entonces $\operatorname{Hom}(\Delta'_n,Z)$ (respectivamente
$\operatorname{Hom}(\Delta''_n,Z)$), al variar $\Delta_n$, es un funtor de
$\mathcal S$ (respectivamente de $\mathcal S'$) a la categoría de conjuntos,
es decir, un conjunto simplicial (respectivamente un conjunto
semisimplicial), que se dice \emph{asociado a la categoría $Z$}; denotémoslos
por $Z'$ y $Z''$. Hay, además, un homomorfismo natural evidente del conjunto
semisimplicial asociado a $Z'$ hacia $Z''$, y este último es un isomorfismo si
y solo si $Z$ es un grupoide. Entendido esto:

\result{PROPOSICIÓN}{4.1} El funtor $Z\longmapsto Z''$ de la categoría de
categorías a la categoría de conjuntos semisimpliciales es plenamente fiel y
define una equivalencia entre la categoría de las \emph{categorías} y la
categoría de los conjuntos semisimpliciales, es decir, los funtores
contravariantes $K$ de $\mathcal S'$ en $(\mathrm{Set})$ que transforman las
sumas amalgamadas $A\amalg_C B$ del tipo especificado anteriormente en productos
fibrados de conjuntos. Asimismo, el funtor $Z\longmapsto Z'$ de la categoría
de grupoides a la categoría de conjuntos simpliciales es plenamente fiel y
define una equivalencia entre la categoría de los \emph{grupoides} y la
categoría de los
\src{109}
conjuntos simpliciales, es decir, los funtores contravariantes $K$ de
$\mathcal S$ en $(\mathrm{Set})$ que transforman las sumas amalgamadas en productos
fibrados.

Así, las categorías pueden considerarse como conjuntos semisimpliciales
particulares, y los grupoides como conjuntos simpliciales particulares,
siempre que, por supuesto, se razone «salvo isomorfismo», como exige la
interpretación de ciertas estructuras en términos de otras. El procedimiento
usual de reducción al caso conjuntista implica entonces:

\result{COROLARIO}{4.2} El enunciado anterior sigue siendo válido cuando las
categorías, los grupoides y los conjuntos simpliciales se sustituyen por
$\mathcal C$-categorías, $\mathcal C$-grupoides y objetos simpliciales en
$\mathcal C$, siempre que existan productos fibrados en $\mathcal C$.

El objeto semisimplicial $K$ en $\mathcal C$ asociado a una categoría
$(X,R,\ldots)$ en $\mathcal C$ puede hacerse explícito interpretando la
componente $K_n=K(\Delta_n)$ de $K$ como la potencia fibrada $n$-ésima de
$(R,p_1)$ sobre $X$, o, mejor, mediante la fórmula de recurrencia
\[
  K_0=X,\qquad
  K_n=(K_{n-1},p_{n-1}^{(n-1)})\times_X(R,p_1),
\]
donde los $p_i^{(n-1)}$ ($0\leq i\leq n-1$) son las proyecciones naturales de
$K_{n-1}$ en $X$ (que se definen también recursivamente).\footnote{El texto impreso lleva $(n+1)$, $K_0=R$ y $p_n^{(n-1)}$ en estos tres lugares. La indexación del nervio exige $n$, $K_0=X$ y $p_{n-1}^{(n-1)}$, respectivamente.} De este modo,
$p_1,p_2,\pi,D,s$ se interpretan como las operaciones simpliciales
correspondientes respectivamente a los morfismos siguientes de $\mathcal S$:
la cara $0$ de $\Delta_1$, la cara $1$ de $\Delta_1$, la cara $(0,2)$ de
$\Delta_2$, la degeneración $\Delta_1\to\Delta_0$ y la simetría de
$\Delta_1$. Todas las demás operaciones semisimpliciales (respectivamente,
simpliciales) se obtienen formalmente a partir de las cuatro (respectivamente,
cinco) anteriores mediante composición y productos fibrados.

Una \emph{prerrelación de equivalencia} en un objeto $X$ de una categoría se
define ahora como el dato de un grupoide cuyo objeto de objetos (por así
decirlo) es $X$. Un dato tal comprende, por tanto, entre otras cosas, un
objeto $R$ y dos morfismos
\[
  p_1,p_2:R\rightrightarrows X.
\]
Pero obsérvese que estos datos por sí solos no determinan la estructura en
cuestión, a diferencia de lo que ocurre con los pares de equivalencia. En esta
exposición nos interesa esta noción para obtener criterios de posibilidad del
paso al cociente, es decir, para formar el conúcleo del par $(p_1,p_2)$. Por
ello, el enunciado del problema no hace intervenir los datos adicionales
inherentes a un grupoide. En las demostraciones de los resultados que siguen,
\src{110}
usaremos, no obstante, esos datos adicionales y, en particular, las
operaciones simpliciales (incluidas las operaciones de simetría) hasta la
dimensión $3$ (en las que interviene la cuarta potencia fibrada de $R$ sobre
$X$).

Una relación de equivalencia en un objeto $X$ de $\mathcal C$ define una
prerrelación de equivalencia: basta verlo en el caso conjuntista, donde a una
relación de equivalencia en un conjunto $X$ se asocia el grupoide cuyo
conjunto de objetos es $X$ y cuyo conjunto de flechas es el grafo de la
relación de equivalencia.

Un monoide de $\mathcal C$ $G$ que actúa sobre un objeto $X$ de $\mathcal C$
define una $\mathcal C$-categoría cuyos objetos subyacentes son $R=G\times X$
y $X$ (siempre que exista $G\times X$), y esta es un $\mathcal C$-grupoide si
y solo si $G$ es un grupo. De nuevo, basta comprobarlo en el caso conjuntista.
La composición de las flechas $(g,a)$ y $(g',g\cdot a)$ se define entonces por
\[
  (g',g\cdot a)(g,a)=(g'g,a),
\]
es decir, si $a,b\in X$, entonces $\operatorname{Hom}(a,b)$ es por definición
el transportador de $a$ a $b$, y los morfismos se componen mediante la
composición de los elementos de $G$.

\medskip
\noindent\textbf{OBSERVACIÓN. ---} Las dificultades lógicas planteadas por un
enunciado como (4.1) pueden evitarse entendiendo que todos los objetos
considerados pertenecen a un «Universo» fijo (que es a su vez un conjunto).

\fsection{5}{Cociente por una relación de equivalencia finita y plana}

\result{TEOREMA}{5.1} Sea $X=\operatorname{Spec}(B)$ un esquema afín y sea
$R$ una prerrelación de equivalencia en $X$ cuya componente $R_1$ sea afín,
digamos $R_1=\operatorname{Spec}(C)$. Supongamos que la primera proyección
$p_1:R_1\to X$ es finita y localmente libre; en otras palabras, supongamos
que el homomorfismo de anillos correspondiente $p_1^*:B\to C$ hace de $C$ un
$B$-módulo proyectivo finitamente generado. Sea $A$ el subanillo de $B$ que
es el núcleo del par de homomorfismos
\[
  p_1^*,p_2^*:B\rightrightarrows C
\]
(el conjunto de elementos $b$ tales que $p_1^*(b)=p_2^*(b)$). Sea
$Y=\operatorname{Spec}(A)$ y sea $f:X\to Y$ el morfismo definido por la
inclusión de $A$ en $B$. En estas condiciones:
\begin{enumerate}
\item[(i)] $B$ es entero sobre $A$, es decir, $f$ es un morfismo entero.

\item[(ii)] El morfismo $f$ es sobreyectivo y sus fibras son las clases de
equivalencia
\src{111}
$p_2(p_1^{-1}(x))$ en el conjunto $X$ módulo $R$; además, la topología de $Y$
es el cociente de la de $X$.

\item[(iii)] $Y$ es el cociente de $X$ por $R$ en la categoría de los
preesquemas.

\item[(iv)] Si $R$ proviene de una \emph{relación de equivalencia}, entonces
el morfismo $f:X\to Y$ es finito localmente libre (es decir, $B$ es un
$A$-módulo proyectivo finitamente generado), y la relación de equivalencia es
efectiva, esto es, $R_1\to X\times_Y X$ es un isomorfismo.
\end{enumerate}

Este teorema generaliza el conocido teorema relativo a un grupo finito $G$ que
actúa por automorfismos en el anillo $B$ y al anillo de invariantes $A$, y su
demostración es análoga a la demostración conocida. El enunciado (iii) puede
precisarse como sigue:

\result{COROLARIO}{5.2} El morfismo canónico
$R_1\to X\times_Y X$ es \emph{sobreyectivo}.

Sea todavía $R$ una prerrelación de equivalencia «finita y localmente libre»
en $X$, pero sea ahora $X$ un preesquema arbitrario. Supongamos que puede
hallarse un preesquema $Y$ y un morfismo $f:X\to Y$ tales que $fp_1=fp_2$ y,
además, tales que la sucesión de homomorfismos de haces de anillos sobre $Y$
\[
  \mathcal O_Y\longrightarrow f_*(\mathcal O_X)
  \rightrightarrows g_*(\mathcal O_R)
\]
sea exacta (donde $g=fp_i$). El teorema proporciona entonces conclusiones
análogas a (i)--(iv) anteriores; en particular, por (iii), $Y$ es el cociente
de $X$ por $R$ y, por tanto, queda determinado salvo isomorfismo único. En
estas condiciones se dice que la prerrelación de equivalencia $R$ en $X$ es
\emph{admisible}. Con esta definición:

\result{TEOREMA}{5.3} Sea $X$ un preesquema y sea $R$ una prerrelación de
equivalencia en $X$ tal que $p_1:R_1\to X$ sea finito y localmente libre.
Entonces $R$ es admisible si y solo si toda clase de equivalencia conjuntista
$p_2(p_1^{-1}(x))$ en $X$ módulo $R$ está contenida en un abierto afín
(condición que siempre se satisface si todo subconjunto finito de $X$ está
contenido en un abierto afín, por ejemplo si $X$ es cuasiproyectivo sobre un
esquema afín).

En efecto, se demuestra sin dificultad que toda clase de equivalencia módulo
$R$ en $X$ está entonces contenida en un abierto afín \emph{estable bajo
$R$}, y se construye el cociente $Y$ pegando las piezas obtenidas del Teorema
5.1.
\src{112}

\result{COROLARIO}{5.4} Supongamos que se satisface esta condición y que,
además, $R$ proviene de una relación de equivalencia. Entonces esa relación es
efectiva, es decir, $R_1\to X\times_Y X$ es un isomorfismo, y $f:X\to Y$ es
finito y localmente libre.

Se concluye inmediatamente por «descenso»:

\result{COROLARIO}{5.5} Bajo las hipótesis de (5.4), $X$ tiene rango $n$ en
todo punto sobre $Y$ si y solo si $(R_1,p_1)$ tiene rango $n$ en todo punto
sobre $X$. Si $X,R_1$ son $Z$-preesquemas y $p_1,p_2$ son $Z$-morfismos, de
modo que $Y$ es un $Z$-preesquema, entonces $X$ es plano sobre $Z$ si y solo
si $Y$ es plano sobre $Z$.

En resumen:

\result{ESCOLIO}{} El dato de un morfismo finito, localmente libre y
sobreyectivo $f:X\to Y$ de preesquemas equivale al dato de un preesquema $X$
provisto de una relación de equivalencia $R$ tal que $p_1:R\to X$ sea finito
y localmente libre y toda clase $p_2(p_1^{-1}(x))$ esté contenida en un
abierto afín.

\medskip
\noindent\textbf{OBSERVACIONES 5.6.}

\smallskip
\noindent $1^\circ$ No se ha necesitado ninguna hipótesis noetheriana.

\smallskip
\noindent $2^\circ$ Esta noción de paso al cociente contiene, como caso
particular, el «descenso inseparable» de CARTIER, que consiste en determinar
los morfismos finitos localmente libres $f:X\to Y$ tales que
$f_*(\mathcal O_X)$ admita una $p$-base relativa a $\mathcal O_Y$ (donde $X$
es un preesquema dado cuyo haz $\mathcal O_X$ es anulado por el número primo
$p>0$). Obsérvese que el resultado se expresa fácilmente sin ninguna
hipótesis de regularidad sobre los anillos locales y sin suponer que $X$ sea
un esquema algebraico sobre un cuerpo. La teoría de JACOBSON--BOURBAKI se
obtiene tomando por $X$ el espectro de un cuerpo de característica $p$.

\smallskip
\noindent $3^\circ$ GABRIEL había obtenido anteriormente un caso particular
del Teorema (5.3) en la teoría del paso al cociente para grupos conmutativos
finitos sobre un cuerpo $k$. [Compárese (7.4).]

\fsection{6}{Cociente por una relación de equivalencia propia y plana}

\result{TEOREMA}{6.1} Sea $S$ un preesquema localmente noetheriano, sea $X$
un $S$-esquema cuasiproyectivo y sea $R$ una prerrelación de equivalencia en
$X$ tal que:
\src{113}
\begin{enumerate}
\item[(a)] $p_1:R_1\to X$ es propio y plano;
\item[(b)] $R_1\to X\times_S X$ es finito o, equivalentemente, por (a), tiene
fibras finitas (condición que se satisface automáticamente si $R$ proviene de
una relación de equivalencia).
\end{enumerate}
En estas condiciones:
\begin{enumerate}
\item[(i)] Existe $X/R=Y$ y, si $S$ es noetheriano y $R$ proviene de una
relación de equivalencia, es cuasiproyectivo sobre $S$.\footnote{El texto
impreso no contiene la segunda hipótesis. La fe de erratas vinculada de la Exposición
212 dice que la cuasiproyectividad parece haberse demostrado entonces solo
cuando $R$ proviene de una relación de equivalencia.}

\item[(ii)] El morfismo canónico $f:X\to Y$ es sobreyectivo, propio y abierto;
sus fibras son las clases de equivalencia $p_2(p_1^{-1}(x))$ en $X$ módulo
$R$. Así, $Y$ se identifica con el espacio topológico cociente de $X$ por la
relación de equivalencia conjuntista definida por $R$. Finalmente,
$R_1\to X\times_Y X$ es sobreyectivo.

\item[(iii)] Si $R$ proviene de una relación de equivalencia, entonces esa
relación es efectiva, es decir, $R_1\to X\times_Y X$ es un isomorfismo, y
además $f:X\to Y$ es plano (por tanto, fielmente plano).
\end{enumerate}

Para la demostración, uno se reduce a (5.1) mediante cuasisecciones adecuadas
de $X$ para $R$; la demostración es análoga a la construcción de grupos
algebraicos cociente en el Seminario Chevalley.

En resumen:

\result{ESCOLIO}{} Sea $X$ cuasiproyectivo sobre un $S$ localmente
noetheriano. El dato de un morfismo propio, fielmente plano y sobreyectivo
$f:X\to Y$ de $X$ a un $S$-preesquema $Y$ equivale al dato de una relación de
equivalencia $R$ en $X$ tal que $p_1:R\to X$ sea propio y plano.

El mismo método da el resultado siguiente:

\result{TEOREMA}{6.2} Sea $S$ un preesquema noetheriano, sea $X$ un
preesquema de tipo finito sobre $S$ y sea $R$ una prerrelación de equivalencia
en el $S$-preesquema $X$. Supongamos que:
\begin{enumerate}
\item[(a)] $p_1:R_1\to X$ es plano y de tipo finito;
\item[(b)] el morfismo $R_1\to X\times_S X$ es cuasifinito (es decir, tiene
fibras finitas).
\end{enumerate}
Entonces existe un abierto denso $U$ de $X$ que está \emph{saturado} para $R$
y tal que:
\begin{enumerate}
\item[(i)] Si $R_U$ es la prerrelación de equivalencia inducida por $R$ en
$U$, entonces existe $U/R_U$ y es de tipo finito sobre $S$.

\item[(ii)] El morfismo canónico $U\to U/R_U$ es sobreyectivo y abierto; sus
fibras son las clases de equivalencia conjuntistas para $R_U$ (así,
$U/R_U$ es el espacio topológico cociente de $U$ por la relación de
equivalencia conjuntista definida por $R_U$). Finalmente, el morfismo
\[
  (R_U)_1\longrightarrow U\times_{U/R_U}U
\]
es sobreyectivo.
\src{114}

\item[(iii)] Si $R$ proviene de una relación de equivalencia, puede disponerse
que $U\to U/R_U$ sea fielmente plano y que $R_U$ sea efectiva.
\end{enumerate}

Este es un resultado de naturaleza esencialmente «birracional».

\medskip
\noindent\textbf{OBSERVACIONES 6.3.}

\smallskip
\noindent $1^\circ$ No sé si la hipótesis (b) es innecesaria en (6.1) y
(6.2). Al tomar cocientes por grupos, nos obliga en la práctica a
restringirnos al caso en que todos los estabilizadores son grupos finitos.

\smallskip
\noindent $2^\circ$ Puede preguntarse si existen resultados análogos a (6.1)
y (6.2) sin hipótesis de planitud. No conozco ningún contraejemplo en esta
dirección. Por otra parte, incluso conservando la hipótesis de planitud y
restringiéndose a relaciones de equivalencia para las cuales $p_1:R\to X$ es
plano y cuasifinito (pero no finito), con $X$ afín, puede ocurrir que $R$ no
sea efectiva: tómense las relaciones de equivalencia inducidas en los
subconjuntos abiertos afines que recubren la variedad de Nagata (en la que un
grupo de orden dos actúa de manera «no admisible»).

\fsection{7}{Aplicaciones}

Como dijimos en la introducción, la aplicación más importante de (6.1) es la
construcción de esquemas de Picard, así como la solución de otros diversos
problemas de módulos, sobre los cuales volveremos más adelante.

Se obtiene una demostración sencilla del resultado siguiente de SHIMURA:

\result{PROPOSICIÓN}{7.1} Sea $A$ un esquema abeliano definido sobre un
anillo de valoración discreta $V$ de cuerpo de fracciones $K$. Entonces todo
esquema abeliano $B'$ sobre $K$ que sea isógeno a un cociente de
$A\otimes_V K$ «tiene buena reducción sobre $V$», es decir, es isomorfo a
algún $B\otimes_V K$,\footnote{El texto impreso contiene $B\otimes_V A$.
Como $B$ es un $V$-esquema y la afirmación concierne a su fibra genérica, el
factor de base debe ser $K$.} donde $B$ es un esquema abeliano sobre $V$
(esencialmente único, como recordamos).

Podemos suponer que $B'$ es el cociente de $A_K$ por un subesquema en grupos
$C'$. (N.B.: $C'$ no será en general «reducido», es decir, sus anillos locales
tendrán elementos nilpotentes.) Considérese el subesquema cerrado $C$ de $A$
que es la «clausura» de $C'$, esto es, el menor subesquema cerrado de $A$ tal
que $C_K$ contiene a $C'$. Se tiene $C_K=C'$, y el hecho de que $V$ sea un
anillo de valoración discreta implica fácilmente que $C$ es un subesquema en
grupos de $A$ sobre $V$. Puesto que $A$ es propio sobre
$\operatorname{Spec}(V)=S$, también lo es $C$; además, $A$ es incluso
proyectivo sobre $S$. Podemos, por tanto, aplicar (6.1) para construir
$A/C=B$,
\src{115}
que es el $B$ buscado.

Por último, argumentos esencialmente conocidos permiten deducir de (6.2) el
resultado siguiente:

\result{TEOREMA}{7.2} Sea $S$ el espectro de un anillo artiniano, sean $F$ y
$G$ esquemas en grupos de tipo finito sobre $S$, y sea $u:F\to G$ un
homomorfismo de esquemas en grupos sobre $S$. Supongamos que:
\begin{enumerate}
\item[(i)] $F$ es plano sobre $S$;
\item[(ii)] el núcleo de $u$ es finito.
\end{enumerate}
En estas condiciones existe el esquema cociente $G/F$; el morfismo canónico
$G\to G/F$ es sobreyectivo y abierto, y sus fibras son las clases de
equivalencia conjuntistas definidas por las acciones por la derecha de $F$ en
$G$. Finalmente, si $u$ es un monomorfismo, entonces $G\to G/F$ es plano y
el morfismo
\[
  G\times F\longrightarrow G\times_{G/F}G
\]
es un isomorfismo. En otras palabras, $G$ es un espacio homogéneo principal
sobre $G/F$ con grupo estructural (que actúa por la derecha) $F$, o, más
exactamente, $F\times_S(G/F)$ considerado como esquema en grupos sobre $G/F$
(véase [1], B, \S~6).

\result{COROLARIO}{7.3} En estas condiciones, $G$ es plano sobre $S$ si y
solo si $G/F$ es plano sobre $S$. Si se cumple esta condición, la formación
del cociente por $F$ conmuta con toda extensión de la base $S$; y, si $F$ es
un subgrupo normal de $G$, entonces $G/F$ puede dotarse de la estructura de
\emph{grupo cociente} de $G$ por $F$.

La situación es especialmente sencilla cuando $S$ es el espectro de un
cuerpo, pues entonces todo $S$-preesquema es automáticamente plano sobre $S$.
Se obtiene:

\result{COROLARIO}{7.4} Sean $F,G$ esquemas en grupos de tipo finito sobre un
cuerpo $k$, y sea $u:F\to G$ un homomorfismo de $k$-grupos. Entonces $u$ se
factoriza como $F\to F'\to G$, donde $F\to F'$ es el homomorfismo cociente
por el subgrupo cerrado $\operatorname{Ker}u$ de $F$, y $F'\to G$ es un
homomorfismo de grupos que es una inmersión cerrada. Existe el cociente
$G/F=G/F'$. El formalismo usual (del tipo proporcionado por los teoremas de
Noether) es válido entre grupos algebraicos sobre $k$.

Este resultado permite tratar de manera uniforme tanto el paso al cociente
entre grupos algebraicos en el sentido clásico (es decir, los que son
irreducibles sobre $k$ y lisos sobre $k$) como el paso al cociente por
subgrupos «infinitesimales»
\src{116}
considerados por CARTIER. Resulta ventajoso considerar las «hiperálgebras»
introducidas por CARTIER, a raíz de los trabajos de DIEUDONNÉ sobre los grupos
formales, como grupos en la categoría de los esquemas formales sobre $k$ y,
cuando proceda (si corresponden a hiperálgebras de rango finito sobre $k$),
como grupos algebraicos finitos sobre $k$.

\fsection{8}{Una conjetura}

Responde a la necesidad de formar cocientes por el grupo proyectivo que actúa
en ciertos subesquemas de «esquemas de Hilbert» (que sustituyen a las
variedades de Chow en la teoría de esquemas).

Sea $S$ un preesquema y sea $n$ un entero. A todo preesquema $S'$ sobre $S$
se asocia el grupo $\operatorname{Gl}(n,\Gamma(S',\mathcal O_{S'}))$ de
matrices invertibles de tamaño $n$ por $n$ con entradas en el anillo de
secciones de $\mathcal O_{S'}$. Esto da un funtor contravariante en $S'$, que
se comprueba fácilmente que es representable; corresponde, por tanto, a un
esquema en grupos sobre $S$, de hecho afín sobre $S$, denotado por
$\operatorname{Gl}(n)_S$. Su formación es compatible con el cambio de base,
de modo que, en realidad, todo procede de un esquema en grupos sobre
$\mathbf Z$, denotado por $\operatorname{Gl}(n)$. El grupo
$\operatorname{Gl}(1)$, llamado grupo multiplicativo y denotado a menudo por
$\mathbf G_m$, corresponde al funtor
$S\longmapsto\Gamma(S,\mathcal O_S)^*$, el grupo de las «unidades» sobre $S$.
Hay un homomorfismo evidente
$\operatorname{Gl}(1)\to\operatorname{Gl}(n)$, y se construye fácilmente el
grupo cociente, denotado por $\operatorname{GP}(n-1)$ y llamado \emph{grupo
proyectivo de grado $n-1$ sobre $\mathbf Z$}. Representa el funtor que asocia
a $S$ el grupo de secciones del haz
$\widetilde{\operatorname{Gl}(n)}_S/
\widetilde{\operatorname{Gl}(1)}_S$, donde
$\widetilde{\operatorname{Gl}(n)}_S$ denota el haz de gérmenes de secciones de
$\operatorname{Gl}(n)_S$ sobre $S$. (¡Hay que tener cuidado: las secciones de
$\operatorname{GP}(n-1)_S$ sobre $S$ no proceden, en general, de secciones de
$\operatorname{Gl}(n)_S$ sobre $S$!) Obsérvese que puede demostrarse que
$\operatorname{GP}(n-1)$ representa también el funtor
$S\longmapsto\operatorname{Aut}_S(\mathbf P_S^{n-1})$, donde
$\mathbf P_S^{n-1}$ es el espacio proyectivo estándar de dimensión relativa
$n-1$ sobre $S$, al menos para $S$ noetheriano. Así es como interviene en la
teoría de módulos.

Sea $S$ un esquema noetheriano, que si se desea puede suponerse afín, y sea
$X$ un $S$-preesquema cuasiproyectivo provisto de un haz invertible
$\mathcal L$ muy amplio con respecto a $S$. Supongamos que el grupo
$G=\operatorname{GP}(n)_S$ actúa en $X$ y al mismo tiempo en $\mathcal L$
(compatiblemente con su acción en $X$), y que actúa \emph{libremente} en
$X$.\footnote{El texto impreso dice «en $S$». La acción considerada y la
condición de cociente conciernen a $X$.}
\src{117}

\result{CONJETURA}{8.1}\footnote{La fe de erratas vinculada de la Exposición 212
afirma que esta conjetura es definitivamente falsa. También señala que el
hecho positivo de la Observación 8.1 fue demostrado posteriormente y remite
al teorema de Mumford y a refinamientos posteriores.} En las condiciones
anteriores:

\smallskip
\noindent $1^\circ$ \emph{La relación de equivalencia definida por $G$ es
efectiva, el cociente $Y=X/G$ es de tipo finito sobre $S$, y el morfismo
canónico $f:X\to Y$ es plano y sobreyectivo} (así, $X$ se convierte en un
fibrado principal homogéneo sobre $Y$ con grupo
$G\times_S Y=\operatorname{GP}(n)_Y$).

\smallskip
\noindent $2^\circ$ Sea $\mathcal L'$ el haz invertible sobre $Y$ deducido de
$\mathcal L$ por «descenso fielmente plano» a lo largo de $f$ (véase [1], B,
\S~teo. 1). Entonces $\mathcal L'$ es \emph{preamplio sobre $Y$ con respecto
a $S$}, es decir, existen un entero $m$ y un morfismo cuasifinito de $Y$ a un
esquema proyectivo estándar adecuado $\mathbf P_S^N$ tales que
$(\mathcal L')^{\otimes m}$ es isomorfo a la imagen inversa de
$\mathcal O_{\mathbf P_S^N}(1)$.

Obsérvese que, incluso si $X$ es separado sobre $S$, puede ocurrir que $Y$ no
sea separado sobre $S$ (situación que se presenta en problemas de módulos
nada patológicos). Si se cumple $1^\circ$, entonces $Y$ es separado si y solo
si la relación de equivalencia definida por $G$ tiene grafo cerrado, es decir,
si $G\times X\to X\times X$ tiene imagen cerrada (es entonces una inmersión
cerrada). Si $Y$ es separado, entonces $\mathcal L'$ es preamplio sobre $Y$
con respecto a $S$ si y solo si es amplio, es decir, si una potencia tensorial
adecuada define una inmersión proyectiva. En los problemas de módulos
mencionados en la introducción puede demostrarse que la relación de
equivalencia obtenida tiene, en efecto, grafo cerrado.

\medskip
\noindent\textbf{OBSERVACIONES 8.1. ---} Hemos supuesto
$G=\operatorname{GP}(n)_S$ para fijar las ideas y porque este es actualmente
el caso más importante en la práctica. Una hipótesis razonable sobre $G$
parecería ser, en cambio, que $G$ fuese una de las «formas» sobre $S$ de uno
de los grupos de Tohokû (cuya construcción sobre los enteros fue realizada
también por CHEVALLEY). El único hecho positivo que conozco en la dirección de
la conjetura anterior es el siguiente: sea $X$ \emph{un esquema afín sobre un
cuerpo $k$ de característica $0$, en el que el grupo
$\operatorname{Gl}(n)_k$ o $\operatorname{GP}(n-1)_k$ actúa libremente.
Entonces la relación de equivalencia definida por $G$ es efectiva, el
cociente $X/G$ es también afín y el morfismo $X\to X/G$ es plano y
sobreyectivo}. La demostración utiliza el hecho siguiente (que actualmente
solo se ha demostrado en característica $0$): cuando se hace actuar a $G$ en
el anillo de coordenadas afín $A$ de $G$, considerado como espacio vectorial
sobre $k$, la representación trivial de $G$ aparece una sola vez (en una
serie de composición de un subespacio vectorial estable bajo $G$ y de
dimensión finita sobre $k$). Parece posible que un uso sistemático de
\src{118}
la teoría de representaciones lineales de $G$ conduzca finalmente a una
demostración de la conjetura, al menos sobre un cuerpo base. Cuando la base ya
no es un cuerpo, el conferenciante no sabe nada.

\section*{BIBLIOGRAFÍA}
\addcontentsline{toc}{section}{Bibliografía}

\begin{description}
\item[{[1]}] GROTHENDIECK (Alexander). --- Techniques de descente et théorèmes
d'existence en géométrie algébrique, I: Généralités, Séminaire Bourbaki,
1959/60, no. 190, 29 pp.

\item[{[2]}] GROTHENDIECK (Alexander). --- Techniques de descente et théorèmes
d'existence en géométrie algébrique, II: Le théorème d'existence en théorie
formelle des modules, Séminaire Bourbaki, 1959/60, no. 195, 22 pp.

\item[{[3]}] GROTHENDIECK (A.) and DIEUDONNÉ (J.). --- Éléments de géométrie
algébrique, I: Le langage des schémas. --- Paris, Presses universitaires de
France, 1960 (Institut des hautes Études scientifiques, Publications
mathématiques, 4).
\end{description}

\def\FgaUnitKey{e212}\clearpage
\chapter*{Fe de erratas de la exposición n.º 212}
\addcontentsline{toc}{chapter}{Fe de erratas de la exposición n.º 212}
\setcounter{footnote}{2}
\src{301}

\noindent\textbf{5. TDTE III: Preesquemas cociente.}

\begin{center}
[Seminario Bourbaki, t. 13, 1960/61, n.º 212, 20 pp.]
\end{center}

\noindent\textbf{Página 15, afirmación (i).} --- El hecho de que $Y=X/R$ sea
cuasiproyectivo sobre $S$ parece, por el momento, haberse demostrado solo
cuando $R$ proviene de una relación de equivalencia.

\noindent\textbf{Páginas 18 y 19.} --- Como señalamos al final de la exposición
siguiente, la conjetura considerada es definitivamente falsa. El «hecho
positivo» señalado en la Observación 8.1 parece haber sido demostrado
simultáneamente por diversos autores (NAGATA, ROSENLICHT, GROTHENDIECK,
\ldots).

Señalemos que la aplicación hecha en TDTE V de la teoría del paso al cociente
desarrollada aquí a la construcción de esquemas de Picard puede sustituirse
también por un uso adecuado de los esquemas de Hilbert (cf. SHMT). Como se
indicó en el § 8, la laguna más importante de la teoría presentada
\src{302}
aquí es la falta de un criterio de existencia de cocientes por una relación de
equivalencia no propia, como las relaciones de equivalencia procedentes de
ciertas acciones del grupo proyectivo. Un teorema importante en esta dirección
ha sido obtenido por MUMFORD\footnote{MUMFORD (David). --- An elementary
theorem in geometric invariant theory, \emph{Bull. Amer. Math. Soc.}, t. 67,
1961, pp. 483--486.}. Para un refinamiento de su resultado y diversas
aplicaciones a la teoría, véase SHMT.

\def\FgaUnitKey{221}\unit{Exposición 221 -- Los esquemas de Hilbert}
\sid{221}
\src{249}
\seminares{13.er año, 1960/61, n.º 221\\Mayo de 1961}
\paperhead{TÉCNICAS DE CONSTRUCCIÓN Y TEOREMAS DE EXISTENCIA EN GEOMETRÍA ALGEBRAICA\\
IV. LOS ESQUEMAS DE HILBERT}%
  {por Alexander GROTHENDIECK}

\subsection*{Introducción}
\addcontentsline{toc}{subsection}{Introducción}

Las técnicas expuestas en [2], I y II, eran en lo esencial independientes de
toda hipótesis de proyectividad sobre los esquemas considerados.
Desafortunadamente, en la actualidad no permiten resolver los problemas de
existencia planteados en la Exposición II de [2]. En la presente exposición y
en la siguiente resolveremos estos problemas bajo hipótesis de proyectividad.
Las técnicas empleadas son típicamente proyectivas y prácticamente no hacen
uso de los resultados de I y II de [2]. Aquí construiremos los «esquemas de
Hilbert», destinados a sustituir el uso de las coordenadas de Chow, como se
indicó en [2], II, n.º 2. En la exposición siguiente, la teoría del paso al
cociente en esquemas desarrollada en III, unida a la teoría de los esquemas de
Hilbert, nos permitirá, por ejemplo, construir los esquemas de Picard
(definidos en [2], II, n.º 3) bajo condiciones bastante generales.

En resumen, puede decirse que ahora se dispone de una técnica de construcciones
proyectivas bastante satisfactoria, salvo que todavía
falta\footnote{Véase el aditivo al final de la exposición.} un teorema de paso
al cociente por grupos tales como el grupo proyectivo, que actúan «sin puntos
fijos» (cf. [2], III, n.º 8). La situación parece incluso algo mejor en
geometría analítica (cuando uno se limita al estudio de espacios analíticos
«proyectivos» sobre un espacio analítico dado), pues para los espacios
analíticos desaparece la dificultad del paso al cociente por un grupo que
actúa convenientemente. Por otra parte, tanto en geometría algebraica como en
geometría analítica quedaría por desarrollar una técnica de construcción
válida sin hipótesis de proyectividad.

\fsection{1}{Familias acotadas de haces: propiedades de permanencia}

Sean $k$ un cuerpo y $X$ un $k$-preesquema, de tipo finito para simplificar.
Para toda extensión $K/k$ se obtiene, pues, un $K$-preesquema
$X_K=X\otimes_k K$. Si $\mathcal F$ es un haz coherente sobre $X_K$ y $K'$ es
una extensión de $K$, entonces $\mathcal F\otimes_K K'=\mathcal F_{K'}$ es un
haz cuasicoherente sobre $X_K\otimes_K K'=X_{K'}$. Dicho esto, si $K$ y $K'$
son dos extensiones cualesquiera de $k$, $\mathcal F$ es un haz cuasicoherente
sobre $X_K$ y $\mathcal F'$
\src{250}
es un haz cuasicoherente sobre $X_{K'}$, diremos que $\mathcal F$ y
$\mathcal F'$ son \emph{equivalentes} si existen $k$-homomorfismos de $K$ y
$K'$ en una misma extensión $K''$ de $k$ tales que $\mathcal F_{K''}$ y
$\mathcal F'_{K''}$ sean isomorfos sobre $X_{K''}$. Esta es una relación de
equivalencia, y nos interesaremos por las clases de equivalencia de haces para
esta relación y por los conjuntos de clases de equivalencia. Obsérvese que, si
$X$ es de tipo finito sobre $k$,\footnote{El texto impreso lleva aquí $X_0$;
solo $X$ está definido en este pasaje.} toda clase de haces coherentes puede
definirse por un haz coherente sobre algún $X_K$, donde $K$ es una extensión
de tipo finito de $k$. Por consiguiente, en la definición de las clases de
haces coherentes puede uno limitarse a las extensiones algebraicamente cerradas
de $k$; también puede limitarse a una extensión algebraicamente cerrada fija
$\Omega$ de $k$, de grado de trascendencia infinito. Dos haces coherentes
$\mathcal F$ y $\mathcal F'$ sobre $X_\Omega$ son entonces equivalentes si y
solo si existe un $k$-automorfismo $\sigma$ de $\Omega$ tal que
$\mathcal F\otimes_\Omega(\Omega,\sigma)$ sea isomorfo a $\mathcal F'$.
Se comprueba así que hay una correspondencia biyectiva entre las clases de
haces coherentes para una u otra definición.

Sean $E$ y $E'$ dos conjuntos de clases de haces coherentes sobre $X$.
Consideremos las clases de todos los haces de la forma
$\mathcal F\otimes\mathcal F'$, donde $\mathcal F$ y $\mathcal F'$ son haces
coherentes sobre un mismo $X_K$, la clase de $\mathcal F$ pertenece a $E$ y la
de $\mathcal F'$ pertenece a $E'$. Se obtiene así un conjunto de clases de
haces coherentes, que denotaremos por $E\otimes E'$. Se definen del mismo modo
los $\operatorname{Tor}_i(E,E')$, etc. En general, a toda función $U$ que
asocia a cada sucesión $\mathcal F_1,\ldots,\mathcal F_n$ de $n$ haces
coherentes sobre un mismo $X_K$ un conjunto
$U(\mathcal F_1,\ldots,\mathcal F_n)$ de haces coherentes sobre $X_K$, y que
posee la evidente propiedad de compatibilidad con los isomorfismos de haces y
las imágenes inversas por cambio de base, se le asocia una función, denotada
por el mismo símbolo $U$, que asigna a cada sucesión de $n$ conjuntos
$E_1,\ldots,E_n$ de clases de haces coherentes un conjunto
$U(E_1,\ldots,E_n)$ de clases de haces coherentes.

Nuestro propósito en este apartado es dar una definición de ciertos conjuntos
de clases de haces, llamados \emph{acotados}, y demostrar que las operaciones
$U$ más usuales, aplicadas a conjuntos acotados, vuelven a dar conjuntos
acotados.

Sea $X$ un preesquema de tipo finito sobre un $S$ noetheriano. Para todo
$s\in S$, la fibra $X_s$ es un preesquema de tipo finito sobre $k(s)$, y
consideraremos clases de haces coherentes sobre $X_s$ en el sentido anterior.
Esto da sentido a la expresión «clase de haces coherentes sobre una fibra de
$X/S$» y a las expresiones análogas. Del mismo modo, procediendo por separado
sobre cada fibra, pueden seguir considerándose operaciones como
$E\otimes E'$, etc., que hacen corresponder a sistemas de conjuntos de
\src{251}
clases de haces coherentes sobre las fibras de $X/S$ otro conjunto de clases
de haces coherentes sobre las fibras de $X/S$.

\result{DEFINICIÓN}{1.1} Sea $E$ un conjunto de clases de haces coherentes
sobre las fibras de $X/S$. Se dice que $E$ está \emph{acotado} si existen un
preesquema $S'$ de tipo finito sobre $S$ y un haz coherente $\mathcal F'$ sobre
$X'=X\times_S S'$ tales que $E$ esté contenido en el conjunto de clases de
haces sobre las fibras de $X/S$ definido por $\mathcal F'$.

Por definición, este último hace corresponder a un $s\in S$ las clases de los
haces $\mathcal F'\otimes_{S'}k(s')$, donde $s'$ recorre los puntos de $S'$
situados sobre $s$ (de modo que $k(s')$ es una extensión de $k(s)$ y
$X_s\otimes_{k(s)}k(s')$ se identifica con la fibra
$X'\otimes_{S'}k(s')=X'_{s'}$ de $X'$ en $s'$). Puede decirse que las familias
acotadas son aquellas contenidas en una \emph{familia algebraica} de haces
coherentes parametrizada por un $S'$ de tipo finito sobre $S$.

Una unión finita de familias acotadas es acotada (tómese el preesquema suma de
los preesquemas de parámetros $S_i$ que definen las familias algebraicas
mayorantes). Un cambio de base $T\longrightarrow S$ transforma una familia
acotada relativa a $X/S$ en una familia acotada relativa a $X_T/T$, y el
recíproco es cierto si $T\longrightarrow S$ es sobreyectivo (o, más
generalmente, si su imagen contiene los $s$ que intervienen efectivamente en la
familia dada $E$ para $X/S$). Esto reduce, en principio, la determinación de
las familias acotadas al caso en que $S$ es el espectro de un álgebra de tipo
finito sobre el anillo de los enteros $\mathbf Z$.

Si $E$ y $E'$ son familias acotadas de clases de haces relativas a $X/S$,
entonces $E\otimes E'$ también está acotada. En efecto, supongamos que $E$ y
$E'$ están mayoradas, respectivamente, por las familias algebraicas definidas
por $T\longrightarrow S$ y $\mathcal F$ sobre $X_T$, y por
$T'\longrightarrow S$ y $\mathcal F'$ sobre $X_{T'}$. Se ve entonces que
$E\otimes E'$ está mayorada por la familia algebraica definida por
$T''=T\times_S T'\longrightarrow S$ y el haz $\mathcal F''$ sobre
$X_{T''}$, producto tensorial de las imágenes inversas de $\mathcal F$ y
$\mathcal F'$ sobre $X_{T''}$. Este razonamiento solo es válido porque el
funtor $\mathcal F\otimes\mathcal F'$ en $\mathcal F,\mathcal F'$ es exacto a
derecha y, por tanto, conmuta con la extensión de la base (y, en particular,
con el paso a las fibras). No es aplicable directamente a operaciones locales
tales como $\operatorname{Tor}_i(E,E')$, $\operatorname{Hom}(E,E')$ y
$\operatorname{Ext}^i(E,E')$. No obstante, puede demostrarse que estas
operaciones siguen transformando conjuntos acotados en conjuntos acotados,
procediendo como para $E\otimes E'$ pero utilizando además resultados del tipo
siguiente (todos ellos contenidos en [3], IV, 6.11): una familia acotada $E$
está siempre mayorada por una familia algebraica definida por un haz coherente
$\mathcal F$ sobre un $X_T$ ($T$ de tipo finito sobre $S$) que es plano
relativamente a $T$. (Se «corta en trozos» el espacio de parámetros inicial.)
Tales
\src{252}
propiedades de planitud de haces adecuados aseguran, en efecto, que operaciones
como $\operatorname{Tor}_i(\mathcal F,\mathcal F')$ conmuten con cualquier
cambio de base. El mismo método se aplica a operaciones de naturaleza global:
imágenes directas e imágenes directas derivadas de haces coherentes por
morfismos propios, $\operatorname{Ext}$ globales relativos a morfismos propios
(cf. [5], III, \S~6), etc. Todas estas operaciones transforman familias
acotadas de haces en familias acotadas de haces. (N. B. -- Aquí los
preesquemas sobre los que se toman los diversos haces pueden cambiar por las
operaciones del tipo considerado.)

Los dos enunciados siguientes se demuestran esencialmente mediante la misma
técnica de planitud; para la descomposición primaria sobre las fibras de un
morfismo de tipo finito, véase en particular [5], IV.

\result{PROPOSICIÓN}{1.2} Sean $E$ y $E'$ conjuntos acotados de clases de
haces sobre las fibras de $X/S$, donde se supone que $X$ es propio sobre $S$.
Entonces:

\medskip
\noindent (i) La familia de núcleos, conúcleos e imágenes de homomorfismos
$\mathcal F\longrightarrow\mathcal F'$, donde la clase de $\mathcal F$
pertenece a $E$ y la de $\mathcal F'$ pertenece a $E'$, está acotada.

\smallskip
\noindent (ii) La familia de haces $\mathcal F''$ que son extensiones de un
$\mathcal F$ por un $\mathcal F'$, donde la clase de $\mathcal F$ pertenece a
$E$ y la de $\mathcal F'$ pertenece a $E'$, está acotada.

Mediante un cambio de base conveniente puede suponerse que $E$ y $E'$ están
definidos, respectivamente, por haces coherentes $\mathcal G$ y $\mathcal G'$
sobre un $X_T/T$, donde $T$ es de tipo finito sobre $S$. Además, puede
suponerse que se satisfacen ciertas hipótesis de planitud que implican que la
formación de los haces
$f_{T*}(\operatorname{Hom}_{X_T}(\mathcal G,\mathcal G'))$ y
$\operatorname{Ext}^{1}_{f_T}(\mathcal G,\mathcal G')$ conmuta con el cambio
de base para un morfismo arbitrario $T'\longrightarrow T$. Asimismo, puede
suponerse que estos haces coherentes sobre $T$ son localmente libres. Sean
$T_0$ y $T_1$ los fibrados vectoriales sobre $T$ cuyos haces de gérmenes de
secciones son, respectivamente, los haces precedentes. Se define entonces
canónicamente un homomorfismo
$\mathcal G_{T_0}\longrightarrow\mathcal G'_{T_0}$ de haces coherentes sobre
$X_{T_0}$ y una extensión
\[
0\longrightarrow\mathcal G'_{T_1}\longrightarrow\mathcal G''
\longrightarrow\mathcal G_{T_1}\longrightarrow0
\]
de haces coherentes sobre $X_{T_1}$ que posee una propiedad universal
evidente.\footnote{El texto impreso invierte los términos extremos de esta sucesión. Se trata de extensiones de $\mathcal G$ por $\mathcal G'$, como indica $\operatorname{Ext}^1(\mathcal G,\mathcal G')$.} El segundo haz define una familia algebraica que mayora la familia
considerada en (ii). Lo mismo vale para el núcleo, el conúcleo y la imagen del
homomorfismo
\src{253}
precedente, cuya consideración demuestra (i) (siempre que se suponga que el
conúcleo es plano relativamente a $T_0$, caso al que se reduce de nuevo
cortando $T_0$ en trozos\ldots).

\result{PROPOSICIÓN}{1.3} Sea $E$ una familia acotada de clases de haces sobre
las fibras de $X/S$. Entonces las clases de los haces estructurales de los
$(\operatorname{supp}\mathcal F)_{\mathrm{red}}$, donde $\mathcal F$ es un
haz coherente sobre un $X_K$, con $K$ algebraicamente cerrado, cuya clase
pertenece a $E$, forman una familia acotada.

Aquí $(\operatorname{supp}\mathcal F)_{\mathrm{red}}$ denota el soporte de
$\mathcal F$ provisto de la estructura reducida inducida; es decir, su haz
estructural es el cociente de $\mathcal O_{X_K}$ por el mayor haz de ideales
que define $\operatorname{supp}\mathcal F$. Puede demostrarse el resultado
análogo a (1.3) para los haces deducidos canónicamente de $\mathcal F$ por la
teoría de la descomposición primaria; por ejemplo, los
$\mathcal F/\mathcal F_i$, donde los $\mathcal F_i$ son los subhaces de
$\mathcal F$ primarios para las componentes del soporte de $\mathcal F$ y
mínimos con esta propiedad; o los $\mathcal O_{X_K}/\mathfrak p$, donde
$\mathfrak p$ es un haz de ideales primos asociado a $\mathcal F$; o los
$\mathcal O_{X_K}/\mathfrak q$, donde $\mathfrak q$ es un haz de ideales
primarios asociado a una componente del soporte de $\mathcal F$ (siendo el
cuerpo de referencia algebraicamente cerrado).

\fsection{2}{Familias acotadas y polinomio de Hilbert}

En lo que sigue suponemos que $X$ es proyectivo sobre $S$ y está provisto de
un haz muy amplio, denotado $\mathcal O_X(1)$. Para toda extensión $K$ de una
extensión residual $k(s)$ de un punto $s$ de $S$, se considera sobre $X_K$ el
haz correspondiente $\mathcal O_{X_K}(1)$, que sigue siendo muy amplio.

A todo haz coherente $\mathcal F$ sobre $X_K$ se le asocia la función
\[
P_{\mathcal F}(n)=\text{característica de Euler--Poincaré de $\mathcal F(n)$
sobre $X_K$},
\]
que es un polinomio en el entero $n$, llamado \emph{polinomio de Hilbert} de
$\mathcal F$. Para valores grandes de $n$, $P(n)$ no es sino la dimensión de
$H^0(X_K,\mathcal F(n))$ sobre $K$, puesto que los
$H^i(X_K,\mathcal F(n))$ se anulan para $i>0$ y $n$ grande.

Supongamos ahora que $\mathcal F$ es un haz coherente sobre $X$, plano sobre
$S$. Entonces los polinomios de Hilbert de los haces $\mathcal F_s$ inducidos
sobre fibras $X_s$ pertenecientes a una misma componente conexa de $S$ son
todos iguales [5], III, \S~7. De ello se sigue (sin hipótesis de planitud) que
el conjunto de polinomios de Hilbert de los haces $\mathcal F_s$, $s\in S$,
es finito para todo haz coherente $\mathcal F$ sobre $X$.
\src{254}
Recordemos además que, si $\mathcal F$ es un haz coherente sobre $X$, es
isomorfo a un haz cociente de un haz $\mathcal O_X(-n)^N$, para $n$ y $N$
suficientemente grandes. Por tanto, los haces $\mathcal F_s$ inducidos sobre
las fibras también son cocientes del haz $\mathcal O(-n)^N$ sobre la fibra.\footnote{El texto impreso omite aquí el exponente $N$; la restricción a la fibra conserva la potencia directa de rango $N$ del haz precedente.}

De estas dos observaciones se deduce la parte «solo si» del teorema siguiente:

\result{TEOREMA}{2.1} Sea $X$ proyectivo sobre un $S$ noetheriano y sea
$\mathcal O_X(1)$ muy amplio sobre $X$ relativamente a $S$. Sea $E$ un
conjunto de clases de haces sobre las fibras de $X/S$. Para que $E$ esté
acotado, es necesario y suficiente\footnote{La exposición impresa dice
solamente «es necesario»; la fe de erratas de 1962 de la Exposición 221 prescribe «es
necesario y suficiente».} que satisfaga las condiciones siguientes:

\medskip
\noindent (a) Existe un haz coherente $\mathcal L$ sobre $X$ (que puede
suponerse de la forma $\mathcal O_X(-n)^N$) tal que $E$ esté contenido en la
familia de clases de haces coherentes cocientes de haces de la forma
$\mathcal L_K$.

\smallskip
\noindent (b) Los polinomios de Hilbert $P_{\mathcal F}$ de los haces
$\mathcal F$ cuya clase pertenece a $E$ son elementos de un mismo conjunto
finito de polinomios.

Queda por demostrar la parte «si», que será un caso particular de un resultado
más preciso. Para todo Módulo coherente $\mathcal F$ sobre un preesquema de
tipo finito sobre un cuerpo $K$, y todo entero $r$, sea $N_r$ el sub-Módulo de
$\mathcal F$ cuyas secciones sobre un abierto son las secciones de $\mathcal F$
sobre dicho abierto cuyo soporte tiene dimensión $<r$. Se tiene, pues,
$N_r=\mathcal F$ para $r>\dim\operatorname{supp}\mathcal F$, mientras que
$N_r=0$ para $r\leq0$, y se obtiene una filtración creciente finita de
$\mathcal F$ cuyos factores $N_{r+1}/N_r$\footnote{El texto impreso invierte
el cociente, escribiendo $N_r/N_{r+1}$, lo que es incompatible con la
filtración creciente que acaba de definirse.} tienen como ciclos primos
asociados exactamente los ciclos primos asociados a $\mathcal F$ que tienen
dimensión $r$. Pondremos
\[
\mathcal F_{(r)}=\mathcal F/N_r.
\]
Entonces los ciclos primos asociados a $\mathcal F_{(r)}$ son exactamente los
ciclos primos asociados a $\mathcal F$ que tienen dimensión $\geq r$; en
particular, es igual a $\mathcal F$ si y solo si los ciclos primos asociados a
$\mathcal F$ tienen dimensión $\geq r$. Con esta notación:

\result{TEOREMA}{2.2} Bajo las hipótesis de (2.1), sea $s$ un entero y
supongamos que $E$ satisface la condición (a) y la siguiente forma debilitada
de (b):

\medskip
\noindent $(b_s)$ Los polinomios de Poincaré $P_{\mathcal F}$ de los haces
$\mathcal F$ cuya clase pertenece a $E$ tienen acotados sus coeficientes en
grados $\geq s-1$.

En estas condiciones, los haces $\mathcal F_{(s)}$ (cuando la clase de
$\mathcal F$ recorre $E$) forman una familia acotada. Además, los coeficientes
de grado $s-2$ de los $P_{\mathcal F}$ están minorados.
\src{255}
Así pues:

\result{COROLARIO}{2.3} Supongamos que los haces $\mathcal F$ cuya clase
pertenece a $E$ son tales que todos sus ciclos primos asociados tienen
dimensión $d$ con $s\leq d\leq r$. Entonces, en (2.1) (b), basta considerar los
coeficientes de los $P_{\mathcal F}$ en grados comprendidos entre $s-1$ y $r$.

El resto de este apartado está dedicado a esbozar la demostración de (2.2).
Los lemas clave son los dos lemas siguientes, el primero de los cuales es bien
conocido (y resume el contenido matemático útil de las coordenadas de Chow):

\result{LEMA 2.4 (CHOW)}{} Consideremos los haces estructurales de subesquemas
$Y$ de fibras $X_K$ ($K$ recorre las extensiones algebraicamente cerradas de
cuerpos residuales de $S$), donde $Y$ es reducido y todas sus componentes
tienen la misma dimensión $r$ (considerándose $\mathcal O_Y$ como un haz
cociente de $\mathcal O_X$). Si los grados de los $Y$ están acotados, los $Y$
forman una familia acotada.

Aquí, el grado $a$ de $Y$ puede definirse del modo más cómodo mediante el
coeficiente del término principal de
\[
P_{\mathcal O_Y}=a\frac{n^r}{r!}+\cdots.
\]

\result{LEMA}{2.5} Sea $\mathcal L$ un haz coherente sobre $X$, y sea $E$ un
conjunto de clases de haces cocientes $\mathcal F$ de haces $\mathcal L_K$
($K$ una extensión residual de $S$). Supongamos que las fibras de $X$ sobre
$S$ tienen dimensión $\leq r$, y escribamos
\[
P_{\mathcal F}(n)=a_{\mathcal F}\frac{n^r}{r!}
+b_{\mathcal F}\frac{n^{r-1}}{(r-1)!}
+\text{términos de grado $<r-1$}.
\]
Entonces el coeficiente $a_{\mathcal F}$ está acotado, mientras que
$b_{\mathcal F}$ está minorado. Si $b_{\mathcal F}$ está acotado, la familia
de los $\mathcal F_{(r)}$ está acotada.

\emph{Demostración.} -- Tras sustituir $S$ por una unión de subesquemas de
$S$ que recubren $S$, puede suponerse que existe un morfismo finito
$f:X\longrightarrow\mathbf P^r_S$ tal que $\mathcal O_X(1)$ es isomorfo a la
imagen inversa de $\mathcal O_{\mathbf P^r_S}(1)$. Por tanto, para todo haz
coherente $\mathcal F$ sobre $X$ se tiene
$P_{\mathcal F}=P_{f_*(\mathcal F)}$. Por otra parte, se comprueba fácilmente
(mediante la técnica del apartado anterior) que un conjunto de haces
$\mathcal F$ sobre $X$ está acotado si y solo si lo está el conjunto de los
$f_*(\mathcal F)$. Finalmente,
\[
f_*(\mathcal F)_{(r)}=f_*(\mathcal F_{(r)}).
\]
\src{256}
Esto nos reduce al caso $X=\mathbf P^r_S$. Además, puede suponerse que
$\mathcal L=\mathcal O_{\mathbf P^r_S}(k)^s$ para ciertos $k$ y $s$. El
coeficiente $a_{\mathcal F}$ satisface
\[
0\leq a_{\mathcal F}\leq s
\]
y, por tanto, está acotado. Es equivalente decir que el coeficiente de
$n^{r-1}$ en los $P_{\mathcal F}(n)$ está minorado (respectivamente, acotado),
o decir lo mismo de
$P_{\mathcal F}(n-k)=P_{\mathcal F(-k)}(n)$. Nos reducimos así al caso
\[
\mathcal L=\mathcal O_{\mathbf P^r_S}^{s}.
\]

Consideremos la sucesión exacta
\[
0\longrightarrow N_r\longrightarrow\mathcal F
\longrightarrow\mathcal F_{(r)}\longrightarrow0,
\]
de donde
\[
P_{\mathcal F}=P_{\mathcal F_{(r)}}+P_{N_r}.
\]
Como el coeficiente de $n^{r-1}$ en $P_{N_r}$ es positivo (pues
$\dim\operatorname{supp}N_r\leq r-1$), se tiene
\[
b_{\mathcal F_{(r)}}\leq b_{\mathcal F}.
\]
Así pues, para demostrar las afirmaciones del lema podemos sustituir
$\mathcal F$ por $\mathcal F_{(r)}$; es decir, podemos suponer que los haces
cocientes $\mathcal F$ de $\mathcal L$ considerados son libres de torsión.

Como $\mathbf P^r_K$ es normal, se sigue que $\mathcal F$ es localmente libre
de rango $a=a_{\mathcal F}$ sobre un abierto $U=\mathbf P^r_K-Y$, donde $Y$
tiene codimensión $\geq2$. Por tanto, $\bigwedge^a\mathcal F$ es un haz sobre
$\mathbf P^r_K$ cuya restricción a $U$ es invertible; luego (como
$\mathbf P^r_K$ es regular y $Y$ tiene codimensión $\geq2$) es isomorfo a la
restricción de un haz invertible sobre $\mathbf P^r_K$, determinado salvo
isomorfismo. Este último es de la forma $\mathcal O_{\mathbf P^r_K}(d)$ para
un entero $d$ bien determinado. Como $\bigwedge^a\mathcal F$ es un cociente
de $\bigwedge^a\mathcal O_{\mathbf P^r_K}^{s}\simeq
\mathcal O_{\mathbf P^r_K}^{N}$\footnote{El texto impreso sustituye aquí el
exponente $s$ por un $n$ no definido; $s$ es el rango del haz libre que se
proyecta sobre $\mathcal F$.}, donde $N=\binom{s}{a}$, posee $N$ secciones
canónicas, y estas definen, por tanto, secciones de
\src{257}
$\mathcal O_{\mathbf P^r_K}(d)$ sobre $U$ que son restricciones de secciones
$s_i$ ($1\leq i\leq N$) de $\mathcal O_{\mathbf P^r_K}(d)$ (pues
$\mathbf P^r_K$ es normal y $Y$ tiene codimensión $\geq2$). Estas secciones
generan $\mathcal O_{\mathbf P^r_K}(d)$ en los puntos de $U$, de modo que no
son todas nulas; esto implica $d\geq0$. Por otra parte, un cálculo sencillo da
\[
b_{\mathcal F}=a_{\mathcal F}(r+1)/2+d.
\]
En particular, esto demuestra que $b_{\mathcal F}>0$, por lo que
$b_{\mathcal F}$ está minorado. Está acotado si y solo si $d$ lo está;
demostremos que, en tal caso, $\mathcal F$ pertenece a una familia acotada.
Podemos suponer que $a_{\mathcal F}$ y $b_{\mathcal F}$ son fijos, digamos
iguales a $a$ y $b$, de modo que $d$ es fijo. El conocimiento de las $N$
secciones $s_i$ de $\mathcal O_{\mathbf P^r_K}(d)$, es decir, de un
homomorfismo
\[
s:\bigwedge^a\mathcal L_K\longrightarrow\mathcal O_{\mathbf P^r_K}(d),
\]
permite recuperar $\mathcal F$ como la coimagen del homomorfismo compuesto
correspondiente
\[
\mathcal L_K\longrightarrow
\operatorname{Hom}\!\left(\bigwedge^{a-1}\mathcal L_K,
\bigwedge^a\mathcal L_K\right)\longrightarrow
\operatorname{Hom}\!\left(\bigwedge^{a-1}\mathcal L_K,
\mathcal O_{\mathbf P^r_K}(d)\right),
\]
donde el primero es el homomorfismo canónico procedente del producto exterior
y el segundo se deduce de $s$. La conclusión se sigue ahora de (1.2) (i).

Combinando los dos lemas precedentes se obtiene:

\result{LEMA}{2.6} Bajo las hipótesis preliminares de (2.1), supongamos que,
para todo $\mathcal F$, se tiene
\[
P_{\mathcal F}(n)=a_{\mathcal F}\frac{n^r}{r!}
+b_{\mathcal F}\frac{n^{r-1}}{(r-1)!}
+\text{términos de grado $<r-1$},
\]
y que los coeficientes $a_{\mathcal F}$ están acotados. Entonces los
coeficientes $b_{\mathcal F}$ están minorados. Si los $b_{\mathcal F}$ están
acotados, los $\mathcal F_{(r)}$ forman una familia acotada.

Podemos suponer que los cuerpos base $K$ de los haces $\mathcal F$ son
algebraicamente cerrados. Dótese cada
$\operatorname{supp}\mathcal F_{(r)}$, unión de las componentes de
$\operatorname{supp}\mathcal F$ de dimensión $r$, de la estructura reducida
inducida. Entonces los grados de los $\operatorname{supp}\mathcal F_{(r)}$
están acotados superiormente por $a$, por lo que, en virtud de (2.4), los
$\operatorname{supp}\mathcal F_{(r)}$ forman un conjunto acotado. Además, para
cada componente de $\operatorname{supp}\mathcal F_{(r)}$, la longitud de
$\mathcal F_{(r)}$ a lo largo de dicha componente es $\leq a$. Por tanto, si
$\mathcal J_{\mathcal F}$ es el ideal que define
$\operatorname{supp}\mathcal F_{(r)}$, entonces $\mathcal F_{(r)}$
\src{258}
puede considerarse como un Módulo sobre el subesquema $Y_{\mathcal F}$ de $X$
definido por $\mathcal J_{\mathcal F}^{a}$. Como en el lema anterior, también
podemos reducirnos al caso $\mathcal F=\mathcal F_{(r)}$, de modo que
$\mathcal F$ proviene de un Módulo sobre $Y_{\mathcal F}$. Los
$Y_{\mathcal F}$ corresponden a una familia acotada de Módulos cocientes de
los $\mathcal O_{X_K}$ y, por consiguiente, provienen de un subesquema cerrado
$Y$ de un esquema $X\times_S T$. Podemos entonces aplicar (2.5) a $Y/T$ y
$\mathcal L\otimes_XY$, lo que da la conclusión.

Ya podemos demostrar (2.2) por inducción sobre la cota superior $r$ de
$\dim\operatorname{supp}\mathcal F$. El enunciado es trivial para $r<0$;
supongamos, pues, que $r\geq0$ y que el enunciado ha sido demostrado para
$r'<r$. Por (2.6), los $\mathcal F_{(r)}$ forman un conjunto acotado; por
tanto, en virtud de (1.2) (i), también lo forman los núcleos de los
homomorfismos $\mathcal L_K\longrightarrow\mathcal F_{(r)}$. Existe, pues, un
Módulo coherente $\mathcal L'$ sobre $X$ tal que los núcleos en cuestión y,
por consiguiente, también los
\[
N_r(\mathcal F)=\operatorname{Ker}(\mathcal F\longrightarrow\mathcal F_{(r)}),
\]
son cocientes de los Módulos $\mathcal L'_K$. Como los
$\mathcal F_{(r)}$ están acotados, los $P_{\mathcal F_{(r)}}$ permanecen
acotados, y la fórmula
\[
P_{\mathcal F}=P_{\mathcal F_{(r)}}+P_{N_r}
\]
demuestra que los $P_{N_r}$ satisfacen la misma condición $(b_s)$ que los
$P_{\mathcal F}$. Por tanto, los $N_r$ satisfacen las condiciones (a) y
$(b_s)$, y por la hipótesis de inducción los $(N_r)_{(s)}$ permanecen acotados.
Ahora bien, $\mathcal F_{(s)}$ es una extensión de $\mathcal F_{(r)}$ por
$(N_r)_{(s)}$; luego, por (1.2) (ii), los $\mathcal F_{(s)}$ permanecen
acotados. Para la última afirmación de (2.2), llevamos simultáneamente la inducción sobre el acotamiento y sobre la cota inferior del coeficiente siguiente. Los núcleos acotados de los morfismos de fuente fija hacia $\mathcal F_{(s)}$ proporcionan un mismo haz coherente que se aplica sobreyectivamente sobre cada $N_s$, pero (1.2) (i) por sí solo no hace acotada la familia de estos cocientes. Reaplicamos por inducción la afirmación sobre los coeficientes a estos cocientes; la aditividad del polinomio de Hilbert y el Lema 2.6 dan entonces la cota inferior requerida para el coeficiente siguiente. Esto termina la demostración.

\fsection{3}{Esquemas de Hilbert: definición y teorema de existencia}

Sea $X$ un preesquema sobre otro preesquema $S$, y sea $\mathcal F$ un Módulo
cuasicoherente sobre $X$. Denotemos por
\[
\operatorname{Quot}(\mathcal F/X/S)
\]
el conjunto de Módulos cuasicoherentes cocientes de $\mathcal F$ que son
planos sobre $S$. Sea ahora $S'\longrightarrow S$ un morfismo de cambio de
base, y pongamos $X'=X\times_S S'$ y
$\mathcal F'=\mathcal F\otimes_{\mathcal O_S}\mathcal O_{S'}$. Así, $X'$ es
un preesquema sobre $S'$ provisto de un Módulo
\src{259}
$\mathcal F'$, y podemos considerar
$\operatorname{Quot}(\mathcal F'/X'/S')$. Pongamos
\[
\underline{\operatorname{Quot}}_{\mathcal F/X/S}(S')
=\operatorname{Quot}(\mathcal F'/X'/S')
\qquad\text{(donde $X'=X\times_SS'$)}.
\]
Si ahora $S''\longrightarrow S'$ es un $S$-morfismo, entonces
$X''=X\times_SS''$ es isomorfo a $X'\times_{S'}S''$, y $\mathcal F''$ es
isomorfo a $\mathcal F'\otimes_{\mathcal O_{S'}}\mathcal O_{S''}$. Como el
funtor de imagen inversa
\[
\mathcal G'\longmapsto
\mathcal G'\otimes_{\mathcal O_{S'}}\mathcal O_{S''}
\]
de la categoría de Módulos cuasicoherentes sobre $X'$ a la categoría de
Módulos cuasicoherentes sobre $X''$ es exacto a derecha y transforma Módulos
planos sobre $S'$ en Módulos planos sobre $S''$, hay una aplicación natural
\[
\operatorname{Quot}(\mathcal F'/X'/S')\longrightarrow
\operatorname{Quot}(\mathcal F''/X''/S''),
\]
por lo que $\underline{\operatorname{Quot}}_{\mathcal F/X/S}(S')$ es un
funtor contravariante de $S'$ (preesquema sobre $S$), con valores en la
categoría de conjuntos. En adelante supondremos que $X$ es proyectivo sobre un
$S$ noetheriano y que $\mathcal F$ es coherente; para simplificar, nos
limitaremos a los $S'$ sobre $S$ que sean localmente noetherianos.

\result{TEOREMA}{3.1} Bajo estas hipótesis, el funtor contravariante
$\underline{\operatorname{Quot}}_{\mathcal F/X/S}$ sobre la categoría de los
$S$-preesquemas localmente noetherianos es representable por un $S$-preesquema
$\operatorname{Quot}_{\mathcal F/X/S}$ que es unión disjunta de una sucesión
de $S$-esquemas proyectivos (en particular,
$\operatorname{Quot}_{\mathcal F/X/S}$ es localmente de tipo finito sobre
$S$).

Obtendremos tal descomposición como sigue. Sea $\mathcal O_X(1)$ un haz
invertible sobre $X$, muy amplio relativamente a $S$. Para todo polinomio
$P(n)$ con coeficientes racionales, sea
$\operatorname{Quot}^P(\mathcal F/X/S)$ el subconjunto de
$\operatorname{Quot}(\mathcal F/X/S)$ formado por los cocientes coherentes
$\mathcal G$ de $\mathcal F$ que son planos sobre $S$ y cuyo polinomio de
Hilbert en todo $s\in S$ es igual a $P$. Pongamos
\[
\underline{\operatorname{Quot}}^P_{\mathcal F/X/S}(S')
=\operatorname{Quot}^P(\mathcal F'/X'/S').
\]
Esto da un subfuntor de
$\underline{\operatorname{Quot}}_{\mathcal F/X/S}$. La propiedad de
invariancia de los polinomios de Hilbert recordada en el Apartado 2 implica lo
siguiente: el funtor $\underline{\operatorname{Quot}}_{\mathcal F/X/S}$ es
representable si y solo si lo son los
$\underline{\operatorname{Quot}}^P_{\mathcal F/X/S}$; en tal caso, el
$S$-preesquema representante $\operatorname{Quot}_{\mathcal F/X/S}$ es
isomorfo a la unión disjunta de los
$\operatorname{Quot}^P_{\mathcal F/X/S}$ que representan los funtores
$\underline{\operatorname{Quot}}^P_{\mathcal F/X/S}$. Así, (3.1) se deduce
del teorema siguiente:
\src{260}
\result{TEOREMA}{3.2} Con las notaciones anteriores, el funtor
$\underline{\operatorname{Quot}}^P_{\mathcal F/X/S}$ está representado por un
$S$-preesquema proyectivo $\operatorname{Quot}^P_{\mathcal F/X/S}$.

El resto de este apartado está dedicado a la demostración de (3.2).

Sea $\nu$ un entero. Para todo $S'$ sobre $S$, denotemos por $A_\nu(S')$ el
conjunto de cocientes $\mathcal G=\mathcal F'/\mathcal K$ de
$\mathcal F'=\mathcal F\otimes_{\mathcal O_S}\mathcal O_{S'}$ que son
coherentes y planos sobre $S'$, y que satisfacen las condiciones siguientes:
\begin{align*}
\text{a.}\quad &R^if'_*(\mathcal G(n))=0 &&\text{para $i>0$ y $n\geq\nu$},\\
\text{b.}\quad &R^if'_*(\mathcal K(n))=0 &&\text{para $i>0$ y $n\geq\nu$},\\
\text{c.}\quad &f'_*(\mathcal K(\nu+k))
=\mathcal S'_k f'_*(\mathcal K(\nu)) &&\text{para $k\geq0$}.
\end{align*}
Para la última relación, se supone que $X$ está escrito como el espectro primo
homogéneo de un álgebra graduada cuasicoherente $\mathcal S_*$ sobre $S$,
graduada positivamente y generada por $\mathcal S_1$, y se pone
$\mathcal S'_* =\mathcal S_*\otimes_{\mathcal O_S}\mathcal O_{S'}$, de modo
que $X'$ es el espectro primo homogéneo de $\mathcal S'_*$. Para demostrar
(3.2), se reduce fácilmente al caso $X=\mathbf P^r_S$ (pues $S$ es una unión
de abiertos $U$ tales que $X|U$ es un subesquema cerrado de $\mathbf P^r_U$ y
$\mathcal O_X(1)$ está inducido por $\mathcal O_{\mathbf P^r_U}(1)$), y al
caso en que $\mathcal F$ es de la forma
$\mathcal O_{\mathbf P^r_S}^{N}$ y es, por tanto, plano sobre $S$. En la
discusión precedente, los haces $\mathcal K$ son entonces también planos sobre
$S'$. Las relaciones de Künneth ([5], III, \S~7) demuestran que las condiciones
(a) y (b) son estables por cambio de base e implican que, para $n\geq\nu$, la
formación de $f'_*(\mathcal G(n))$ y $f'_*(\mathcal K(n))$ conmuta con la
extensión de la base (loc. cit.); por tanto, (c) también es estable por
extensión de la base.

En otras palabras, $A_\nu(S')$ es un funtor contravariante de $S'$, más
precisamente un subfuntor de
$A(S')=\underline{\operatorname{Quot}}^P_{\mathcal F/X/S}(S')$. Al variar
$\nu$, se obtiene una sucesión creciente de subconjuntos $A_\nu(S')$ de
$A(S')$ cuya unión es $A(S')$, por un teorema bien conocido de Serre [5], III,
\S~2. Supongamos ahora que $\mathcal G$ es un cociente coherente de
$\mathcal F'$, plano sobre $S'$, y que $s$ es un punto de $S'$ tal que el
cambio de base $\operatorname{Spec}(k(s))\longrightarrow S'$ da un cociente
$\mathcal G_s$ de $\mathcal F_s$ que satisface las condiciones (a), (b) y (c),
es decir, que pertenece a $A_\nu(\operatorname{Spec}(k(s)))$. Existe entonces
un entorno abierto $U$ de $s$ sobre el cual se satisfacen las mismas
condiciones para $\mathcal G|f'^{-1}(U)$; es decir, este cociente pertenece a
$A_\nu(U)$. Para (a) y (b), esto resulta del «teorema de las funciones
holomorfas» [5], III, \S~7; la condición (c) resulta del lema de Nakayama y del
hecho de que, en cualquier caso, se sabe que
$f'_*(\mathcal K(n+k))=\mathcal S'_k f'_*(\mathcal K(n))$
\src{261}
para $n$ suficientemente grande y $k\geq0$ (cf. [5], III, \S~2). De estas
observaciones se deduce lo siguiente (compárese [4], IV): el funtor $A$ es
representable si y solo si los funtores $A_\nu$ son representables; en tal
caso, el $S$-preesquema $Q$ que representa $A$ es unión creciente de
subesquemas abiertos $Q_\nu$ que representan los $A_\nu$.

Sea
\[
M_* =\bigoplus_{n\geq0}f_*(\mathcal F(n))=\mathcal S_*^{N},
\]
de modo que
\[
M'_* =M_*\otimes_{\mathcal O_S}\mathcal O_{S'}
=\bigoplus_{n\geq0}f'_*(\mathcal F'(n))=\mathcal S'_*{}^{N}.
\]
De (a) se sigue que

\noindent (a') $f'_*(\mathcal G(n))$ es localmente libre de rango $P(n)$ para
$n\geq\nu$, [5], III, \S~7; y de (b), con $i=1$, que

\noindent (a'') $f'_*(\mathcal G(n))$ es un Módulo cociente de $M'_n$.

Además, el conocimiento de este Módulo cociente para $n=\nu$ determina, por
(c), los sub-Módulos $f'_*(\mathcal K(n))$ de $M'_n$ para $n\geq\nu$ y, por
tanto, determina $\mathcal K$ y luego $\mathcal G$. Se obtiene así una
aplicación inyectiva
\[
A_\nu(S')\longrightarrow\operatorname{Grass}_{P(\nu)}(M'_\nu)
\]
de $A_\nu(S')$ en el conjunto de Módulos cocientes localmente libres de
$M'_\nu$ de rango $P(\nu)$ y, por tanto, un homomorfismo funtorial
\[
i_\nu:A_\nu(S')\longrightarrow
\underline{\operatorname{Grass}}_{P(\nu)}(M_\nu)(S'),
\]
donde el funtor de la derecha está representado por el esquema de Grassmann
$\operatorname{Grass}_{P(\nu)}(M_\nu)$ (compárese [4], V), que es proyectivo
sobre $S$. Afirmo:

\result{LEMA}{3.3} El funtor $A_\nu(S')$ es representable, y el morfismo
$Q_\nu\longrightarrow\operatorname{Grass}_{P(\nu)}(M_\nu)$ que representa el
homomorfismo $i_\nu$ es una inmersión (lo que implica que $Q_\nu$ es
cuasiproyectivo sobre $S$).

Esta afirmación equivale a la siguiente (compárese [4], IV). Supongamos dado
un Módulo cociente $N$ de $M'_\nu$, localmente libre de rango $P(\nu)$. Existe
entonces un subpreesquema $Z$ de $S'$ tal que, para todo preesquema localmente
noetheriano $T'$
\src{262}
sobre $S'$, la imagen inversa de $N$ sobre $T'$ pertenece a
$\operatorname{Im}A_\nu(T')$ si y solo si $T'\longrightarrow S'$ se factoriza
a través del subpreesquema $Z$. Cambiando las notaciones, podemos suponer que
$S'=S$; es decir, se nos da un cociente $N_\nu$ de $M_\nu$ por un sub-Módulo
$R_\nu$. Este proviene de un elemento de $A_\nu(S)$ si y solo si satisface las
dos condiciones siguientes:

\medskip
\noindent (i) $M_{\nu+k}/\mathcal S_kR_\nu$ es localmente libre de rango
$P(\nu+k)$ para $k\geq0$.

\smallskip
\noindent (ii) El subhaz $\mathcal K$ de $\mathcal F$ definido por el
sub-Módulo graduado
\[
R_* =\sum_{k\geq0}\mathcal S_kR_\nu
\]
de $M_*$ ([5], II, \S~3), y el cociente
$\mathcal G=\mathcal F/\mathcal K$, satisfacen las condiciones (a) y (b)
anteriores.

Estas condiciones son manifiestamente necesarias. Recíprocamente, si se
satisfacen, el haz $\mathcal G$ definido en (ii), al ser isomorfo al haz
asociado al $\mathcal S_*$-Módulo graduado $N_*$ que es la suma de los
$M_{\nu+k}/\mathcal S_kR_\nu$, es plano sobre $S$ (pues sus fibras son
sumandos directos de localizaciones de $N_*$\footnote{El texto impreso lleva $N$; aquí se localiza el Módulo graduado $N_*$ que acaba de definirse.} en ideales primos homogéneos de
$\mathcal S_*$), y tiene polinomio de Hilbert $P$ por (i). En virtud de (ii),
se satisfacen (a') y (a''), de modo que, para $n\geq\nu$, el homomorfismo
natural $N_n\longrightarrow f_*(\mathcal G(n))$ es un homomorfismo
sobreyectivo entre Módulos localmente libres del mismo rango y, por tanto, un
isomorfismo. Así,
\[
f_*(\mathcal K(\nu+k))=\mathcal S_kR_\nu
\quad\text{para todo $k\geq0$}.
\]
Esto demuestra que $\mathcal G\in A_\nu(S)$ y que $R_\nu$ es el elemento de
$\operatorname{Grass}_{P(\nu)}(M_\nu)$ definido por $\mathcal G$.

El criterio (i)--(ii) también se aplica después de un cambio de base
$S'\longrightarrow S$. Demostraremos primero que la condición (i) después del
cambio de base $S'\longrightarrow S$ puede expresarse exigiendo que
$S'\longrightarrow S$ se factorice a través de cierto subpreesquema $Z$ de
$S$. Una vez hecho esto, sustituir $S$ por $Z$ nos reduce al caso en que (i) ya
se satisface sobre $S$; puesto que es estable por cambio de base, queda por
expresar la condición (ii). Sea $U$ el conjunto de los $s\in S$ tales que la
cohomología en dimensión $>0$ de los haces inducidos sobre la fibra $X_s$ por
$\mathcal G(n)$ y $\mathcal K(n)$ se anula para $n\geq\nu$. Ya hemos observado
que $U$ es abierto, y la condición (ii) se satisface tras un cambio de base
$S'\longrightarrow S$ si y solo si $S'\longrightarrow S$ se factoriza a
través de $U$. Esto demuestra el Lema 3.3. Queda por demostrar el lema
siguiente:

\result{LEMA}{3.4} Sea $S$ un preesquema localmente noetheriano provisto de un
álgebra cuasicoherente graduada positivamente $\mathcal S_*$, generada por
$\mathcal S_1$; sea $M_*$ un $\mathcal S_*$-Módulo graduado cuasicoherente de
tipo finito, $P$ un polinomio con coeficientes racionales y $\nu$ un entero.
Existe entonces un subpreesquema manifiestamente único $Z$ de $S$
\src{263}
con la propiedad siguiente: para todo preesquema $S'$ sobre $S$, el Módulo
$M_n\otimes_{\mathcal O_S}\mathcal O_{S'}$ es localmente libre de rango $P(n)$
para todo $n\geq\nu$ si y solo si $S'\longrightarrow S$ se factoriza a través
de $Z$.

Podemos suponer claramente que $S$ es afín y, por tanto, noetheriano. Entonces:

\result{LEMA}{3.5} Para todo entero $N\geq\nu$, sea $U_N$ el subconjunto
abierto de $S$ formado por los $s\in S$ tales que
\[
\operatorname{rank}_{k(s)}
M_{ns}\otimes_{\mathcal O_{S,s}}k(s)\leq P(n)
\quad\text{para todo $\nu\leq n\leq N$}.
\]
Entonces la sucesión decreciente de abiertos $U_N$ es estacionaria.

Se sabe [3], IV, que $S$ posee una partición finita en subesquemas reducidos
$S_i$ tales que cada $M\otimes_{\mathcal O_S}\mathcal O_{S_i}$ es plano sobre
$S_i$. Podemos, pues, suponer que $M$ es plano y, por tanto, que los $M_n$ son
planos. Finalmente, podemos suponer claramente que $S$ es conexo. Existen
entonces un entero $n_0$ y un polinomio $Q$ tales que
\[
\operatorname{rank}_{k(s)}
M_{ns}\otimes_{\mathcal O_{S,s}}k(s)=Q(n)
\quad\text{para $n\geq n_0$},
\]
[5], III, \S~7. Supongamos primero que $P<Q$\footnote{El texto impreso lleva
$P\neq Q$; la fe de erratas de 1962 prescribe $P<Q$, sustituye el caso «$P=Q$» por
el caso complementario $P(n)\geq Q(n)$ para $n$ grande, y sustituye las
desigualdades estrictas de (*) por desigualdades no estrictas.}, de modo que
$P(n)<Q(n)$ para $n$ grande. Entonces, evidentemente, $U_N=\varnothing$ para
$N$ grande, y la sucesión $U_N$ es estacionaria. En el caso complementario,
$P(n)\geq Q(n)$ para $n$ grande, de modo que $U_N=U_{n_0}$ para $N\geq n_0$,
y la sucesión $U_N$ vuelve a ser estacionaria.

En particular, el conjunto $U_\infty$ de los $s\in S$ tales que
\[
\tag{*}
\operatorname{rank}_{k(s)}
M_{ns}\otimes_{\mathcal O_{S,s}}k(s)\leq P(n)
\quad\text{para todo $n\geq\nu$}
\]
es abierto, pues es la intersección de los $U_N$. Para demostrar (3.4),
podemos, por tanto, sustituir $S$ por el abierto $U_\infty$\footnote{El texto impreso lleva $U$; el abierto considerado aquí es $U_\infty$, la intersección de los $U_N$.}, reduciéndonos al caso en
que la desigualdad (*) se satisface en todo $s\in S$. Además:

\result{LEMA}{3.6} Sea $M$ un Módulo sobre un preesquema localmente
noetheriano $S$, y sea $r$ un entero. Existe entonces un subpreesquema
manifiestamente único $Z$ de $S$ con la propiedad siguiente: para todo $S'$
sobre $S$, el Módulo $M\otimes_{\mathcal O_S}\mathcal O_{S'}$ es localmente
libre de rango $r$ si y solo si $S'\longrightarrow S$ se factoriza a través de
$Z$. Si
\[
\operatorname{rank}_{k(s)}
M_s\otimes_{\mathcal O_{S,s}}k(s)\leq r
\quad\text{para todo $s$},
\]
entonces $Z$ es un subpreesquema cerrado de $S$ (suponiéndose $M$ coherente).
\src{264}
En efecto, el razonamiento anterior nos reduce al caso en que la desigualdad
(*) se satisface para todo $s\in S$ (sustituyendo $S$, si es necesario, por el
subconjunto abierto de los puntos donde se satisface). Podemos entonces suponer
que $M$ interviene en una sucesión exacta
\[
\mathcal O_S^q\longrightarrow\mathcal O_S^r\longrightarrow M\longrightarrow0.
\]
La condición considerada para $S'$ sobre $S$ significa también que, en la
sucesión exacta correspondiente
$\mathcal O_{S'}^q\longrightarrow\mathcal O_{S'}^r
\longrightarrow M'\longrightarrow0$, la segunda flecha es un isomorfismo; es
decir, que la primera es nula. Por consiguiente, el subpreesquema cerrado $Z$
de $S$ definido por el ideal generado por las entradas de la matriz que define
$\mathcal O_S^q\longrightarrow\mathcal O_S^r$ posee la propiedad requerida.

Volvamos ahora a la demostración de (3.4). Sea $Z_n$ el subesquema cerrado de
$S$ asociado por (3.6) al Módulo $M_n$ y al entero $r=P(n)$, y sea $Z'_N$ el
ínfimo de los $Z_n$ para $\nu\leq n\leq N$. Los $Z'_N$ forman una sucesión
decreciente de subesquemas cerrados de $S$\footnote{El texto impreso lleva $Z$; todos los $Z_n$ y sus ínfimos $Z'_N$ son subesquemas cerrados de $S$.} y, por tanto, es necesariamente
estacionaria. Sea $Z$ su valor constante para $N$ grande. Este es el $Z$
buscado en (3.4). Esto termina la demostración de (3.4) y, por consiguiente,
también la de (3.3).

Hemos demostrado, pues, que
$\underline{\operatorname{Quot}}^P_{\mathcal F/X/S}$ está representado por un
$S$-preesquema $Q$ que es unión creciente de subesquemas abiertos $Q_\nu$
cuasiproyectivos sobre $S$. Para ir más lejos, invoquemos el Teorema 2.1, del
cual se deduce fácilmente que $Q$ es cuasicompacto (pues es la imagen del
preesquema $S'$ de tipo finito sobre $S$ que parametriza la familia de haces
cocientes de los $\mathcal F_K$ con polinomio de Hilbert $P$). Así, $Q$ es
igual a uno de los $Q_\nu$ y, por tanto, es cuasiproyectivo sobre $S$. Para
demostrar que es proyectivo sobre $S$, queda por demostrar que es propio sobre
$S$. Para ello basta invocar el criterio valuativo de la propiedad de ser
propio en la forma
[5], II, 7.3.8. Solo hay que verificar:

\result{LEMA}{3.7} Sea $S$ el espectro de un anillo de valoración discreta,
sea $s$ su punto genérico, sea $X$ un preesquema sobre $S$, sea $\mathcal F$
un Módulo cuasicoherente sobre $X$, y sea $\mathcal G_s$ un Módulo
cuasicoherente cociente de
$\mathcal F_s=\mathcal F\otimes_{\mathcal O_S}k(s)$ sobre $X_s$. Existe
entonces un único Módulo cuasicoherente cociente $\mathcal G$ de $\mathcal F$,
plano sobre $S$, cuya restricción a $X_s$ es $\mathcal G_s$.
\src{265}
En efecto, si $\mathcal G_s=\mathcal F_s/\mathcal K_s$, basta tomar el mayor
subhaz $\mathcal K$ de $\mathcal F$ que induce $\mathcal K_s$, [5], I, 9.4.2,
y poner $\mathcal G=\mathcal F/\mathcal K$. Se comprueba fácilmente que este
haz posee las propiedades requeridas.

El Teorema 3.2 y, por consiguiente, (3.1), quedan ahora completamente
demostrados.

La demostración establece al mismo tiempo:

\result{PROPOSICIÓN}{3.8} Bajo las hipótesis de (3.2), sean
$Q=\operatorname{Quot}^P_{\mathcal F/X/S}$, $X_Q=X\times_SQ$ y
$\mathcal F_Q=\mathcal F\otimes_{\mathcal O_S}\mathcal O_Q$; y sea
$\mathcal G$ el cociente coherente de $\mathcal F_Q$, plano sobre $Q$, con
polinomio de Hilbert relativo $P$, tal que $(Q,\mathcal G)$ representa el
funtor $\underline{\operatorname{Quot}}^P_{\mathcal F/X/S}$. Existe un entero
$\nu$ tal que, para $n\geq\nu$, $(f_Q)_*(\mathcal G(n))$ es un Módulo sobre
$Q$, localmente libre de rango $P(n)$ y muy amplio relativamente a $S$; es
decir, define una inmersión de $Q$ en un esquema de Grassmann
$\operatorname{Grass}_{P(n)}(M)$ sobre $S$. En particular, para $n\geq\nu$,
el haz
\[
\bigwedge^{P(n)}(f_Q)_*(\mathcal G(n))
\]
sobre $Q$ es invertible y muy amplio relativamente a $S$.

Como en (3.2), podemos reducirnos al caso en que $\mathcal F$ es plano sobre
$S$; entonces basta elegir un entero $\nu$ tal que $A_\nu=A$ (con las
notaciones anteriores).

La aplicación más importante de (3.2) es el caso
$\mathcal F=\mathcal O_X$. Escribimos entonces
\[
\operatorname{Quot}_{\mathcal O_X/X/S}=\operatorname{Hilb}_{X/S},
\qquad
\operatorname{Quot}^P_{\mathcal O_X/X/S}=\operatorname{Hilb}^P_{X/S},
\]
y obtenemos una descomposición
\[
\operatorname{Hilb}_{X/S}=\coprod_P\operatorname{Hilb}^P_{X/S}.
\]
Por definición, $\operatorname{Hilb}_{X/S}$ representa el funtor
$\underline{\operatorname{Hilb}}_{X/S}(S')$, conjunto de los subesquemas
cerrados de $X'=X\times_SS'$ que son planos sobre $S'$; y
$\operatorname{Hilb}^P_{X/S}$ representa el subfuntor correspondiente de
subesquemas cerrados con polinomio de Hilbert dado $P$. Estos preesquemas se
llaman también, respectivamente, \emph{preesquema de Hilbert} de $X$ sobre
$S$ y \emph{preesquema de Hilbert de índice $P$}. La terminología se justifica
por el papel desempeñado en la teoría por los polinomios de Hilbert. Su
diferencia de naturaleza respecto de las variedades de Chow clásicas (que
parametrizan ciclos y no variedades) es del mismo orden que la existente entre
\src{266}
el anillo de Chow de las clases de ciclos sobre una variedad y el anillo de
clases de haces sobre la variedad (tal como aparece en el teorema de
Riemann--Roch [1]). En efecto, cuando $X=\mathbf P^r_S$ y $S$ es el espectro
de un cuerpo, conocer el polinomio de Hilbert de un Módulo coherente
$\mathcal F$ sobre $X$ equivale también a conocer las clases de Chern de
$\mathcal F$, o la clase de $\mathcal F$ en el anillo de clases de haces
coherentes sobre $X$.

\result{OBSERVACIONES}{3.9} Obsérvese asimismo que la construcción de
$\operatorname{Quot}_{\mathcal F/X/S}$ y
$\operatorname{Quot}^P_{\mathcal F/X/S}$ se redujo al caso
$X=\mathbf P^r_S$, $\mathcal F=\mathcal O_X^N$, con $\mathcal O_X(1)$ el haz
muy amplio usual; más precisamente, los
$\operatorname{Quot}_{\mathcal F/X/S}$ generales aparecen como subesquemas
cerrados de estos. Como la formación de los
$\operatorname{Quot}^P_{\mathcal F/X/S}$ es manifiestamente compatible con
los cambios de base $S'\longrightarrow S$, nos reducimos además a
$S=\operatorname{Spec}(\mathbf Z)$. En definitiva, pues, se reduce al estudio
de los esquemas proyectivos siguientes sobre $\mathbf Z$:
\[
\mathfrak Q^P_{r,N}=
\operatorname{Quot}^{P}_{(\mathcal O_{\mathbf P^r_{\mathbf Z}})^N/
\mathbf P^r_{\mathbf Z}/\operatorname{Spec}(\mathbf Z)},
\]
y, más particularmente, de los esquemas de Hilbert absolutos
\[
\operatorname{Hilb}^P_r=\mathfrak Q^P_{r,1}.
\]
Sería muy interesante un estudio más detallado de estos esquemas, comenzando
por la determinación de sus componentes conexas (¿son conexos?),\footnote{La
fe de erratas de 1962 señala que Hartshorne demostró la conexidad de los
$\operatorname{Hilb}^P_r$ sobre un cuerpo algebraicamente cerrado y determinó
los pares $(r,P)$ para los cuales
$\operatorname{Hilb}^P_r\neq\varnothing$.} y de sus componentes irreducibles
(por Serre [7], puede haber componentes irreducibles situadas enteramente sobre
un primo no nulo $p$). Recordemos la pregunta de Weil de si las componentes
irreducibles de las fibras de $\operatorname{Hilb}^P_r$ sobre
$s\in\operatorname{Spec}(\mathbf Z)$ corresponden a extensiones «regulares»
del cuerpo primo, es decir, si son «relativamente conexas». Estas cuestiones
quizá sean más abordables para los esquemas de Hilbert que para las
«variedades de Chow».

\fsection{4}{Variantes}

\noindent a. Bajo las hipótesis de (3.1), sea $U$ un subconjunto abierto de
$X$, y sea $A'$ el subfuntor de
$A=\underline{\operatorname{Quot}}_{\mathcal F/X/S}$ tal que $A'(S')$ es el
conjunto de Módulos cocientes $\mathcal G$ de $\mathcal F'$, planos sobre
$S'$, cuyo soporte está contenido en $U'$. Manifiestamente, $A'$ está
representado por un subconjunto abierto del preesquema
$\operatorname{Quot}_{\mathcal F/X/S}$ que representa $A$. De ello resulta
que los Teoremas (3.1) y (3.2)
\src{267}
siguen siendo válidos cuando se supone que $X$ es cuasiproyectivo sobre $S$ en
lugar de proyectivo sobre $S$, siempre que en las conclusiones también se
sustituya «proyectivo» por «cuasiproyectivo» y que
$\underline{\operatorname{Quot}}_{\mathcal F/X/S}(S')$ denote ahora el
conjunto de cocientes coherentes $\mathcal G$ de $\mathcal F'$, planos sobre
$S'$, cuyo soporte es propio sobre $S'$.

\medskip
\noindent b. Más generalmente, pueden imponerse a los cocientes $\mathcal G$
de $\mathcal F'=\mathcal F\otimes_{\mathcal O_S}\mathcal O_{S'}$, planos
sobre $S'$, toda clase de condiciones naturales adicionales estables por
cambio de base, obteniéndose así subfuntores de
$\underline{\operatorname{Quot}}_{\mathcal F/X/S}$ cuya representabilidad
puede buscarse. En muchos casos, el criterio usual demuestra que estos
funtores vuelven a estar representados por subconjuntos abiertos de
$\operatorname{Quot}_{\mathcal F/X/S}$. Esto ocurre, en particular, cuando se
impone una de las propiedades adicionales siguientes:

\smallskip
\noindent $1^\circ$ Las dimensiones de los ciclos primos asociados a los
Módulos $\mathcal G_{s'}$ ($s'\in S'$) inducidos sobre las fibras $X'_{s'}$
pertenecen a un conjunto dado de enteros.

\smallskip
\noindent $2^\circ$ (Cuando $\mathcal F=\mathcal O_X$, de modo que
$\mathcal G$ corresponde a un subpreesquema cerrado $Y$ de $X'$): $Y$ es un
preesquema liso ([3], IV) sobre $S'$, respectivamente normal sobre $S'$ (es
decir, las fibras $Y_{s'}$ son normales «sobre $k(s')$», lo que significa
normales después de toda extensión del cuerpo base), respectivamente (cuando
$X$ es plano sobre $S$) una intersección completa local relativa en $X'$ sobre
$S'$ (es decir, las fibras $Y_{s'}$ son intersecciones completas locales en los
$X'_{s'}$).\footnote{El texto impreso omite las primas en estas propiedades de $Y$ sobre la base variable $S'$ y dentro de $X'$. La hipótesis relativa al $X$ original plano sobre $S$ permanece sin primas.}

Otras condiciones harían intervenir propiedades cohomológicas de los Módulos
$\mathcal G_{s'}$ inducidos sobre los $X'_{s'}$, etc. Por supuesto, una
conjunción de condiciones, cada una de ellas representada por un subconjunto
abierto $U_i$ de $\operatorname{Quot}_{\mathcal F/X/S}$, está representada por
la intersección de estos subconjuntos abiertos. Por ejemplo, asignar a todo
$S'$ sobre $S$ el conjunto de subesquemas cerrados $Y$ de
$X'=X\times_SS'$ que son recubrimientos étale ([3], I) de $S'$ de rango dado
$r$ da un funtor contravariante representable de $S'$.

\medskip
\noindent c. Los preesquemas $\operatorname{Hom}_S(X,Y)$,
$\Pi_{X/S}(Z/X)$ y $\operatorname{Isom}_S(X,Y)$ definidos en [2], II, C,
n.º 2 existen bajo hipótesis de proyectividad convenientes y aparecen como
subconjuntos abiertos de preesquemas de Hilbert adecuados. Puesto que
\[
\operatorname{Hom}_S(X,Y)=\Pi_{X/S}((X\times_SY)/X),
\]
el caso de $\operatorname{Hom}_S(X,Y)$ se reduce al de
$\Pi_{X/S}(Z/X)$. Para todo $S'$ sobre $S$, el conjunto de secciones de
$Z'=Z\times_SS'$ sobre $X'=X\times_SS'$ está en correspondencia biyectiva con
el conjunto de subpreesquemas $\Gamma$ de $Z'$ (necesariamente cerrados si $Z$
es separado sobre $X$) tales que el morfismo $\Gamma\longrightarrow X'$
inducido por $Z'\longrightarrow X'$ sea un isomorfismo. De este modo,
\src{268}
cuando $X$ es plano y propio sobre $S$ y $Z$ es cuasiproyectivo sobre $S$,
$\Pi_{X/S}(Z/X)$ existe y aparece como subpreesquema abierto de
$\operatorname{Hilb}_{Z/S}$. Por consiguiente, cuando $X$ es proyectivo y
plano sobre $S$ e $Y$ es cuasiproyectivo sobre $S$,
$\operatorname{Hom}_S(X,Y)$ existe y aparece como subpreesquema abierto de
$\operatorname{Hilb}_{(X\times_SY)/S}$. Cuando $X$ e $Y$ son ambos
proyectivos sobre $S$, se sigue inmediatamente que
$\operatorname{Isom}_S(X,Y)$ también existe y está representado por un
subconjunto abierto de $\operatorname{Hom}_S(X,Y)$. Análogamente, cuando $X$
es plano y proyectivo sobre $S$ e $Y$ es cuasiproyectivo sobre $S$, el
$S$-preesquema $\operatorname{Imm}_S(X/Y)$ correspondiente al subfuntor de
$\operatorname{Hom}_S(X,Y)$ cuyos elementos son los $S'$-homomorfismos
$X'\longrightarrow Y'$ que son inmersiones también está representado por un
subconjunto abierto de $\operatorname{Hom}_S(X,Y)$.

Sea $\mathcal L$ (respectivamente, $\mathcal M$) un haz invertible sobre $X$
(respectivamente, $Y$), muy amplio relativamente a $S$. Entonces
$\mathcal L\otimes_{\mathcal O_S}\mathcal M$ es un haz muy amplio sobre
$X\times_SY$ relativamente a $S$. Por consiguiente, para todo polinomio $P$
con coeficientes racionales, $\operatorname{Hilb}^P_{(X\times_SY)/S}$ está
definido y es un preesquema cuasiproyectivo sobre $S$. Induce, por tanto,
sobre $\operatorname{Hom}_S(X,Y)$ un subconjunto abierto y cerrado,
cuasiproyectivo sobre $S$, que denotaremos por
$\operatorname{Hom}_S(X,Y)^P$. Así, las secciones de
$\operatorname{Hom}_S(X,Y)^P$ sobre $S$ son los $S$-morfismos
$g:X\longrightarrow Y$ tales que, para todo entero $n$,
\[
\chi\!\left(\left(\mathcal L\otimes_{\mathcal O_X}g^*(\mathcal M)\right)^{
\otimes n}\right)=P(n).
\]

De este modo se obtienen generalizaciones del teorema de Matsusaka que afirma
que los automorfismos de una variedad proyectiva «polarizada» forman un grupo
algebraico. Aquí la afirmación tiene un sentido manifiestamente más preciso,
pues este grupo se define como solución de un problema universal. Obsérvese
también que, sobre un cuerpo algebraicamente cerrado, el grupo de
automorfismos considerado anteriormente se obtiene del «verdadero» grupo
definido aquí pasando al cociente por los elementos nilpotentes. Esto explica
por qué las construcciones antiguas tienen pocas posibilidades de funcionar
sobre un cuerpo base no perfecto: el ideal de elementos nilpotentes que
aparece después de una extensión del cuerpo base no tiene por qué estar
«definido sobre $k$». La misma observación se aplica a la mayoría de las
construcciones de estilo antiguo.

\fsection{5}{Estudio infinitesimal de los esquemas de Hilbert}

Se deduce del resultado siguiente:

\src{269}
\result{PROPOSICIÓN}{5.1} Sea $S$ un preesquema, sea $S_0$ un subpreesquema
definido por un ideal cuasicoherente $\mathcal J$ de cuadrado nulo, sea $X$ un
$S$-preesquema y sea $\mathcal F$ un Módulo cuasicoherente sobre $X$. Pongamos
$X_0=X\times_SS_0$ y
$\mathcal F_0=\mathcal F\otimes_{\mathcal O_S}\mathcal O_{S_0}$, y sea
$\mathcal G_0=\mathcal F_0/\mathcal K_0$ un Módulo cuasicoherente cociente de
$\mathcal F_0$, plano sobre $S_0$. Para todo subconjunto abierto $U$ de $X$,
sea $\mathcal E(U)$ el conjunto de Módulos cuasicoherentes cocientes
$\mathcal G$ de $\mathcal F|U$, planos sobre $S$, tales que
$\mathcal G\otimes_{\mathcal O_S}\mathcal O_{S_0}=\mathcal G_0$. Al variar
$U$, los $\mathcal E(U)$ son, pues, las secciones de un haz $\mathcal E$ sobre
$X$.\footnote{El texto impreso lleva $U$; al variar los abiertos $U$ de $X$, los conjuntos $\mathcal E(U)$ definen un haz sobre $X$.} El haz de grupos
\[
\mathcal A=\operatorname{Hom}_{\mathcal O_{X_0}}
(\mathcal K_0,\mathcal G_0\otimes_{\mathcal O_{S_0}}\mathcal J)
\]
actúa naturalmente sobre $\mathcal E$, que se convierte así en un haz
«formalmente homogéneo principal bajo $\mathcal A$» (es decir, para todo
abierto $U$ de $X$, $\mathcal E(U)$ es vacío o es un conjunto homogéneo
principal bajo $\mathcal A(U)$).

De ello se sigue que:

\result{COROLARIO}{5.2} Supongamos que localmente sobre $X$ existe una
extensión $\mathcal G$ de $\mathcal G_0$ a un cociente de $\mathcal F$ plano
sobre $S$ (es decir, que las fibras del haz $\mathcal E$ no son vacías).
Existe entonces una clase de obstrucción canónica
\[
c(\mathcal G_0)\in H^1(X,\mathcal A)
\]
cuya anulación es necesaria y suficiente para que exista una extensión global
$\mathcal G$ de $\mathcal G_0$ a un cociente de $\mathcal F$ plano sobre $S$.
Si esta clase se anula, el conjunto $\mathcal E(X)$ de todas las extensiones
posibles es un conjunto homogéneo principal bajo
\[
\mathcal A(X)=\operatorname{Hom}_{\mathcal O_{X_0}}
(\mathcal K_0,\mathcal G_0\otimes_{\mathcal O_{S_0}}\mathcal J).
\]
\footnote{En esta última fórmula, el texto impreso lleva $\mathcal O_X$ en
lugar de $\mathcal O_{X_0}$, aunque la definición inmediatamente anterior de
$\mathcal A$ fija $\mathcal O_{X_0}$.}
En particular, existe una extensión global si $H^1(X,\mathcal A)=0$.

\result{COROLARIO}{5.3} Supongamos que
$Q=\operatorname{Quot}_{\mathcal F/X/S}$ existe (cf. 4 (a)); por ejemplo, que
$X$ es cuasiproyectivo sobre un $S$ localmente noetheriano y que $\mathcal F$
es coherente. Sea $x\in Q$ correspondiente a una extensión residual
$K=k(x)$ de $k(s)$ ($s\in S$):
\src{270}
así, $x$ está definido por un Módulo coherente cociente
$\mathcal G_0=\mathcal F_0/\mathcal K_0$ del Módulo
$\mathcal F_0=\mathcal F_K$ sobre el $K$-preesquema $X_K$. Sea $\mathcal A$
el haz coherente sobre $X_K$ definido por
\[
\mathcal A=\operatorname{Hom}_{\mathcal O_{X_0}}(\mathcal K_0,\mathcal G_0).
\]
Entonces el espacio tangente de Zariski a la fibra $Q_s$ en $x$ (el dual
sobre $K$ de $\mathfrak m/\mathfrak m^2$, donde $\mathfrak m$ es el ideal
maximal de $\mathcal O_{Q_K,x}$) es canónicamente isomorfo a
$H^0(X_K,\mathcal A)$.

La descripción del espacio tangente de Zariski admite la generalización
siguiente. Dado un $S$-morfismo $g:S'\longrightarrow Q$, es decir, una sección
$g'$ de $Q'=Q\times_SS'$ sobre $S'$, el Módulo
\[
\Omega=g^*(\Omega^1_{Q/S})=g'^*(\mathcal J/\mathcal J^2)
\]
(donde $\mathcal J$ es el ideal sobre $Q'$ definido por la sección $g'$ de
$Q'$ sobre $S'$) está caracterizado por la fórmula siguiente, funtorial en el
Módulo coherente $\mathcal M$ sobre $S'$:
\[
\operatorname{Hom}_{\mathcal O_{S'}}(\Omega,\mathcal M)
\simeq H^0(X',\mathcal A),
\]
donde $\mathcal A$ es el Módulo sobre $X'=X\times_SS'$ definido por
\[
\mathcal A=\operatorname{Hom}_{\mathcal O_{X'}}
(\mathcal K,\mathcal G\otimes_{\mathcal O_{S'}}\mathcal M).
\]
\footnote{El texto impreso pone $\mathcal O_S$ en este producto tensorial;
aquí $\mathcal G$ y $\mathcal M$ son Módulos sobre $S'$, y la construcción de
deformación de cuadrado nulo exige el producto tensorial sobre
$\mathcal O_{S'}$.}
Aquí $\mathcal G=\mathcal F'/\mathcal K$ es el Módulo cociente de
$\mathcal F'=\mathcal F\otimes_{\mathcal O_S}\mathcal O_{S'}$ que corresponde
a $g$. En efecto, aplíquese (5.1) sustituyendo $S_0$ por $S'$ y $S$ por el
preesquema $D(\mathcal M)=(S',\mathcal O_{S'}+\mathcal M)$, donde
$\mathcal M$ se considera como un ideal de cuadrado nulo.

Cuando $\mathcal F=\mathcal O_X$ en (5.1), el dato de $\mathcal G_0$
corresponde a un subpreesquema cerrado $Y_0$ de $X_0$, plano sobre $S_0$,
definido por el ideal $\mathcal J_0=\mathcal K_0$. La fórmula (*) da entonces
\[
\mathcal A=\operatorname{Hom}_{\mathcal O_{X_0}}
(\mathcal J_0/\mathcal J_0^2,
\mathcal O_{Y_0}\otimes_{\mathcal O_{S_0}}\mathcal J),
\]
donde $\mathcal J_0/\mathcal J_0^2$\footnote{El texto impreso omite el subíndice $0$ en las dos menciones del haz conormal y en $\mathcal J_0^{-1}$ más abajo. El Ideal de $Y_0$ en $X_0$ es $\mathcal J_0$; el Ideal $\mathcal J$ de la base es distinto.} se interpreta como el haz conormal de $Y_0$
en $X_0$, denotado también por
\src{271}
$\check{\mathcal N}_{Y_0/X_0}$. Conviene entonces considerar $\mathcal A$
como un Módulo sobre $Y_0$ y calcular $H^0$ y $H^1$ sobre $Y_0$. Además, si
$Y_0$ es localmente una intersección completa en $X_0$, con $X$ plano sobre
$S$, queda garantizada la existencia de extensiones locales en (5.1); asimismo,
$\mathcal J_0/\mathcal J_0^2$ es localmente libre sobre $Y_0$, y puede escribirse
\[
\mathcal A=\mathcal N_{Y_0/X_0}\otimes_{\mathcal O_{S_0}}\mathcal J,
\]
donde el primer factor de la derecha es el haz normal de $Y_0$ en $X_0$.
Utilizando el criterio fundamental de lisura ([3], III, 3.1), se obtiene, por
ejemplo:

\result{COROLARIO}{5.4} Bajo las hipótesis de (5.3), supongamos que
$\mathcal F=\mathcal O_X$, que $X$ es plano sobre $S$ y que el subpreesquema
cerrado $Y_0$ de $X_0$ correspondiente a $\mathcal G_0$ es localmente una
intersección completa. Entonces el espacio tangente de Zariski a $Q_s$ en $x$
es canónicamente isomorfo a $H^0(Y_0,\mathcal N_{Y_0/X_0})$. Si
$H^1(Y_0,\mathcal N_{Y_0/X_0})=0$, entonces el preesquema de Hilbert
$\operatorname{Hilb}_{X/S}$ es liso sobre $S$ en $x$.
($\mathcal N_{Y_0/X_0}$ es el haz normal de $Y_0$ en $X_0$.)

\result{OBSERVACIÓN}{5.5} Este enunciado se aplica, en particular, cuando
$Y_0$ es una intersección completa en $X_0$ definida por una ecuación, es
decir, un «divisor de Cartier» positivo. Entonces $\mathcal N_{Y_0/X_0}$ es
isomorfo al haz sobre $Y_0$ inducido por el haz invertible
$\mathcal J_0^{-1}$ sobre $X_0$ definido por el divisor $Y_0$. Esta situación
se presenta especialmente en el estudio de familias de divisores positivos
sobre una variedad proyectiva no singular $X_0$. El isomorfismo entre el
espacio tangente de Zariski en el punto $x$ de $Q$ (o, equivalentemente, del
subconjunto abierto $D$ de $Q$ correspondiente a los divisores) y
$H^0(Y_0,\mathcal N_{Y_0/X_0})$ era conocido en geometría algebraica clásica
como el «homomorfismo característico» del primero al segundo. Solo estaba
definido cuando $x$ era un punto liso de la variedad de parámetros $T$ de una
«familia continua completa» de divisores, es decir, desde nuestro punto de
vista, de una componente irreducible del esquema $D$ provisto de su estructura
reducida inducida. El
espacio tangente a $T$ en $x$ es entonces un subespacio del espacio tangente a
$D$ en $x$, por lo que el homomorfismo característico clásico es, en efecto,
inyectivo, pero solo es sobreyectivo bajo condiciones adicionales; por ejemplo,
cuando $D$ es íntegro en $x$. De hecho, Zappa [8] construyó un ejemplo
\src{272}
(con $X$ una superficie proyectiva no singular sobre los números complejos) en
el que el homomorfismo característico no es sobreyectivo ni siquiera en el
punto genérico de $T$. Así, $D$ no es íntegro ni siquiera en el punto genérico
de la componente irreducible considerada. Esto demuestra de manera
especialmente llamativa hasta qué punto las variedades con elementos
nilpotentes son necesarias para comprender fenómenos de la teoría más clásica
de superficies.\footnote{La fe de erratas de 1962 añade que Mumford acababa de
construir una componente irreducible del esquema de Hilbert de
$\mathbf P^3_{\mathbf C}$ cuyos puntos generales representan curvas no
singulares de grado 14 y género 24, que es no reducida en su punto genérico;
al hacer explotar una de estas curvas se obtiene también una variedad triple
proyectiva regular cuyo esquema formal de módulos es genéricamente no
reducido.}

\subsection*{5.6}

En (5.4) dimos un criterio de lisura aplicable, en particular, a esquemas de
divisores. Kodaira dio en [6] un criterio diferente, a saber, la anulación de
$H^1(X_0,\mathcal L)$, donde $\mathcal L=\mathcal J_0^{-1}$ es el haz
invertible sobre $X_0$ definido por el divisor $Y_0$. Este criterio es válido
cuando $S$ es el espectro de un cuerpo de característica cero, y fue demostrado
en [6] por métodos trascendentes cuando el cuerpo base es $\mathbf C$. Más
generalmente, siendo de nuevo $S$ arbitrario, la condición de Kodaira basta
para que el morfismo canónico
\[
D\longrightarrow\operatorname{Pic}_{X/S}
\]
del preesquema de divisores al preesquema de Picard de $X/S$ sea liso en el
punto $x$ considerado (como se deduce fácilmente del criterio usual de lisura
una vez conocida la existencia de $\operatorname{Pic}_{X/S}$). Si, además,
$\operatorname{Pic}_{X/S}$ es liso sobre $S$ en la imagen de $x$ (por ejemplo,
si $\operatorname{Pic}_{X/S}$ es liso sobre $S$), entonces $D$ es liso sobre
$S$ en $x$. Por otra parte, Cartier demostró que todo preesquema en grupos
localmente de tipo finito sobre un cuerpo $k$ de característica cero es liso
sobre $k$. Combinando estos dos resultados se recupera el resultado de
Kodaira. Estas observaciones demuestran también que, sobre un cuerpo $K$ de
característica $p>0$, cuando $\operatorname{Pic}_{X/S}$ no es liso sobre $K$\footnote{El texto impreso lleva $k$; el cuerpo base de esta frase en característica positiva es $K$.}
(como cuando $X$ es la superficie de Igusa), la condición
$H^1(X_0,\mathcal L)=0$ implica, por el contrario, que $D$ no es liso en $x$,
y que incluso es no reducido en $x$ si $K$ es algebraicamente cerrado.

Finalmente, registremos el resultado siguiente, que desempeña un papel
importante en el estudio infinitesimal de los espacios fibrados:

\result{PROPOSICIÓN}{5.7} Sea $X$ un preesquema finito y plano sobre un $S$
localmente noetheriano, y sea $Z$ un preesquema sobre $X$ tal que
$\Pi_{X/S}(Z/X)$ exista (como ocurre si $Z$ es cuasiproyectivo sobre $X$). Si
$Z$ es liso sobre $X$, entonces $\Pi_{X/S}(Z/X)$ es liso sobre $S$.

Esto resulta inmediatamente de la definición y del criterio usual de lisura
([3], III, 3.1). Obsérvese que, cuando $X$ es finito y plano sobre $S$, la
existencia de $\Pi_{X/S}(Z/X)$\footnote{El texto impreso lleva aquí $Z/K$; no
se ha definido ningún $K$, y el funtor de secciones considerado desde (5.7) es
$\Pi_{X/S}(Z/X)$.} puede tratarse de manera muy elemental, sin
\src{273}
utilizar la teoría de los esquemas de Hilbert. Por ejemplo, si $X$ es radicial
sobre $S$, entonces $\Pi_{X/S}(Z/X)$ existe sin ninguna restricción sobre $Z$.
Como otro ejemplo, sea $T$ un $S$-preesquema y sea $T_n$ el «entorno
infinitesimal de orden $n$» de la diagonal de $T\times_ST$ en $T\times_ST$,
provisto de los morfismos $p_1,p_2:T_n\longrightarrow T$ inducidos por las dos
proyecciones. Consideremos $T_n$ como un preesquema finito sobre $T$ por medio
de $p_1$, y supongamos además que es plano sobre $T$ (como ocurre cuando $T$
es liso sobre $S$). Para todo preesquema $X$ sobre $T$, pongamos
\[
(X/T/S)^{(n)}=
\Pi_{T_n/T}\bigl(p_2^*(X/T)/T_n\bigr).
\]
\footnote{El texto impreso coloca $T_n/S$ bajo el símbolo $\Pi$. La
Proposición 5.7 se aplica aquí al morfismo finito y plano
$p_1:T_n\to T$, y el resultado es explícitamente un preesquema sobre $T$; por
tanto, el denominador es $T_n/T$.}
Este es un preesquema sobre $T$ llamado \emph{fibrado de gérmenes de secciones
de orden $n$ de $X$ sobre $T$} (relativo a $S$). Depende funtorialmente de $X$
y es liso sobre $T$ si $X$ lo es.

\fsection{6}{Relaciones con la norma y los productos simétricos}

Sea $S$ un preesquema, sean $X$ e $Y$ $S$-preesquemas, y sea
\[
u:(X/S)^n\longrightarrow Y
\]
un $S$-morfismo simétrico de la $n$-ésima potencia cartesiana de $X/S$ en
$Y$. Para simplificar, supongamos que $S$ es localmente noetheriano y que $X$
e $Y$ son de tipo finito sobre $S$. A todo Módulo coherente $\mathcal F$ sobre
$X$ cuyo soporte es finito sobre $S$, que es plano sobre $S$ y tiene rango
relativo $n$ sobre $S$ (es decir, $f_*(\mathcal F)$ es un Módulo localmente
libre de rango $n$ sobre $S$), puede asociársele entonces de manera natural una
sección de $Y$ sobre $S$:
\[
\pi^u_{X/S}(\mathcal F)\in\Gamma(Y/S).
\]
No daremos aquí la definición formal, sino que nos limitaremos a observar que
el formalismo resultante es una generalización natural del formalismo usual de
normas y trazas. Para todo $X/S$ separado, la ley de norma tiene como blanco canónico el espacio de potencias divididas $\Gamma^n_S(X)$. Si $X$ es plano sobre $S$ y existe la $n$-ésima potencia simétrica (por ejemplo, cuando las órbitas de $\mathfrak S_n$ en $(X/S)^n$ están contenidas en abiertos afines), el morfismo canónico $\operatorname{Symm}^n_S(X)\to\Gamma^n_S(X)$ es un isomorfismo, y se obtiene un elemento canónico
\[
\pi_{X/S}(\mathcal F)\in\Gamma(\operatorname{Symm}^n_S(X)/S),
\]
\src{274}
del cual pueden recuperarse los $\pi^u_{X/S}(\mathcal F)$. Otro caso
importante es aquel en que $X$ es un monoide conmutativo sobre $S$, $X=Y$, y
$u$ procede de la ley de composición de $X$. Escribimos entonces simplemente
$\pi(\mathcal F)$ para la sección de $X$ sobre $S$ asociada al Módulo
$\mathcal F$ sobre $X$.

Supongamos ahora que $\mathcal F$ es un Módulo coherente sobre $X$ tal que
$\operatorname{Quot}_{\mathcal F/X/S}$ existe, o al menos tal que el funtor
$\underline{\operatorname{Quot}}^n_{\mathcal F/X/S}$ que asigna a todo $S'$
sobre $S$ el conjunto de haces cocientes coherentes $\mathcal M$ de
$\mathcal F'=\mathcal F\otimes_{\mathcal O_S}\mathcal O_{S'}$ que son planos
sobre $S'$ y tienen rango relativo $n$ está representado por un $S$-preesquema
$\operatorname{Quot}^n_{\mathcal F/X/S}$. (Cuando $X$ es cuasiproyectivo
sobre $S$, $\operatorname{Quot}^n_{\mathcal F/X/S}$ existe y, con las
notaciones del Apartado 3, es $\operatorname{Quot}^P_{\mathcal F/X/S}$, donde
$P$ es el polinomio constante $n$.) Como la formación de
$\pi^u_{X/S}(\mathcal M)$ es compatible con el cambio de base, se obtiene un
morfismo canónico
\[
\pi^u_{X/S}:\operatorname{Quot}^n_{\mathcal F/X/S}\longrightarrow Y,
\]
y, en general,
\[
\pi_{X/S}:\operatorname{Quot}^n_{\mathcal F/X/S}
\longrightarrow\Gamma^n_S(X).
\]
Cuando $X$ es plano sobre $S$, identificamos este blanco con $\operatorname{Symm}^n_S(X)$.
Independientemente de la identificación precedente, cuando $X$ es separado sobre $S$, la construcción mediante la familia finita y plana para $\mathcal F=\mathcal O_X$ da un morfismo
\[
\pi_{X/S}:\operatorname{Hilb}^n_{X/S}
\longrightarrow\operatorname{Symm}^n_S(X).
\]
Este es manifiestamente un isomorfismo para $n=0$ o $n=1$. Pero para
$n>1$\footnote{El texto impreso repite $n\geq1$, contradiciendo el
isomorfismo que acaba de afirmarse para $n=1$.}, incluso cuando $S$ es el
espectro de un cuerpo $k$ y $X$ es liso sobre $S$, no es en general un
isomorfismo, ni siquiera un morfismo inyectivo: un subesquema de dimensión cero
de $X$ (correspondiente, por ejemplo, a un ideal $\mathfrak i$ primario para el
ideal maximal en un anillo local $\mathcal O_{X,x}$, para un punto cerrado $x$
de $X$) no está determinado por el ciclo que define (en este caso, por la
codimensión sobre $k$ de $\mathfrak i$ en $\mathcal O_{X,x}$). Solo puede
afirmarse lo siguiente, siendo de nuevo $S$ arbitrario:

\medskip
\noindent a. Si $X$ es liso sobre $S$, el morfismo norma induce un isomorfismo
del subconjunto abierto de $\operatorname{Hilb}^n_{X/S}$ que clasifica los
recubrimientos étale de rango $n$ contenidos en $X$ (cf. n.º 4, (b)) sobre el
subconjunto abierto de $\operatorname{Symm}^n_S(X)$ correspondiente a los
$n$-ciclos sin componentes múltiples.
\src{275}

\medskip
\noindent b. Si, además, $X$ tiene dimensión relativa 1 sobre $S$, entonces el
morfismo norma es un isomorfismo de $\operatorname{Hilb}^n_{X/S}$ sobre
$\operatorname{Symm}^n_{X/S}$.

Esto se debe a que un subesquema de dimensión cero de una curva algebraica no
singular está determinado por su ciclo correspondiente. La misma observación
se aplica, más generalmente, a los divisores de Cartier positivos sobre un
esquema algebraico no singular (y es posible que, en este caso tan especial,
la variedad de Chow dé el mismo objeto que el esquema de Hilbert).

\fsection{7}{Complementos y cuestiones}

Como observó J.-P. Serre, un ejemplo bien conocido de Nagata demuestra que
pueden encontrarse un esquema $S$, espectro de un cuerpo $k$; un $S$-esquema
$S'$, espectro de una extensión cuadrática $k'$ de $k$; y un $S'$-esquema
propio y liso (pero no proyectivo) $X$ de dimensión 3 tal que
$\Pi_{S'/S}(X/S')$ no existe. Se sigue, a fortiori, que el esquema de
Hilbert $\operatorname{Hilb}^2_{X/S}$ no existe (ni siquiera el $k$-esquema
que representaría los recubrimientos étale de $S$ de rango 2 contenidos en
$X$, y a fortiori tampoco el cuadrado simétrico de $X$; cf. el apartado
anterior). Esto impone serias limitaciones a la posibilidad de realizar
construcciones no proyectivas en geometría algebraica. (Sin embargo, es
plausible que tales limitaciones no aparezcan en geometría analítica, como
tampoco aparecen en geometría formal (cf. [2], II).) Por otra parte, si $X$ es
un esquema propio sobre $S=\operatorname{Spec}k$ y $Z$ es cuasiproyectivo
sobre $X$, entonces $\Pi_{X/S}(Z/X)$ existe y es un esquema que es unión
disjunta de una sucesión de esquemas cuasiproyectivos sobre $S$ (como en el
caso proyectivo (3.1)). En efecto, se reduce al caso en que el propio $X$ es
proyectivo dominando $X$ por un $S$-esquema proyectivo $X'$. No daremos los
detalles de la demostración, que utiliza también el resultado de factorización
para un morfismo finito recogido en [2], I, A, n.º 2 (b). El método funciona
porque, siendo $S$ el espectro de un cuerpo, el $X'$ que interviene en el lema
de Chow es automáticamente plano sobre $S$.

No sé si el resultado sigue siendo cierto sin hipótesis sobre $S$, suponiendo
únicamente que $X$ es propio y plano sobre $S$ y que $Z$ es cuasiproyectivo
sobre $X$. Un caso importante en las aplicaciones es aquel en que $Z$ es un
subesquema cerrado de $X$; si $\Pi_{X/S}(Z/X)$ existe, es necesariamente un
subesquema cerrado de $S$. Puede
\src{276}
construirse directa y muy sencillamente cuando $X$ es proyectivo sobre $S$,
sin utilizar la teoría de los esquemas de Hilbert. Más generalmente, el mismo
método demuestra que, si $Z$ es afín sobre $X$, entonces $\Pi_{X/S}(Z/X)$
existe y es afín sobre $S$. También demuestra que, si $X$ es propio y plano
sobre $S$ (no necesariamente proyectivo sobre $S$), entonces, para todo
fibrado vectorial localmente trivial $Z$ sobre $X$, $\Pi_{X/S}(Z/X)$ existe y
es un fibrado vectorial sobre $S$. Sería deseable retomar y unificar estos
resultados.

\section*{BIBLIOGRAFÍA}
\addcontentsline{toc}{section}{Bibliografía}

\begin{description}
\item[{[1]}] BOREL (A.) y SERRE (J.-P.). -- Le théorème de Riemann--Roch,
\emph{Bull. Soc. math. France}, vol. 86, 1958, pp. 97--136.

\item[{[2]}] GROTHENDIECK (Alexander). -- Technique de construction et
théorèmes d'existence en géométrie algébrique, I, II, III, \emph{Séminaire
Bourbaki}, vol. 13, 1959/60, n.º 190, 29 pp., n.º 195, 22 pp.; y vol. 14,
1960/61, n.º 212, 20 pp.

\item[{[3]}] GROTHENDIECK (Alexander). -- Séminaire de Géométrie algébrique,
I, II, III, IV. -- París, Institut des hautes Études scientifiques, 1960/61
(mimeografiado).

\item[{[4]}] GROTHENDIECK (Alexander). -- Technique de construction en
géométrie analytique, IV, V, \emph{Séminaire Cartan}, vol. 13, 1960/61, n.os 11
y 12.

\item[{[5]}] GROTHENDIECK (A.) y DIEUDONNÉ (J.). -- \emph{Éléments de
géométrie algébrique}, I: Le langage des schémas. -- París, Presses
universitaires de France, 1960 (Institut des hautes Études scientifiques,
Publications mathématiques, 4); II, III, IV (de próxima aparición).

\item[{[6]}] KODAIRA (K.). -- Characteristic linear systems of complete
continuous systems, \emph{Amer. J. of Math.}, vol. 78, 1956, pp. 716--744.

\item[{[7]}] SERRE (Jean-Pierre). -- Exemples de variétés projectives en
caractéristique $p$ non relevables en caractéristique 0, \emph{Proc. Nat. Acad.
Sc. U.S.A.}, vol. 47, 1961, pp. 108--109.

\item[{[8]}] ZAPPA (G.). -- Sull'esistenza, sopra le superficie algebriche, di
sistemi continui completi infiniti, la cui curva è a serie caratteristica
incompleta, \emph{Pont. Acad. Sc. Acta}, vol. 9, 1945, pp. 91--93.
\end{description}

\section*{ADITIVO}
\addcontentsline{toc}{section}{Aditivo}
\begin{center}
{\small [añadido durante la corrección de pruebas]}
\end{center}

Parece ahora que las conjeturas de [2], III, n.º 8 son falsas, incluso para
variedades no singulares sobre un cuerpo de característica cero, tanto en lo
que respecta a la existencia como a la cuasiproyectividad del cociente, e
incluso cuando $G$ actúa con grafo cerrado.

\def\FgaUnitKey{e221}\clearpage
\chapter*{Fe de erratas de la exposición n.º 221}
\addcontentsline{toc}{chapter}{Fe de erratas de la exposición n.º 221}
\setcounter{footnote}{3}
\src{302}

\noindent\textbf{6. TDTE IV: Los esquemas de Hilbert.}

\begin{center}
[Seminario Bourbaki, t. 13, 1960/61, n.º 221, 28 pp.]
\end{center}

\noindent\textbf{Página 6, Teorema 2.1, 3.ª línea.} --- En lugar de
«es necesario», léase «es necesario y suficiente».

\noindent\textbf{Página 15, línea 13.} --- En lugar de «$P\neq Q$», léase
«$P<Q$».

\noindent\textbf{Línea 15.} --- En lugar de «Si, por el contrario, $P=Q$»,
léase «En el caso contrario, se tendrá $P(n)\geq Q(n)$ para $n$ grande».

\noindent\textbf{Línea 18, fórmula (*).} --- En lugar de «$<$» y «$>$»,
léase «$\leq$» y «$\geq$».

\noindent\textbf{5.ª línea desde abajo.} --- En lugar de «$r$», léase «$Z$».

\noindent\textbf{Página 18, Observación 3.9.} --- Señalemos que Hartshorne ha
estudiado las componentes conexas de los esquemas de Hilbert sobre un cuerpo
algebraicamente cerrado. Este autor demuestra que los
$\operatorname{Hilb}^{P}_{r}$ son conexos y determina los pares $(r,P)$ para
los cuales $\operatorname{Hilb}^{P}_{r}\neq\varnothing$\footnote{HARTSHORNE
(R.). --- Connectedness of the Hilbert scheme (Tesis, Harvard, 1962).}.

\noindent\textbf{Página 23, 5.ª línea desde abajo.} --- En lugar de «una
componente», léase «de una componente».

\noindent\textbf{Páginas 23 y 24, Observación 5.5.} --- A propósito del ejemplo
de Zappa, señalemos que Mumford acaba de construir una componente irreducible
del esquema de Hilbert de $\mathbf P^3_{\mathbf C}$ (cuyos puntos generales
representan curvas no singulares de grado 14 y género 24) que es no reducida
en su punto genérico. Haciendo explotar una de las curvas obtenidas, obtiene
también un esquema proyectivo regular de dimensión 3 sobre $\mathbf C$, cuyo
esquema formal de módulos es no reducido en su punto genérico o, lo que
equivale a lo mismo, cuya variedad local de módulos, en el sentido de la
geometría analítica sobre $\mathbf C$ (cf. Seminario Cartan, t. 13, 1960/61),
es no reducida en todos sus puntos.

\def\FgaUnitKey{232}\unit{Exposición 232 -- Esquemas de Picard: teoremas de existencia}
\sid{232}
\src{143}
\seminares{14.\textsuperscript{o} año, 1961/62, n.\textsuperscript{o} 232\\Febrero de 1962}
\paperhead{TÉCNICA DE DESCENSO Y TEOREMAS DE EXISTENCIA EN GEOMETRÍA ALGEBRAICA\\
V. ESQUEMAS DE PICARD: TEOREMAS DE EXISTENCIA}%
  {por Alexander GROTHENDIECK}

\fsection{1}{Grupos y funtores de Picard relativos}

Para todo preesquema (más generalmente, todo espacio anillado) $X$, llamamos al
grupo de clases, consideradas salvo isomorfismo, de Módulos invertibles (es
decir, Módulos localmente isomorfos a $\mathcal O_X$) sobre $X$ el
\emph{grupo de Picard} (absoluto) de $X$, y lo denotamos por
$\operatorname{Pic}(X)$. Se tiene, por tanto, un isomorfismo canónico
\begin{equation}
\operatorname{Pic}(X)\xrightarrow{\sim}H^1(X,\mathcal O_X^*),
\tag{1.1}
\end{equation}
donde $\mathcal O_X^*$ denota el haz de unidades de $\mathcal O_X$ (que se
identifica, en efecto, con el haz de automorfismos del Módulo invertible
estándar $\mathcal O_X$). Es claro que
$X\longmapsto\operatorname{Pic}(X)$ es un funtor contravariante en $X$, y que
el isomorfismo (1.1) es funtorial.

Si $X$ es un preesquema sobre otro preesquema $S$, al variar $S'$ en la
categoría $(\mathrm{Sch})/S$ de preesquemas sobre $S$, la construcción anterior
proporciona un funtor contravariante
$S'\longmapsto\operatorname{Pic}(X\times_SS')$. Este funtor no tiene ninguna
posibilidad de ser «representable» ([4], II, A), pues los automorfismos no
triviales de los Módulos invertibles que se pretende clasificar impiden que sea
de «naturaleza local» ([5], IV, 5.4). Hay, pues, que «hacerlo local». De modo
general, para todo preesquema relativo $X/S$, introduzcamos el grupo relativo
\begin{equation}
\operatorname{Pic}'(X/S)=H^0\bigl(S,R^1f_*(\mathcal O_X^*)\bigr),
\tag{1.2}
\end{equation}
donde $f:X\longrightarrow S$ es el morfismo estructural (compárese [4], II,
C 3). En el pasaje citado se llama a este grupo grupo de Picard relativo; por
razones que aparecerán enseguida, conviene llamarlo aquí \emph{grupo de Picard
relativo restringido} de $X/S$. Al variar $S'$ en $(\mathrm{Sch})/S$, la
asignación $S'\longmapsto\operatorname{Pic}'(X\times_SS'/S')$ es un funtor
contravariante en $S'$, denotado asimismo por $\operatorname{Pic}'_{X/S}$ y
definido, por tanto, esencialmente por
\begin{equation}
\operatorname{Pic}'_{X/S}(S')=
\operatorname{Pic}'(X\times_SS'/S').
\tag{1.3}
\end{equation}

\src{144}
Este funtor es ahora de «naturaleza local», pues se ha efectuado el ajuste
necesario. Intuitivamente, el segundo miembro de (1.3) se interpreta como el
conjunto de «familias algebraicas» de clases de haces invertibles sobre (las
fibras de) $X/S$, indexadas por el preesquema de parámetros $S'/S$. Cuando el
funtor $\operatorname{Pic}'$ es representable, el preesquema sobre $S$ que lo
representa se denota por $\operatorname{Pic}'_{X/S}$ y se llama
\emph{preesquema de Picard} de $X$ sobre $S$; se tiene entonces
\begin{equation}
\operatorname{Hom}_S(S',\operatorname{Pic}'_{X/S})
\simeq\operatorname{Pic}'_{X/S}(S')
=\operatorname{Pic}'(X\times_SS'/S').
\tag{1.4}
\end{equation}

Hay, sin embargo, casos importantes en los que $\operatorname{Pic}'_{X/S}$ no
es representable (por ejemplo, una variedad de «Brauer--Severi» sobre un
cuerpo $k$ sin punto racional sobre $k$), aunque sigue habiendo una definición
natural de preesquema de Picard relativo. Ello se debe a que, al definir el
funtor $\operatorname{Pic}'$ a partir de los grupos de Picard absolutos
$\operatorname{Pic}(X\times_SS'/S')$, aún no se ha localizado suficientemente;
más precisamente, $\operatorname{Pic}'$ no es, en general, «compatible con el
descenso fielmente plano». Hagámoslo explícito.

Sea $\mathfrak M$ la clase de los morfismos de preesquemas que son
\emph{fielmente planos y cuasicompactos}; esta clase es estable por cambio de
base y por composición. Sea $P$ un funtor contravariante de $(\mathrm{Sch})/S$
a la categoría de conjuntos y, para todo $S$-morfismo
$u:T'\longrightarrow T$ con $u\in\mathfrak M$, consideremos el diagrama
\begin{equation}
P(T)\longrightarrow P(T')\rightrightarrows P(T'\times_TT')
\tag{1.5}
\end{equation}
obtenido al aplicar $P$ a
\[
T\longleftarrow T'\mathrel{\substack{\xleftarrow{\,\mathrm{pr}_1\,}\\[-0.35ex]
\xleftarrow[\,\mathrm{pr}_2\,]{}}}T'\times_TT'.
\]
Si $P$ es representable, la teoría del descenso ([4], I, B, teorema 2) implica
que (1.5) es exacto para todo $u\in\mathfrak M$. Esto se expresa diciendo que
$P$ es compatible con $\mathfrak M$; en este caso, que $P$ es «compatible con
el descenso fielmente plano», o también que el «prehaz» $P$ sobre
$(\mathrm{Sch})/S$ es un «haz» para la noción de localización proporcionada
por $\mathfrak M$. Para un $P$ arbitrario, un procedimiento estándar, familiar
por la localización topológica ordinaria, permite asociarle un «haz»
$\mathcal P$ y un homomorfismo de funtores $P\longrightarrow\mathcal P$
universal en el sentido evidente. El cálculo de $\mathcal P$ puede hacerse
explícito como sigue. Para definir $\mathcal P(T)$, para cada $T'$ sobre $T$
tal que $u:T'\longrightarrow T$ pertenezca a $\mathfrak M$, sea

\src{145}
$\overline H^0(T'/T,P)$ el subconjunto de $P(T')$ formado por los elementos
$\xi$ cuyas dos imágenes $\xi_1,\xi_2$ en $P(T'\times_TT')$ tienen la misma
imagen en $P(T'')$ para algún morfismo
$v:T''\longrightarrow T'\times_TT'$ con $v\in\mathfrak M$. [N.B. -- El
conjunto $\overline H^0$ así definido es, por tanto, mayor que el conjunto
$H^0(T'/T,P)$ introducido en [4], I, A, 4(a).] Al variar $T'$ sobre un $T$ fijo
(siempre con $u\in\mathfrak M$), los $\overline H^0(T'/T,P)$ forman un sistema
inductivo cuando se dota al conjunto de los $T'$ del preorden definido por la
dominación, y se pone
\begin{equation}
\mathcal P(T)=\varinjlim_{T'}\overline H^0(T'/T,P).
\tag{1.6}
\end{equation}
La ley funtorial en $T$ de esta expresión es inmediata.

Cuando
\[
P(T)=\operatorname{Pic}(X\times_ST),
\]
el funtor contravariante sobre $(\mathrm{Sch})/S$ definido por (1.6) se llama
\emph{funtor de Picard relativo} de $X$ sobre $S$ y se denota por
$\operatorname{Pic}_{X/S}$; el grupo $\operatorname{Pic}_{X/S}(S)$ se llama
\emph{grupo de Picard relativo} de $X$ sobre $S$ y se denota por
$\operatorname{Pic}(X/S)$. Se obtiene entonces la biyección evidente
\begin{equation}
\operatorname{Pic}_{X/S}(T)\simeq\operatorname{Pic}(X\times_ST/T).
\tag{1.7}
\end{equation}

Así, un elemento de $\operatorname{Pic}(X/S)$ viene definido por un elemento
$\xi'$ de un grupo $\operatorname{Pic}(X\times_SS')$, donde
$S'\longrightarrow S$ es fielmente plano y cuasicompacto, tal que existe un
morfismo fielmente plano y cuasicompacto
$S''\longrightarrow S'\times_SS'$ para el cual coinciden las dos imágenes
inversas de $\xi'$ en $\operatorname{Pic}(X\times_SS'')$. Un elemento $\xi'$ de
$\operatorname{Pic}(X\times_SS')$ y un elemento $\xi_1$ de
$\operatorname{Pic}(X\times_SS_1)$ que satisfagan las condiciones anteriores
definen el mismo elemento de $\operatorname{Pic}(X/S)$ si y solo si existe un
morfismo fielmente plano y cuasicompacto
$S'_1\longrightarrow S'\times_SS_1$ tal que las imágenes de ambos elementos en
$\operatorname{Pic}(X\times_SS'_1)$ sean iguales. A menudo sigue siendo
conveniente trabajar con el funtor $P'=\operatorname{Pic}'_{X/S}$ introducido
antes; se comprueba inmediatamente que el morfismo canónico
$P\longrightarrow P'$ induce un isomorfismo
\begin{equation}
\mathcal P\xrightarrow{\sim}\mathcal P',
\tag{1.8}
\end{equation}
lo que proporciona una descripción, con frecuencia más cómoda, de
$\operatorname{Pic}_{X/S}$ en términos de
$\operatorname{Pic}'_{X/S}=P'$. En virtud de 2.3, al sustituir $P$ por $P'$ en
la descripción anterior de $\operatorname{Pic}(X/S)$, se puede de hecho
\src{146}
tomar $S''=S'\times_SS'$ y $S'_1=S'\times_SS_1$, al menos bajo las condiciones
especificadas en el pasaje citado.

Cuando el funtor $\operatorname{Pic}_{X/S}$ es representable, se dice que
$X/S$ admite un preesquema de Picard, y el preesquema sobre $S$ que representa
el funtor se llama \emph{preesquema de Picard} de $X$ sobre $S$ y sigue
denotándose por $\operatorname{Pic}_{X/S}$. Basta evidentemente que
$P'=\operatorname{Pic}'_{X/S}$ sea representable, pues entonces $P'$ ya es un
«haz», y (1.8) muestra que $P'\longrightarrow\mathcal P'$ es el morfismo
canónico
\begin{equation}
\operatorname{Pic}'_{X/S}\longrightarrow\operatorname{Pic}_{X/S},
\tag{1.9}
\end{equation}
que es entonces un isomorfismo. Nuestra terminología concuerda así con la
introducida anteriormente en (1.4). En general, cuando
$\operatorname{Pic}_{X/S}$ existe, queda definido por el isomorfismo funtorial
\begin{equation}
\operatorname{Hom}_S(S',\operatorname{Pic}_{X/S})\xrightarrow{\sim}
\operatorname{Pic}(X\times_SS'/S').
\tag{1.10}
\end{equation}

\fsection{2}{Relaciones entre los diversos grupos de Picard relativos y absolutos}

\result{PROPOSICIÓN}{2.1} Sea $f:X\longrightarrow S$ un morfismo tal que
$\mathcal O_S\xrightarrow{\sim}f_*(\mathcal O_X)$. Se tiene entonces una
sucesión exacta
\[
0\longrightarrow\operatorname{Pic}(S)\longrightarrow\operatorname{Pic}(X)
\longrightarrow\operatorname{Pic}'(X/S).
\]
Cuando $X$ admite una sección sobre $S$, el último morfismo es sobreyectivo;
es decir, se tiene un isomorfismo
\[
\operatorname{Pic}'(X/S)\xrightarrow{\sim}
\operatorname{Pic}(X)/\operatorname{Pic}(S).
\]

La sucesión exacta puede considerarse como la sucesión exacta en grados bajos
asociada a la sucesión espectral de Leray para $f$ y $\mathcal O_X^*$. La
segunda afirmación es igualmente formal.

\result{PROPOSICIÓN}{2.2} Sea $f:X\longrightarrow S$ un morfismo cuasicompacto
y separado tal que $\mathcal O_S\xrightarrow{\sim}f_*(\mathcal O_X)$, y sea
$S'\longrightarrow S$ fielmente plano y cuasicompacto. Entonces:

\noindent (i) $\operatorname{Pic}'(X/S)\longrightarrow
\operatorname{Pic}'(X\times_SS'/S')$ es inyectivo;

\noindent (ii) si $X$ admite localmente una sección sobre $S$ (es decir, todo
$s\in S$ posee un entorno abierto $U$ tal que $X|U$ admite una sección sobre
$U$), entonces el diagrama
\[
\operatorname{Pic}'(X/S)\longrightarrow
\operatorname{Pic}'(X\times_SS'/S')\rightrightarrows
\operatorname{Pic}'(X\times_SS''/S'')
\]

\src{147}
donde $S''=S'\times_SS'$, es exacto.

La primera afirmación se deduce de las propiedades elementales del descenso
fielmente plano y de la observación general siguiente. Si
$f:X\longrightarrow S$ satisface
$\mathcal O_S\xrightarrow{\sim}f_*(\mathcal O_X)$, entonces el funtor
$\mathcal F\longmapsto f^*(\mathcal F)$ de los Módulos localmente libres de
tipo finito sobre $S$ a los Módulos localmente libres de tipo finito sobre $X$
es plenamente fiel, y su imagen esencial está formada por los Módulos
$\mathcal G$ sobre $X$ tales que $f_*(\mathcal G)$ es localmente libre y el
homomorfismo canónico
\[
f^*\bigl(f_*(\mathcal G)\bigr)\longrightarrow\mathcal G
\]
es un isomorfismo. La segunda afirmación fue demostrada mediante la teoría del
descenso en [4], I, B, 4.

Los resultados de 2.2 también pueden enunciarse como sigue.

\result{COROLARIO}{2.3} Bajo las condiciones de 2.2, el homomorfismo canónico
(1.9) $\operatorname{Pic}'_{X/S}\longrightarrow\operatorname{Pic}_{X/S}$ es
inyectivo, y es biyectivo si $X$ admite localmente una sección sobre $S$.
(Así, en este último caso, el grupo de Picard relativo
$\operatorname{Pic}(X/S)$ y el grupo de Picard relativo restringido
$\operatorname{Pic}'(X/S)$ son idénticos.)

Combinando esto con 2.1 se obtiene:

\result{COROLARIO}{2.4} Bajo las condiciones de 2.2, se tiene una sucesión
exacta
\[
0\longrightarrow\operatorname{Pic}(S)\longrightarrow\operatorname{Pic}(X)
\longrightarrow\operatorname{Pic}(X/S).
\]
Cuando $X$ admite una sección sobre $S$, el último homomorfismo es
sobreyectivo; es decir, se tiene entonces un isomorfismo
\[
\operatorname{Pic}(X/S)\xrightarrow{\sim}
\operatorname{Pic}(X)/\operatorname{Pic}(S).
\]

\result{OBSERVACIÓN}{2.5} Sea $f:X\longrightarrow S$ un morfismo tal que
$\mathcal O_S\xrightarrow{\sim}f_*(\mathcal O_X)$, y sea $g$ una sección de
$X$ sobre $S$. Para un Módulo invertible $\mathcal L$ sobre $X$, llamemos al
isomorfismo $\mathcal O_S\xrightarrow{\sim}g^*(\mathcal L)$ una
$g$-rigidificación de $\mathcal L$, y llamemos a un Módulo invertible
$\mathcal L$ sobre $X$ provisto de una $g$-rigidificación \emph{Módulo
invertible $g$-rigidificado}. Todo automorfismo de una estructura de este
tipo es trivial, y $\operatorname{Pic}'(X/S)$ se identifica con el conjunto de
clases de isomorfismo de Módulos invertibles $g$-rigidificados sobre
$X$\footnote{El texto impreso dice «sobre $S$»; la rigidificación mediante una
sección $g:S\to X$ concierne aquí al Módulo $\mathcal L$ sobre $X$.}. (Este es
el hecho que permitió emplear la teoría del descenso en la demostración de
2.2(ii).) Esto da una nueva interpretación de $\operatorname{Pic}(X/S)$, al
menos cuando $f$ es además cuasicompacto y separado, de modo que
$\operatorname{Pic}(X/S)\simeq\operatorname{Pic}'(X/S)$ por 2.3.

\src{148}
\result{OBSERVACIÓN}{2.6} Sea $f:X\longrightarrow S$ como en 2.2, y sea
$S'\longrightarrow S$ fielmente plano y cuasicompacto, de modo que exista un
$S$-morfismo $S'\longrightarrow X$, o, equivalentemente, que
$X'=X\times_SS'$ admita una sección sobre $S'$. Pongamos
$S''=S'\times_SS'$ y $X''=X\times_SS''$, y consideremos la sucesión exacta
\[
\operatorname{Pic}(X/S)\longrightarrow\operatorname{Pic}(X'/S')
\rightrightarrows\operatorname{Pic}(X''/S'').
\]
Aplicando 2.4 se obtiene una sucesión exacta
\[
\operatorname{Pic}(X/S)\longrightarrow
\operatorname{Pic}(X')/\operatorname{Pic}(S')
\rightrightarrows
\operatorname{Pic}(X'')/\operatorname{Pic}(S''),
\]
y, en particular, todo elemento del grupo de Picard relativo ya «proviene»
de un elemento de $\operatorname{Pic}(X')$. Esto simplifica sustancialmente la
descripción del grupo de Picard relativo dada en la sección anterior, y también
la del funtor de Picard de $X$ sobre $S$, pues para todo $T$ sobre $S$ el
argumento precedente se aplica a $X\times_ST/T$ y al morfismo
$T'=S'\times_ST\longrightarrow T$. Si, por ejemplo, $f$ mismo es fielmente
plano, se puede tomar $S'=X$. Cuando $f$ es además de tipo finito
(respectivamente liso, etc.), esto permite limitarse, en la descripción del
funtor de Picard relativo
$T\longmapsto\operatorname{Pic}(X\times_ST/T)$, a cambios de base
$T'\longrightarrow T$ que sean de tipo finito (respectivamente lisos, etc.).
Cuando $f$ es proyectivo y plano y $S$ es localmente noetheriano, se demuestra
que el morfismo $S'\longrightarrow S$ anterior puede elegirse de modo que $S'$
sea la unión disjunta de recubrimientos planos $S'_i$ de abiertos $S_i$ que
cubren $S$; si $f$ es además separable, los $S'_i$ pueden tomarse étale sobre
$S_i$.

\fsection{3}{El teorema principal de existencia: enunciado}

No se dispone, ni siquiera a título de conjetura, de un enunciado de existencia
de preesquemas de Picard que englobe todos los casos conocidos. Una condición
«prácticamente necesaria», si puede decirse así, es que
$f:X\longrightarrow S$ sea propio, asegurando las propiedades esenciales de
finitud, y plano. Estas condiciones no son suficientes, incluso si $S$ es el
espectro del álgebra de los números duales $k[t]/(t^2)$ sobre un cuerpo $k$
(por ejemplo, el cuerpo $\mathbf C$ de los números complejos) y $X$ tiene
dimensión 1. En el momento de redactar esta exposición, los teoremas de
existencia más importantes para el preesquema de Picard se deducen del teorema
siguiente.

\result{TEOREMA}{3.1} Sea $f:X\longrightarrow S$ un morfismo de preesquemas
localmente noetherianos. Supongamos:

\noindent (i) $f$ es proyectivo;

\src{149}
\noindent (ii) $f$ es plano;

\noindent (iii) las fibras geométricas de $f$ son íntegras.

Bajo estas condiciones, $\operatorname{Pic}_{X/S}$ existe.

La demostración, que se esbozará en las dos secciones siguientes, mostrará al
mismo tiempo lo que sigue. Sea $\xi$ la sección de
$\operatorname{Pic}_{X/S}$ correspondiente a un haz muy amplio
$\mathcal O_X(1)$ sobre $X/S$, es decir, inducido por una inmersión proyectiva
$X\longrightarrow\mathbf P(\mathcal E)$. Existe entonces un abierto $U$ de
$\operatorname{Pic}_{X/S}$, unión disjunta de abiertos cuasiproyectivos sobre
$S$, tal que $U$ es estable por traslación por $\xi$, y
$\operatorname{Pic}_{X/S}$ es la unión creciente de los abiertos $U-n\xi$,
todos isomorfos a $U$. En particular, bajo las condiciones de 3.1,
$\operatorname{Pic}_{X/S}$ es separado sobre $S$.

\result{OBSERVACIÓN}{3.2} Algunos ejemplos muestran (por ejemplo, con $S$ el
espectro de un anillo de valoración discreta y $X$ de dimensión relativa 1
sobre $S$) que, si se omite la hipótesis (iii) de 3.1 y se sustituye por la
hipótesis más débil de que, para todo $s\in S$, el homomorfismo
\[
k(s)\longrightarrow H^0(X_s,\mathcal O_{X_s})
\]
es un isomorfismo, entonces $\operatorname{Pic}_{X/S}$ no es necesariamente
separado sobre $S$. Esto ocurre tanto cuando las fibras geométricas de $f$ son
reducidas pero una fibra geométrica genérica íntegra «estalla» por
especialización en dos componentes irreducibles, como cuando las fibras
geométricas de $f$ son irreducibles pero una fibra geométrica genérica íntegra
se especializa en una «fibra múltiple». El primer caso se presenta, por ejemplo,
con una cónica que degenera en dos rectas secantes; D. Mumford proporcionó un
ejemplo del segundo, con una curva elíptica que degenera en una curva elíptica
doble. Estos ejemplos son válidos en toda característica.

\result{OBSERVACIÓN}{3.3} Bajo las condiciones de 3.1, no sé si
$\operatorname{Pic}_{X/S}$ es una unión disjunta de abiertos de tipo finito, y
por tanto cuasiproyectivos, sobre $S$.\footnote{La fe de erratas de 1962 informa de
que D. Mumford respondió afirmativamente a esta cuestión mediante la teoría de
los esquemas de Picard desarrollada aquí.} Los polinomios de Hilbert
$Q\in\mathbf Q[t]$, como en el caso de los esquemas de Hilbert ([4], IV), dan
una descomposición de $\operatorname{Pic}_{X/S}$ en una unión disjunta de
abiertos $\operatorname{Pic}^{Q}_{X/S}$, y parece plausible que estos sean
cuasiproyectivos sobre $S$; esto se verá al menos en la exposición siguiente
cuando $f$ sea liso. Obsérvese que, si se sustituye la hipótesis (i) por la de
que $X$ sea localmente proyectivo sobre $S$ (lo cual basta para 3.1, pues la
existencia de $\operatorname{Pic}_{X/S}$ es manifiestamente local sobre $S$),
resulta, por el contrario, fácil

\src{150}
dar ejemplos en los que $\operatorname{Pic}_{X/S}$ contiene componentes conexas
que no son de tipo finito sobre $S$. Sea, por ejemplo, $X_0$ una variedad
algebraica no singular y proyectiva sobre un cuerpo algebraicamente cerrado
$k$, provista de un automorfismo $u$ y de un elemento $\xi$ del grupo de
Néron--Severi de $X_0$ tales que los $u^n(\xi)$ sean dos a dos distintos. Se
puede tomar para $X_0$ el producto de una curva elíptica $E$ por sí misma, y
para $u$ el automorfismo
\[
(x,y)\longmapsto(x,y+x)
\]
de $E\times E$. Sea $S$ la unión de dos curvas irreducibles no singulares que
se cortan en dos puntos $a$ y $b$. Existe un recubrimiento principal conexo de
grupo $\mathbf Z$, $P$, sobre $S$; usando las acciones de $\mathbf Z$ sobre $X_0$
definidas por $u$, se obtiene un fibrado asociado sobre $S$ de fibra $X_0$,
trivial sobre $S-a$ y $S-b$ y, en el caso particular considerado, de hecho un
esquema abeliano sobre $S$. Se comprueba fácilmente que
$\operatorname{Pic}_{X/S}$, que también es el fibrado asociado a $P$ y a las
acciones de $\mathbf Z$ sobre $\operatorname{Pic}_{X_0/k}$ por medio de $u$,
contiene una componente conexa isomorfa a
\[
P\times\operatorname{Pic}^{0}_{X_0/k},
\]
donde $\operatorname{Pic}^{0}$ denota la componente conexa de la identidad en
$\operatorname{Pic}$; esta componente no es de tipo finito sobre $S$.
[Fenómenos análogos se producen también en diversos casos de preesquemas de
Picard no separados sobre $S$, como los considerados en 3.2.]

\fsection{4}{Divisores de Cartier relativos y fibrados proyectivos}

Solo necesitaremos divisores positivos y omitiremos este calificativo adicional
en toda la sección.

Sea $X$ un preesquema. Un \emph{divisor de Cartier} sobre $X$, o simplemente un
divisor sobre $X$, es un subpreesquema cerrado $D$ de $X$ definido por un Ideal
$\mathcal J$ que es un Módulo invertible, es decir, localmente generado por una
sección que no es divisor de cero en $\mathcal O_X$. A $D$ asociamos el Módulo
invertible
\[
\mathcal L(D)=\mathcal J^{-1}.
\]
La inyección canónica $\mathcal J\longrightarrow\mathcal O_X$ proporciona un
homomorfismo canónico
\[
s_D:\mathcal O_X\longrightarrow\mathcal J^{-1}=\mathcal L(D),
\qquad\text{es decir,}\quad s_D\in\Gamma(X,\mathcal L(D)).
\]
Además, dar un divisor equivale esencialmente a dar un Módulo invertible
$\mathcal L$ sobre $X$ junto con una sección $s$ que no sea divisor de cero en
ningún punto, asociando recíprocamente a tal par $(\mathcal L,s)$ el
«divisor» de $s$, denotado por $\operatorname{div}(s)$. Para un $\mathcal L$
invertible fijo sobre $X$, el conjunto de

\src{151}
divisores $D$ que definen $\mathcal L$ está en biyección con el conjunto
cociente
\[
\Gamma(X,\mathcal L)^*/\Gamma(X,\mathcal O_X^*),
\]
donde $\Gamma(X,\mathcal L)^*$ es el subconjunto de
$\Gamma(X,\mathcal L)$ formado por las secciones que no son divisores de cero
en ningún punto.

Supongamos ahora dado un morfismo localmente de tipo finito
$f:X\longrightarrow S$, con $S$ localmente noetheriano para simplificar. Sean
$\mathcal J$ un Ideal coherente sobre $X$, $D$ el subesquema de $X$ que define,
$x\in X$, y pongamos $s=f(x)$. Las condiciones siguientes son equivalentes:

\noindent (i) $\mathcal J$ es invertible en $x$ (es decir, $\mathcal J_x$ está
generado por un elemento regular de $\mathcal O_{X,x}$), y $D$ es plano sobre
$S$ en $x$;

\noindent (ii) $X$ y $D$ son planos sobre $S$ en $x$, y $D_s$ es un divisor de
Cartier sobre la fibra $X_s$ en $x$;

\noindent (iii) $X$ es plano sobre $S$ en $x$, y $\mathcal J_x$ está generado
por un elemento $f_x$ que induce sobre $X_s$ un germen que no es divisor de
cero.

Se dice entonces que $D$ es un \emph{divisor de Cartier relativo} sobre $X/S$
(o simplemente un divisor relativo sobre $X/S$) en el punto considerado. La
condición (i) muestra que $D$ es también un divisor relativo en los puntos
vecinos de $x$. Por tanto, si $X$ y $D$ son planos sobre $S$ y $D$ es propio
sobre $S$, el conjunto de los $s\in S$ tales que $D_s$ es un divisor de Cartier
en $X_s$ (es decir, tales que $D$ es un divisor de Cartier relativo en los
puntos de $X_s$) es abierto en $S$. Además, la definición anterior se formuló
precisamente para que los divisores de Cartier relativos sean estables por todo
cambio de base $S'\longrightarrow S$. Sea $\operatorname{Div}(X/S)$ el
conjunto de divisores relativos sobre $X/S$, y definamos el funtor
contravariante en $S'$ sobre $S$ mediante
\[
\operatorname{Div}_{X/S}(S')=
\operatorname{Div}(X\times_SS'/S').
\]

Supongamos que $X$ es plano y propio sobre $S$. Por la caracterización (ii) de
los divisores de Cartier relativos, $\operatorname{Div}_{X/S}$ puede
considerarse como un subfuntor del funtor $\operatorname{Hilb}_{X/S}$ definido
en [4], IV, y la inclusión
\[
\operatorname{Div}_{X/S}\longrightarrow\operatorname{Hilb}_{X/S}
\]
es «representable por inmersiones abiertas» (cf. [5], IV, 3.13), por las
observaciones precedentes. El teorema principal de existencia de [4], IV, da
entonces:

\src{152}
\result{PROPOSICIÓN}{4.1} Supongamos que $f:X\longrightarrow S$ es proyectivo
y plano. Entonces el funtor $\operatorname{Div}_{X/S}$ es representable; más
precisamente, está representado por un abierto de
$\operatorname{Hilb}_{X/S}$.

Utilizando, para un haz muy amplio fijo $\mathcal O_X(1)$ sobre $X/S$, la
descomposición canónica de $\operatorname{Hilb}_{X/S}$ en los abiertos
$\operatorname{Hilb}^{Q}_{X/S}$ correspondientes a los polinomios de Hilbert
$Q\in\mathbf Q[t]$, se obtiene asimismo la descomposición análoga
\[
\operatorname{Div}_{X/S}=
\coprod_{Q\in\mathbf Q[t]}\operatorname{Div}^{Q}_{X/S}
\]
como unión disjunta de abiertos cuasiproyectivos sobre $S$.

La asignación $D\longmapsto\mathcal L(D)$ proporciona además un homomorfismo
funtorial
\begin{equation}
\operatorname{Div}_{X/S}\longrightarrow\operatorname{Pic}_{X/S}
\tag{+}
\end{equation}
que estudiaremos ahora; bajo condiciones bastante generales resultará
relativamente representable ([5], IV, 3). Partamos de un elemento $\xi$ de
$\operatorname{Pic}_{X/S}(S')$, tomando $S'=S$ para simplificar la notación, y
mostremos que el subfuntor correspondiente de $\operatorname{Div}_{X/S}$ es
representable. Consideremos primero el caso en que $\xi$ viene definido por un
Módulo invertible $\mathcal L$ sobre $X$. Supongamos que $X$ es propio y plano
sobre $S$ y que sus fibras geométricas sobre $S$ son íntegras. Esto implica
también ([2], III, párr. 7) que
$\mathcal O_S\xrightarrow{\sim}f_*(\mathcal O_X)$ y que esta relación sigue
siendo válida tras todo cambio de base $S'\longrightarrow S$. Los divisores de
Cartier relativos $D$ sobre $X/S$ para los que $\mathcal L(D)$ y $\mathcal L$
definen el mismo elemento de
$\operatorname{Pic}(X/S)=\operatorname{Pic}_{X/S}(S)$, o, equivalentemente por
2.4, para los que $\mathcal L(D)$ y $\mathcal L$ son localmente isomorfos sobre
$S$, están en biyección con las secciones del haz cociente
$f_*(\mathcal L)^*/\mathcal O_S^*$. Esta correspondencia es compatible con el
cambio de base. Consideraciones generales de tipo «Künneth» del pasaje citado
(véase también [6]) muestran que la propiedad de ser propio de $X/S$ y la
planitud de $\mathcal L$ sobre $S$ implican la existencia de un Módulo coherente
$\mathcal Q$ sobre $S$, definido salvo isomorfismo único, y de un isomorfismo
de haces
\[
f_*(\mathcal L)\xrightarrow{\sim}
\operatorname{Hom}_{\mathcal O_S}(\mathcal Q,\mathcal O_S).
\]
\footnote{El texto impreso coloca $\mathcal O_X$ como subíndice de
$\operatorname{Hom}$; puesto que $\mathcal Q$ es un Módulo sobre $S$, su dual
se toma sobre $\mathcal O_S$.} La formación de $\mathcal Q$ es compatible con
el cambio de base. Aquí $f_*(\mathcal L)^*$ denota el subhaz de conjuntos de
$f_*(\mathcal L)$ cuyas secciones sobre $U$ son las secciones de $\mathcal L$
sobre $f^{-1}(U)$ que definen divisores de Cartier relativos sobre

\src{153}
$f^{-1}(U)/U$, es decir, que inducen secciones que no son divisores de cero
sobre los $X_s$ para $s\in U$. Como las fibras $X_s$ son íntegras, esto
significa simplemente que las secciones inducidas sobre las fibras $X_s$ no
son idénticamente nulas o, equivalentemente, en términos de homomorfismos
locales $\mathcal Q\longrightarrow\mathcal O_S$, que estos homomorfismos son
sobreyectivos, por el lema de Nakayama. Así, el conjunto de secciones de
$f_*(\mathcal L)^*/\mathcal O_S^*$ está en biyección con el conjunto de
Módulos cociente invertibles de $\mathcal Q$, o, por la definición del fibrado
proyectivo $\mathbf P(\mathcal Q)$ asociado al Módulo coherente $\mathcal Q$
(cf. [5], V, 2), con el conjunto de secciones de $\mathbf P(\mathcal Q)$ sobre
$S$. Esta descripción es compatible con las imágenes inversas y proporciona,
por tanto, el teorema siguiente.

\result{TEOREMA}{4.3} Sea $f:X\longrightarrow S$ un morfismo propio y plano de
fibras geométricas íntegras, donde $S$ es localmente noetheriano, y sea
$\mathcal L$ un Módulo invertible sobre $X$. Para todo $S'$ sobre $S$, sea
$T(S')$ el conjunto de divisores relativos $D$ sobre $X\times_SS'/S'$ tales
que $\mathcal L(D)$ sea isomorfo a
$\mathcal L\otimes_{\mathcal O_S}\mathcal O_{S'}$ localmente sobre $S'$, es decir, tales que
$\mathcal L(D)$ y $\mathcal L\otimes_{\mathcal O_S}\mathcal O_{S'}$ definan
el mismo elemento de $\operatorname{Pic}(X\times_SS'/S')$. Existe entonces un
Módulo coherente $\mathcal Q$ sobre $S$, determinado salvo isomorfismo único,
tal que el funtor $T$ está representado por el fibrado proyectivo
$\mathbf P(\mathcal Q)$.

\result{COROLARIO}{4.4} Si, además, $f$ es proyectivo, el homomorfismo funtorial
$\operatorname{Div}_{X/S}\longrightarrow\operatorname{Pic}_{X/S}$ es
representable por morfismos proyectivos.

En efecto, cuando $X$ admite una sección (respectivamente, localmente una
sección) sobre $S$, el homomorfismo anterior está representado por fibrados
proyectivos (respectivamente, fibrados proyectivos locales), en virtud de 4.3 y
2.1. Cuando $f$ es cuasiproyectivo, se reduce fácilmente al caso anterior
mediante un argumento de descenso que utiliza cuasisecciones finitas y planas
de $X$ localmente sobre $S$.

\result{OBSERVACIÓN}{4.5} Bajo las condiciones de 4.3, el Módulo $\mathcal Q$
sobre $S$ no es, en general, localmente libre, como expresa el hecho de que la
dimensión de las fibras reducidas de $\mathcal Q$, es decir, la dimensión de
$H^0(X_s,\mathcal L_s)$ para $s\in S$ variable, pueda dar saltos. Dado un
Módulo coherente $\mathcal Q$ sobre un preesquema localmente noetheriano $S$,
se comprueba fácilmente que, para un $s\in S$ fijo, $\mathcal Q$ es libre en
$s$ si y solo si $\mathbf P(\mathcal Q)$ es plano sobre $S$ en los puntos
situados sobre $s$; en tal caso es incluso liso sobre $S$ en esos puntos. Aquí,
con $\mathcal Q$ definido en términos de $\mathcal L$ como antes, esto significa
también que la formación de la imagen directa $f_*(\mathcal L)$ «conmuta con el
cambio de base
\src{154}
en un entorno de $s$», o, equivalentemente, que
$f_*(\mathcal L)_s\longrightarrow H^0(X_s,\mathcal L_s)$ es sobreyectivo. Esto
ocurre, por ejemplo, si $H^1(X_s,\mathcal L_s)=0$. Suponiendo que existan los
preesquemas en cuestión, estos criterios se aplican en particular a la situación
universal $\operatorname{Div}_{X/S}\longrightarrow\operatorname{Pic}_{X/S}$ y
dan, respectivamente, una condición necesaria y suficiente y una condición
suficiente para que este morfismo sea liso en un punto dado de
$\operatorname{Div}_{X/S}$.

\fsection{5}{Demostración del teorema principal de existencia}

Bajo las condiciones de 3.1, elijamos un Módulo muy amplio $\mathcal O_X(1)$
sobre $X/S$, y sea $\xi$ el elemento correspondiente de
$\operatorname{Pic}(X/S)$. Pongamos, para abreviar,
\[
P(S')=\operatorname{Pic}(X\times_SS'/S'),
\]
y supongamos primero, para simplificar, que $X/S$ admite una sección. Sea
$P^+(S')$ el subconjunto de $P(S')$ formado por las clases de Módulos
invertibles $\mathcal L$ sobre $X\times_SS'$ tales que
\[
R^if'_*(\mathcal L(n))=0\quad\text{para $i>0$ y todo $n\geq0$},
\qquad
f'_*(\mathcal L(n))\neq0\quad\text{para todo $n\geq0$}.
\]
Estas condiciones son estables por cambio de base y definen, por tanto, un
subfuntor $P^+$ de $P$, manifiestamente estable por traslación por $\xi$.
Utilizando los «teoremas A y B» de Serre ([2], III, 2) y las generalidades de
[5], IV, 5, se ve fácilmente que $P$ es representable si y solo si lo es $P^+$;
en ese caso, $P^+$ está representado por un abierto $U$ del preesquema
$\operatorname{Pic}_{X/S}$ que representa a $P$, y este es la unión creciente
de los abiertos $U-n\xi$.

Pongamos, para abreviar, $D=\operatorname{Div}_{X/S}$, y sea $D^+$ la imagen
inversa de $P^+$ por el morfismo canónico $D\longrightarrow P$. Se tiene así un
morfismo
\begin{equation}
D^+\longrightarrow P^+,
\tag{+}
\end{equation}
y ya se sabe que $D^+$ está representado por un abierto $D^+$ del preesquema
$D=\operatorname{Div}_{X/S}$, cuya existencia garantiza 4.1. Se deduce
fácilmente de 4.3 que (+) está representado por morfismos lisos, proyectivos y
sobreyectivos; más precisamente, por fibrados proyectivos asociados a Módulos
localmente libres y no nulos en ningún punto. En efecto, si $\mathcal L$ sobre
$X\times_SS'$ es como al comienzo de esta sección, $f'_*(\mathcal L)$ es un
Módulo localmente libre no nulo cuya formación conmuta con el cambio de base;
en la notación de 4.3, $\mathcal Q$ es entonces isomorfo al dual de
$f'_*(\mathcal L)$. Deduciremos la representabilidad de $P^+$ del criterio
general [5], IV, 4.7. En el pasaje citado, tomemos para $\mathfrak C$ la clase
de los morfismos de preesquemas fielmente

\src{155}
planos y cuasicompactos, que son epimorfismos efectivos por [4], I B. Acabamos
de verificar la condición (a) del pasaje citado, a saber, que (+) es
representable por morfismos pertenecientes a $\mathfrak C$. También se ha
verificado la condición (b), que dice que el funtor $P^+$ es compatible con el
descenso fielmente plano y cuasicompacto; esto es inmediato. Queda por probar
la condición (c) de dicho pasaje: la relación de equivalencia $R$ en el
preesquema $D^+$ deducida del morfismo (+), representable por $\mathfrak C$, es
$\mathfrak C$-efectiva, es decir, es efectiva y
$D^+\longrightarrow D^+/R$ pertenece a $\mathfrak C$. Obsérvese primero que
los abiertos $D^{+Q}$ de $D^+$ correspondientes a los distintos polinomios de
Hilbert virtuales $Q\in\mathbf Q[t]$ son estables por $R$, pues las fibras de
$R$ son conexas. Basta, por tanto, mostrar, para cada $Q$, que la relación de
equivalencia $R^Q$ inducida en $D^{+Q}$ es $\mathfrak C$-efectiva. Ahora bien,
$D^{+Q}$ es cuasiproyectivo, mientras que la relación de equivalencia $R^Q$ es
proyectiva y plana. Se cumplen, pues, las hipótesis de [4], III, 6.1, que dan el
resultado requerido.

En el caso general en que $X/S$ no admite necesariamente una sección, se reduce
fácilmente al caso anterior por descenso, o se repite la demostración precedente
con la modificación necesaria en la definición de $P^+$.

\result{OBSERVACIONES}{5.1} El método empleado es esencialmente el método de
Matsusaka para la construcción proyectiva de las variedades de Picard. El
resultado de [4], III, invocado para pasar al cociente efectivo, también puede
deducirse fácilmente del teorema de existencia de los esquemas de Hilbert
(véase, por ejemplo, [6]); clásicamente, estos cocientes se construían con
coordenadas de Chow. Obsérvese que la formación del abierto
$\operatorname{Pic}^{+}_{X/S}$ de $\operatorname{Pic}_{X/S}$ y su
descomposición en abiertos $\operatorname{Pic}^{+Q}_{X/S}$ cuasiproyectivos
sobre $S$, según los polinomios de Hilbert de los divisores que definen los
Módulos invertibles considerados, es compatible con el cambio de base, lo que
permite aplicar el argumento de descenso.

\result{OBSERVACIONES}{5.2} Es posible que $\operatorname{Pic}_{X/S}$ exista
siempre que $f:X\longrightarrow S$ sea propio y plano y los homomorfismos
\[
k(s)\longrightarrow H^0(X_s,\mathcal O_{X_s})\qquad(s\in S)
\]
sean isomorfismos; esta última condición significa entonces también que
$\mathcal O_S\xrightarrow{\sim}f_*(\mathcal O_X)$ y que esta relación sigue
siendo válida tras todo cambio de base $S'\longrightarrow S$. Esto es, al
menos, lo que se demuestra para espacios analíticos cuando $f$ es además
proyectivo, mediante un método distinto (el de Chow, salvo error), expuesto en
[5], IX, 3.1.\footnote{La fe de erratas de 1962 indica que la conjetura de existencia
formulada al comienzo de esta observación es falsa, aunque los métodos de
Mumford demuestran un teorema ligeramente más débil.} En ese método, el paso al

\src{156}
cociente por una relación de equivalencia propia y plana en un abierto del
esquema de divisores se sustituye por el paso al cociente por el grupo
proyectivo en el esquema de inmersiones de $X$ en $\mathbf P^r_S$. Este método
puede probablemente adaptarse a los esquemas mediante los resultados de
Mumford sobre el paso al cociente por el grupo proyectivo [8]. Por el momento,
solo existe una demostración escrita cuando $X$ tiene «muchas secciones
locales» sobre $S$, por ejemplo, cuando $X$ es separable sobre un anillo local
completo con cuerpo residual algebraicamente
cerrado.\footnote{El texto impreso omite «con cuerpo residual algebraicamente
cerrado»; la fe de erratas dice que el contraejemplo de Mumford hace indispensable
esta restricción.} En principio, este método tendría mayor alcance, pues también
da la existencia de preesquemas de Picard en casos en los que no son separados
y en los que, por tanto, el primer método necesariamente falla. (Técnicamente,
la dificultad consiste en que, cuando las fibras geométricas de $f$ dejan de
ser íntegras, el funtor considerado en 4.3 no está representado por el propio
fibrado proyectivo $\mathbf P(\mathcal Q)$, sino por un abierto suyo, lo que
conduce al delicado problema de tomar el cociente por una relación de
equivalencia plana pero no propia.)

\result{OBSERVACIÓN}{5.3} Obsérvese que la demostración dada aquí no utiliza ni
una construcción previa de las jacobianas de curvas o familias de curvas, ni la
teoría de variedades abelianas o esquemas abelianos. Difiere, por ello,
esencialmente de exposiciones tradicionales como el libro de Lang [7] o la
exposición de Chevalley [1], que siguen el camino esbozado por A. Weil. Incluso
para las jacobianas de curvas no singulares sobre un cuerpo algebraicamente
cerrado, digamos el cuerpo de los números complejos, la construcción dada aquí
es la única conocida que proporciona las propiedades muy fuertes adoptadas
como definición en la sección 1: esencialmente las de Chevalley, pero permitiendo
«variedades de parámetros» con elementos nilpotentes. Que la construcción de
los esquemas de Picard debe preceder, y no seguir, a la teoría de variedades
abelianas es claro a priori, pues los esquemas de Picard no son, en general,
esquemas abelianos ni se reducen a ellos. Esto ya se ve para curvas singulares
sobre un cuerpo algebraicamente cerrado, donde aparecen las «jacobianas
generalizadas» de Rosenlicht, que no son variedades abelianas. Además, la
teoría de variedades abelianas, y más generalmente de esquemas abelianos, se
simplifica mucho cuando se dispone de una teoría general de esquemas de Picard.
En particular, la teoría de dualidad para esquemas abelianos, especialmente los
resultados de tipo Cartier, se vuelve entonces casi formal (véase, por ejemplo,
[10]).

\src{157}
\result{OBSERVACIONES}{5.4} El «principio de compatibilidad» de Igusa para la
jacobiana de una curva que degenera en una curva singular solo puede
comprenderse adecuadamente como un teorema de existencia del esquema de Picard
de un esquema relativo en curvas $X/S$, no necesariamente liso sobre $S$. Es,
por tanto, un caso particular del teorema principal de existencia 3.1 cuando la
curva especializada es íntegra, es decir, en lenguaje clásico, irreducible de
multiplicidad 1. En la actualidad, el caso de una curva especial reducible,
incluso cuando sus componentes tienen multiplicidad 1, es decir, cuando la
curva especial es separable sobre el cuerpo residual, queda fuera de los
teoremas de existencia conocidos, salvo sobre un anillo de valoración discreta
completo con cuerpo residual algebraicamente cerrado; compárese 5.2. Esta
cuestión de existencia aparecerá ciertamente en una construcción geométrica, en
teoría de esquemas, de las «compactificaciones» de Baily--Satake de los esquemas
de módulos de curvas de género $g$. (Por el momento, esta compactificación solo
se conoce para $g=1$, gracias a los trabajos de Igusa.)

\fsection{6}{Teoremas de existencia relativos}

Esbozaremos aquí algunos casos útiles en los que la existencia de ciertos
esquemas de Picard implica la de otros, lo que permite deducir diversos
teoremas de existencia adicionales del teorema principal 3.1.

\result{PROPOSICIÓN}{6.1} Sea $f:X\longrightarrow S$ un morfismo plano y
proyectivo tal que, en la factorización de Stein $f=f''f'$, el morfismo
$f':X\longrightarrow S'$ es plano, con fibras geométricas íntegras, y satisface
por tanto las hipótesis de 3.1, mientras que el morfismo finito
$f'':S'\longrightarrow S$ es plano. Entonces $\operatorname{Pic}_{X/S}$ existe
y, con la notación introducida en [4], II, C 2, se tiene un isomorfismo canónico
\[
\operatorname{Pic}_{X/S}\xrightarrow{\sim}
\Pi_{S'/S}\operatorname{Pic}_{X/S'}.
\]

Para demostrarlo, se establece primero un isomorfismo de funtores
\[
\operatorname{Pic}_{X/S}\xrightarrow{\sim}
\Pi_{S'/S}\operatorname{Pic}_{X/S'},
\]
se aplica 3.1, que implica que $\operatorname{Pic}_{X/S'}$ es representable, y
se usa después la estructura de $\operatorname{Pic}_{X/S'}$ descrita en la
sección 3, que implica que todo subconjunto finito de una fibra de
$\operatorname{Pic}_{X/S'}$ sobre $S$ está contenido en un abierto afín, para
obtener la existencia de $\Pi_{S'/S}\operatorname{Pic}_{X/S'}$.

\src{158}
Por ejemplo, cuando $X$ es un esquema que es unión disjunta de esquemas $X_i$
sobre $S$ que satisfacen las condiciones de 3.1, el enunciado 6.1 se reduce al
enunciado trivial
\[
\operatorname{Pic}_{X/S}\xrightarrow{\sim}
\prod_i\operatorname{Pic}_{X_i/S}.
\]

\result{COROLARIO}{6.2} Sea $f:X\longrightarrow S$ un morfismo proyectivo y
plano con fibras geométricas localmente íntegras (por ejemplo, un morfismo
proyectivo normal). Entonces $\operatorname{Pic}_{X/S}$ existe.

En este caso, además, $S'$ es un recubrimiento étale de $S$; esto ocurre
siempre que $f$ sea separable, es decir, plano con fibras geométricamente
reducidas. Se comprueba que el teorema de estructura enunciado en la sección 3
para $\operatorname{Pic}_{X/S}$ sigue siendo válido, en virtud de la estructura
análoga de $\operatorname{Pic}_{X/S'}$.

Por otra parte, un procedimiento de descenso proporciona un teorema de
existencia relativo cuyo alcance depende de la solución de las cuestiones de
descenso no plano planteadas en [4], I, A (c). Solo explicitaremos el caso
particular siguiente:

\result{PROPOSICIÓN}{6.3} Sea $f:X\longrightarrow S$ un morfismo propio, y sean
$X_1$ y $X_2$ dos subesquemas de $X$, planos sobre $S$, definidos por dos
Ideales coherentes $\mathcal I_1$ y $\mathcal I_2$ tales que
$\mathcal I_1\cap\mathcal I_2=(0)$ y
$\mathcal O_X/(\mathcal I_1+\mathcal I_2)$ es plano sobre $S$; es decir, el
subesquema de $X$ que es el supremo de $X_1$ y $X_2$ es $X$, mientras que su
ínfimo $Z$ es plano sobre $S$.\footnote{La letra final aparece impresa como
«$X$», pero la condición inmediatamente anterior concierne a la planitud de
$\mathcal O_X/(\mathcal I_1+\mathcal I_2)$ sobre $S$ y, por tanto, define $Z$
como plano sobre $S$.} Supongamos además que, para todo $s\in S$, los
homomorfismos
\[
k(s)\longrightarrow H^0(X_{i,s},\mathcal O_{X_{i,s}}),\qquad i=1,2,
\]
son biyectivos. Entonces el homomorfismo natural de funtores
\[
\operatorname{Pic}_{X/S}\longrightarrow
\operatorname{Pic}_{X_1/S}\times\operatorname{Pic}_{X_2/S}
\]
es representable por morfismos afines. Por consiguiente, si
$\operatorname{Pic}_{X_1/S}$ y $\operatorname{Pic}_{X_2/S}$ existen, también
existe $\operatorname{Pic}_{X/S}$, y el morfismo canónico
\[
\operatorname{Pic}_{X/S}\longrightarrow
\operatorname{Pic}_{X_1/S}\times\operatorname{Pic}_{X_2/S}
\]
es afín.

Por descenso fielmente plano se reduce al caso en que $Z$ admite una sección
sobre $S$, que define así secciones de $X$, $X_1$ y $X_2$ sobre $S$ y permite
eliminar los automorfismos de las estructuras consideradas como se explicó en
2.5. La demostración consiste entonces en observar que los datos de un Módulo
invertible

\src{159}
«rigidificado» $\mathcal L$ sobre $X$ equivalen a los datos de una terna
$(\mathcal L_1,\mathcal L_2,u)$, donde $\mathcal L_i$ es un Módulo invertible
«rigidificado» sobre $X_i$ y $u$ es un isomorfismo de $\mathcal L_1|Z$ con
$\mathcal L_2|Z$ compatible con las rigidificaciones. Solo hay que comprobar
que, para $\mathcal L_1$ y $\mathcal L_2$ fijos, el dato de $u$ se expresa por
una sección de un esquema adecuado sobre $S$ que es afín sobre $S$; esto es
fácil. De 6.3 se deduce inmediatamente:

\result{COROLARIO}{6.4} Sea $X$ un esquema propio y separable sobre un cuerpo
$k$, y sean $X_i$ las componentes irreducibles de $X$. Si existen los
$\operatorname{Pic}_{X_i/k}$, también existe $\operatorname{Pic}_{X/k}$, y el
morfismo canónico
\[
\operatorname{Pic}_{X/k}\longrightarrow\prod_i\operatorname{Pic}_{X_i/k}
\]
es afín.

Combinado con 6.2, esto demuestra, por ejemplo, la existencia de
$\operatorname{Pic}_{X/k}$ siempre que $X$ sea un esquema proyectivo y
separable sobre un cuerpo $k$. Cuando $X$ deja de ser separable sobre $k$,
existe asimismo un resultado de reducción que utiliza un argumento debido a
Oort [11]. El método se aplica también sobre un esquema de base arbitrario, caso
útil, por ejemplo, para demostrar en la exposición siguiente el resultado de
finitud anunciado en 3.3. Para evitar un enunciado excesivamente técnico,
limitémonos a un cuerpo de base:

\result{PROPOSICIÓN}{6.5} Sea $X$ un esquema propio sobre un cuerpo $k$, y sea
$X_0$ un subesquema con el mismo conjunto subyacente, definido por tanto por un
Ideal nilpotente sobre $X$. Entonces el morfismo funtorial
\[
\operatorname{Pic}_{X/k}\longrightarrow\operatorname{Pic}_{X_0/k}
\]
es representable por morfismos afines. En particular, si existe
$\operatorname{Pic}_{X_0/k}$, también existe $\operatorname{Pic}_{X/k}$, y el
morfismo $\operatorname{Pic}_{X/k}\longrightarrow\operatorname{Pic}_{X_0/k}$
es afín.

Combinando esto con 6.4 se concluye inmediatamente:

\result{COROLARIO}{6.6} Sea $X$ un preesquema proyectivo sobre un cuerpo $k$.
Entonces $\operatorname{Pic}_{X/k}$ existe.

\result{OBSERVACIÓN}{6.6} Es sumamente plausible que
$\operatorname{Pic}_{X/k}$ exista para todo esquema propio $X$ sobre un cuerpo
$k$.\footnote{La fe de erratas de 1962 informa de que Murre acababa de responder
afirmativamente a esta cuestión.} Los resultados anteriores reducen esta
cuestión al caso en que $k$ es algebraicamente cerrado y $X$ es íntegro. Por el
lema de Chow, existe entonces un esquema íntegro $X'$, proyectivo sobre $k$, y
un morfismo dominante $g:X'\longrightarrow X$. Bastaría, por tanto, mostrar que
el

\src{160}
morfismo funtorial correspondiente
\[
\operatorname{Pic}_{X/k}\longrightarrow\operatorname{Pic}_{X'/k}
\]
es representable, y sin duda representable por morfismos afines, puesto que ya
se sabe que $\operatorname{Pic}_{X'/k}$ es representable. Esto plantea
cuestiones de descenso no plano que por el momento siguen sin resolverse.
Obsérvese que, si se restringe el funtor $\operatorname{Pic}_{X/k}$ a los
preesquemas reducidos, con $X$ propio e íntegro sobre un cuerpo algebraicamente
cerrado $k$, sí se obtiene un funtor representable, como demostró Chevalley [1]
cuando $X$ es normal y como demostró Seshadri [12] en general mediante un
método de descenso. Sin embargo, con nuestra notación, el esquema construido
por estos autores no es $\operatorname{Pic}_{X/k}$, sino
$(\operatorname{Pic}_{X/k})_{\mathrm{red}}$, el esquema reducido correspondiente
a $\operatorname{Pic}_{X/k}$.

\result{OBSERVACIÓN}{6.7} Más generalmente, sea $f:X'\longrightarrow X$ un
morfismo sobreyectivo de preesquemas propios sobre $k$. Consideraciones de
descenso no plano llevan a conjeturar que
\[
\operatorname{Pic}_{X/k}\longrightarrow\operatorname{Pic}_{X'/k}
\]
es un morfismo afín.\footnote{La fe de erratas de 1962 informa de que esta cuestión
tiene respuesta afirmativa y remite al último párrafo de los Comentarios a la
exposición siguiente.} Esto implicaría, en particular, tras dividir por las
componentes conexas de los elementos neutros, que el homomorfismo
correspondiente de grupos de Néron--Severi es inyectivo módulo torsión. Esto se
puede verificar mediante la teoría de intersecciones cuando $X$ y $X'$ son no
singulares. La respuesta no parece conocida en ningún otro caso.

\result{OBSERVACIÓN}{6.8} Contra lo que podría pensarse, los esquemas de Picard
de esquemas algebraicos con elementos nilpotentes son útiles, e incluso
indispensables, en diversas cuestiones. Así, cuando $X$ es un esquema
proyectivo, digamos liso sobre $k$,\footnote{El texto impreso dice «regular»; la
fe de erratas prescribe «liso sobre $k$».} y $Y$ es una sección hiperplana, hay que
considerar los «entornos infinitesimales» $X_n$ de $Y$ de todos los órdenes y
los esquemas de Picard $\operatorname{Pic}_{X_n/k}$. Cuando $X$ es irreducible
de dimensión $\geq4$ (respectivamente $\geq3$), el morfismo canónico
\[
\operatorname{Pic}_{X/k}\longrightarrow\operatorname{Pic}_{X_n/k}
\qquad(n\text{ grande})
\]
es un isomorfismo (respectivamente, induce un isomorfismo sobre las imágenes
inversas de los subgrupos de torsión de Néron--Severi), resultado que será útil
en el estudio cualitativo de los esquemas de Picard en la exposición siguiente.
Asimismo, considerar los esquemas de Picard de ciertas curvas con elementos
nilpotentes, junto con los teoremas fundamentales de geometría formal [3],
permite verificar, en el caso equicaracterístico, una conjetura de Mumford:
para todo anillo local noetheriano normal completo $A$ de dimensión 2, el grupo
de clases de divisores de $A$ puede considerarse como el conjunto de puntos
racionales sobre $k$ de un

\src{161}
grupo algebraico $G$ sobre el cuerpo residual $k$; además, $G$ queda
determinado canónicamente una vez elegido un cuerpo de coeficientes en $A$.
Cuando $A$ tiene dimensión arbitraria, es plausible que exista un grupo
proalgebraico sobre $k$ que desempeñe el mismo papel que $G$. Cuando
$\operatorname{Spec}(A)$ puede «desingularizarse», dicho grupo se construye como
límite proyectivo de esquemas de Picard de esquemas proyectivos adecuados sobre
$k$ con elementos nilpotentes.

\section*{BIBLIOGRAFÍA}
\addcontentsline{toc}{section}{Bibliografía}

\begin{description}
\item[{[1]}] CHEVALLEY (Claude). -- Sur la théorie de Picard,
\emph{Amer. J. of Math.}, vol. 82, 1960, pp. 435--490.

\item[{[2]}] GROTHENDIECK (A.) y DIEUDONNÉ (J.). -- Éléments de géométrie
algébrique, Chapitre I et suivants. -- París, Presses universitaires de France,
1960--1961 (Institut des Hautes Études Scientifiques, Publications
mathématiques 4, 8, \ldots).

\item[{[3]}] GROTHENDIECK (Alexander). -- Géométrie formelle et géométrie
algébrique, \emph{Séminaire Bourbaki}, vol. 11, 1958/59,
n.\textsuperscript{o} 182, 28 pp.

\item[{[4]}] GROTHENDIECK (Alexander). -- Techniques de descente et théorèmes
d'existence en géométrie algébrique, I--IV, \emph{Séminaire Bourbaki}, vol. 12,
1959/60, n.\textsuperscript{o} 190, 29 pp. y n.\textsuperscript{o} 195, 22 pp.;
vol. 13, 1960/61, n.\textsuperscript{o} 212, 20 pp. y
n.\textsuperscript{o} 221, 28 pp.

\item[{[5]}] GROTHENDIECK (Alexander). -- Techniques de construction en
géométrie analytique, I--X, \emph{Séminaire Cartan}, vol. 13, 1960/61:
Déformation des structures analytiques complexes,
n.\textsuperscript{o} 7--17.

\item[{[6]}] GROTHENDIECK (Alexander). -- Séminaire de géométrie algébrique,
editado por Lichtenbaum, Harvard University, 1961 (de próxima aparición).

\item[{[7]}] LANG (Serge). -- \emph{Abelian varieties}. -- Nueva York,
Interscience Publishers, 1959 (Interscience Tracts in pure and applied
Mathematics, 7).

\item[{[8]}] MUMFORD (D.). -- An elementary theorem in geometric invariant
theory, \emph{Bull. Amer. math. Soc.}, vol. 67, 1961, pp. 483--486.

\item[{[9]}] MUMFORD (D.). -- The topology of normal singularities of an
algebraic surface and a criterion for simplicity. -- París, Presses
universitaires de France, 1961 (Institut des Hautes Études Scientifiques,
Publications mathématiques, 9).

\item[{[10]}] MUMFORD (D.) y TATE (J.). -- Séminaire de géométrie algébrique,
Harvard University, Spring Term 1962 (de próxima aparición).

\item[{[11]}] OORT (F.). -- Sur le schéma de Picard,
\emph{Bull. Soc. math. France}, vol. 90, 1962 (de próxima aparición).

\item[{[12]}] SESHADRI (C. S.). -- Tesis, de próxima aparición en
\emph{Annali di Matematica}.
\end{description}

\def\FgaUnitKey{e232}\clearpage
\chapter*{Fe de erratas de la exposición n.\textsuperscript{o} 232}
\addcontentsline{toc}{chapter}{Fe de erratas de la exposición n.\textsuperscript{o} 232}
\src{303}

\noindent\textbf{7. TDTE V: Esquemas de Picard. Teoremas de existencia.}

\begin{center}
[Seminario Bourbaki, vol. 14, 1961/62, n.\textsuperscript{o} 232, 19 pp.]
\end{center}

\noindent\textbf{Página 7, observación 3.3.} --- La cuestión planteada aquí fue
resuelta afirmativamente por MUMFORD (véanse los Comentarios a la exposición
siguiente).

\noindent\textbf{Página 13, observación 5.2.} --- Como señalo al comienzo de la
exposición siguiente, la conjetura de existencia formulada aquí es falsa; no
obstante, MUMFORD consigue demostrar por sus métodos un teorema ligeramente más
débil.

\noindent\textbf{Página 14, línea 7.} --- Después de «completo», añádase «con
cuerpo residual algebraicamente cerrado». El contraejemplo de MUMFORD recién
mencionado muestra, además, que esta restricción es indispensable.

\noindent\textbf{Página 17, observación 6.6.} --- Como señalamos al comienzo de
la exposición siguiente, la cuestión planteada aquí acaba de ser resuelta
afirmativamente por MURRE.

\noindent\textbf{Página 18, observación 6.7.} --- La cuestión planteada aquí
tiene respuesta afirmativa (cf. el último párrafo de los Comentarios a la
exposición siguiente).

\noindent\textbf{Página 18, observación 6.8, cuarta línea.} --- En lugar de
«regular», léase «liso sobre $k$».

\def\FgaUnitKey{236}\unit{Exposición 236 -- Esquemas de Picard: propiedades generales}
\sid{236}
\src{221}
\seminares{14.º año, 1961/62, n.º 236\\Mayo de 1962}
\paperhead{TÉCNICA DE DESCENSO Y TEOREMAS DE EXISTENCIA EN GEOMETRÍA ALGEBRAICA\\
VI. ESQUEMAS DE PICARD: PROPIEDADES GENERALES}%
  {por Alexander GROTHENDIECK}

\fsection{0}{Complementos a la exposición anterior ([3], V)\textsuperscript{(1)}}

Se han producido algunos avances en las cuestiones de existencia de los
preesquemas de Picard planteadas en [3], V:

\medskip
\noindent a. (Mumford). No es cierto en general que, cuando
$f:X\longrightarrow S$ es un morfismo proyectivo separable (= plano con fibras
separables), exista el preesquema de Picard $\operatorname{Pic}_{X/S}$, ni
siquiera si las fibras de $f$ tienen dimensión 1 y $S$ es el espectro de un
anillo de valoración discreta completo. Se obtiene un ejemplo tomando
$S=\operatorname{Spec}\mathbf R[[t]]$ y tomando como $X$ el subesquema de
$\mathbf P^2_S$ definido por la ecuación (donde $x,y,z$ son variables
homogéneas)
\[
x^2+y^2=tz^2.
\]
Se trata de una cónica que degenera geométricamente en dos rectas que se
cortan, aunque la fibra especial sobre $\mathbf R$ es irreducible (está dada
por $x^2+y^2=0$ sobre $\mathbf R$). Se ve fácilmente que, tras la extensión
étale $S'\longrightarrow S$, con $S'=\operatorname{Spec}\mathbf C[[t]]$, existe
el preesquema de Picard de $X'/S'$ y admite una descripción explícita como
unión disjunta de copias de $\widetilde S$, donde $\widetilde S$ se obtiene de
$S$ duplicando el origen un número infinito de veces. Se comprueba fácilmente
que el dato de descenso sobre $\operatorname{Pic}_{X'/S'}$ para
$S'\longrightarrow S$ (dado aquí por la acción del grupo de Galois
$\mathbf Z/2\mathbf Z$ de $S'$ sobre $S$) no es efectivo: el grupo permuta
ciertos puntos duplicados, de modo que algunas órbitas no están contenidas en
ningún abierto afín. Mumford puede demostrar, sin embargo, que si
$f:X\longrightarrow S$ es un morfismo proyectivo separable tal que, para todo
$s\in S$, las componentes irreducibles de $X_s$ son geométricamente
irreducibles sobre $k(s)$, entonces existe $\operatorname{Pic}_{X/S}$. Su
demostración se apoya en un refinamiento de su teorema del cociente; cf. [7].
Obsérvese, por otra parte, que sin ninguna hipótesis sobre las componentes
irreducibles de las fibras $X_s$, sigue siendo posible que exista, no obstante,
el esquema $\operatorname{Pic}^{\tau}_{X/S}$ (que se introducirá más adelante).

\src{222}
\noindent b. (Murre). Sea $X$ un esquema propio sobre un cuerpo. Entonces
existe $\operatorname{Pic}_{X/k}$. La demostración sigue en parte la de
Chevalley [1] y utiliza de manera esencial la estructura de grupo del funtor de
Picard.

Para algunos complementos relativos a la teoría de los esquemas de Picard, en
particular en relación con los esquemas abelianos, será útil consultar [7]. Por
último, una laguna notable de la presente exposición es la ausencia de
«criterios de equivalencia» que permitan comparar el esquema de Picard de un
esquema proyectivo con los de sus secciones hiperplanas. Los teoremas clave
para desarrollar tales criterios se encuentran en [5], a los cuales hay que
añadir los teoremas de existencia para los esquemas de Picard.

\fsection{1}{Propiedades topológicas de los preesquemas conmutativos en grupos}

Sea primero $k$ un cuerpo y $G$ un preesquema en grupos sobre $k$. Como el
elemento neutro $e$, al ser racional sobre $k$, es necesariamente cerrado, se
deduce inmediatamente que la diagonal de $G\times_kG$ es cerrada y, por tanto,
que $G$ es separado: todo preesquema en grupos sobre un cuerpo es separado.
Denotemos por $G^\circ$ la componente conexa del elemento neutro $e$. Como $e$
es racional sobre $k$, $G^\circ$ es de hecho geométricamente conexa; es decir,
la formación de $G^\circ$ es compatible con la extensión del cuerpo de base.
Se deduce asimismo que $G^\circ$ es estable por multiplicación como conjunto y,
cuando $G$ es localmente noetheriano, de modo que $G^\circ$ es abierto, puede
considerarse un subgrupo abierto de $G$. Supongamos en lo sucesivo que $G$ es
localmente de tipo finito sobre $k$. Entonces $G^\circ$ es geométricamente
irreducible y de tipo finito sobre $k$. En efecto, se puede suponer que $k$ es
algebraicamente cerrado y después que $G$ es reducido (pues
$G_{\mathrm{red}}$ es entonces un subgrupo de $G$, ya que
$G_{\mathrm{red}}\times_kG_{\mathrm{red}}$ vuelve a ser reducido); por tanto,
es liso sobre $k$ en un abierto no vacío y, mediante traslación de ese
abierto, en todas partes. Así, $G$ es localmente irreducible, por lo que sus
componentes irreducibles coinciden con sus componentes conexas y $G^\circ$ es
irreducible. Sea $U$ un entorno afín de $e$ en $G^\circ$. Utilizando la
irreducibilidad de $G^\circ$, se ve inmediatamente que $U\cdot U=G^\circ$, lo
que demuestra que $G^\circ$ es cuasicompacto y, por tanto, de tipo finito sobre
$k$.

Para simplificar, supongamos que $G$ es conmutativo. Para todo entero $n>0$,
sea $G^{(n)}$ la imagen inversa de $G^\circ$ por el homomorfismo de potencia
$n$-ésima $\varphi_n$ de $G$; así, $G^{(n)}$ es un subgrupo abierto de $G$.
Pongamos
\[
G^\tau=\bigcup_nG^{(n)}.
\]

\src{223}
\[
G^\sigma=\bigcup_{(n,p)=1}G^{(n)},
\qquad
G^\rho=\bigcup_hG^{(p^h)},
\]
donde $p$ es el exponente característico del cuerpo $k$. Estos son subgrupos
abiertos de $G$ que satisfacen
\[
G^\sigma\cap G^\rho=G^\circ,
\qquad
G^\sigma G^\rho=G^\tau.
\]

\result{OBSERVACIÓN}{} Se puede construir el esquema en grupos cociente
$G/G^\circ=N$ (cf. [3], IV), y definir entonces $G^\tau$, $G^\sigma$ y
$G^\rho$ como las imágenes inversas en $G$, respectivamente, del subgrupo de
torsión de $N$, de la suma de sus componentes $q$-primarias para los primos
$q\neq p$ y de su componente $p$-primaria.\footnote{El texto impreso intercambia
las dos últimas correspondencias: asigna a $G^\sigma$ la componente
$p$-primaria y a $G^\rho$ su complemento primo con $p$. Las definiciones
precedentes de $G^\sigma$ y $G^\rho$ obligan al orden corregido que se da aquí.}
Obsérvese a este respecto que $N$ es un esquema en
grupos discreto separable sobre $k$. Así, tras elegir una clausura algebraica
$\overline k$ de $k$, con grupo de Galois $\pi$, se identifica con un grupo
discreto ordinario sobre el cual $\pi$ actúa por automorfismos. Esto da la
interpretación evidente de la construcción del subgrupo de torsión y de su
descomposición en componentes $q$-primarias. Cuando $G$ es el esquema de
Picard de un esquema propio $X$ sobre $k$, se puede llamar a $N$ el
\emph{esquema de Néron--Severi} (reducido) de $X$ sobre $k$. Cuando
$G^\circ_{\mathrm{red}}$ es un subesquema en grupos de $G^\circ$, lo que ocurre
en particular siempre que $k$ sea perfecto o $G^\circ$ sea propio sobre $k$
(por ejemplo, cuando $X$ sea geométricamente normal), conviene introducir
también el cociente $N'=G/G^\circ_{\mathrm{red}}$, que tiende a comportarse
mejor que $N$ por especialización, es decir, cuando $X$ varía en una familia
de esquemas algebraicos.

Sea ahora $S$ un preesquema localmente noetheriano y sea $G$ un preesquema en
grupos sobre $S$, localmente de tipo finito sobre $S$. No suponemos que $G$ sea
de tipo finito sobre $S$ ni separado sobre $S$. Pongamos
\[
G^\circ=\bigcup_{s\in S}(G_s)^\circ,
\]
y, cuando $G$ sea conmutativo,
\[
G^\tau=\bigcup_{s\in S}(G_s)^\tau,
\qquad
G^\sigma=\bigcup_{s\in S}(G_s)^\sigma,
\qquad
G^\rho=\bigcup_{s\in S}(G_s)^\rho.
\]

\src{224}
Estos son subconjuntos de $G$ estables por la multiplicación en $G$, lo que
evidentemente no implica que puedan ser definidos por subpreesquemas en grupos
de $G$. En particular, parece que en general no existe un subpreesquema en
grupos de $G$ cuyo conjunto subyacente sea $G^\circ$. Por supuesto, si uno de
estos conjuntos es abierto, entonces, con su estructura inducida, es un
subpreesquema abierto en grupos de $G$. Veremos que esto siempre ocurre para
$G^\tau$. Así, desde el punto de vista de los funtores representables, y en
particular desde el punto de vista de la especialización, la equivalencia
numérica se comporta de manera más satisfactoria que la equivalencia
algebraica. He aquí las principales propiedades generales de los conjuntos que
acabamos de definir:

\result{TEOREMA}{1.1} Los conjuntos $G^\circ$, $G^\tau$, $G^\sigma$ y $G^\rho$
son localmente constructibles. Además:

\noindent (i) $G^\circ$ es cuasicompacto sobre $S$. Cuando los $G_s^\circ$ son
propios y $G$ es separado sobre $S$, $G^\circ$ es propio sobre $S$ y, por
tanto, cerrado en $G$.

\noindent (ii) $G^\tau$ es abierto. Si $G^\circ$ es cerrado, también lo es
$G^\tau$.

\noindent (iii) Si $G^\circ$ es cerrado, también lo es $G^\sigma$, siempre que
se esté en característica igual, es decir, que todos los cuerpos residuales de
$S$ tengan la misma característica. Si $G^\circ$ es cerrado y
$G\longrightarrow S$ es universalmente abierto en los puntos de $G^\sigma$
(cf. Corolario 1.5 más adelante), entonces $G^\sigma\longrightarrow S$ es
universalmente abierto.

\noindent (iv) Si $G^\circ$ es cerrado, también lo es $G^\rho$. Supongamos que
se está en característica igual y que, para todo entero $n>0$ con $(n,p)=1$,
el homomorfismo de potencia $n$-ésima de $G$ es abierto. Entonces $G^\rho$ es
abierto.

Damos indicaciones para la demostración. La constructibilidad local de
$G^\circ$ y su cuasicompacidad sobre $S$ están contenidas en el lema siguiente:

\result{LEMA}{1.2} Supongamos $S\neq\varnothing$. Entonces existen un abierto
no vacío $U$ de $S$, un esquema en grupos $H$ de tipo finito sobre $U$ con
fibras conexas y un homomorfismo de esquemas en grupos
$H\longrightarrow G|U$ que es una inmersión abierta de imagen $G^\circ|U$.

Para demostrar el lema, se puede suponer que $S$ es irreducible y denotar por
$\eta$ su punto genérico. Tras el cambio de base
$S'=\operatorname{Spec}(\mathcal O_{S,\eta})\longrightarrow S$, se obtiene un
esquema en grupos $G'$ sobre el anillo local artiniano
$\mathcal O_{S,\eta}=A$, con un subesquema abierto en grupos $G'{}^\circ$ de
tipo finito sobre $A$, como se ha indicado antes. Por tanto, este procede de un
esquema en grupos $H$ de tipo finito sobre un entorno abierto $U$ de $\eta$, y
la inmersión canónica $G'{}^\circ\longrightarrow G'$ procede de una

\src{225}
inmersión abierta $H\longrightarrow G|U$, que es un homomorfismo de esquemas en
grupos después de hacer $U$ suficientemente pequeño. Si $U$ es suficientemente
pequeño, las fibras de $H$ son conexas y tienen todas la misma dimensión,
igual a la de las fibras de $G$. De ello se sigue que, para todo $s\in U$, la
imagen de $H_s$ en $G_s$ es exactamente $G_s^\circ$ (tras reducir de nuevo
$U$), lo que demuestra 1.2.

La segunda afirmación de (i) está contenida en el lema siguiente, aplicado a
un entorno abierto cuasicompacto de $G^\circ$ en $G$:

\result{LEMA}{1.3} Sea $X$ un preesquema separado de tipo finito sobre un
preesquema localmente noetheriano $S$, sea $g$ una sección de $X$ sobre $S$ y
sea $X^\circ$ la unión, en los $X_s$, de las componentes conexas que contienen
$g(s)$. Si $s\in S$ y $X_s^\circ$ es propio sobre $k(s)$, existe un entorno
abierto $U$ de $s$ tal que $X^\circ|U$ es propio sobre $U$ y, por tanto,
cerrado en $X|U$.

Por descenso fielmente plano sobre la base, se reduce al caso en que $S$ es el
espectro de un anillo local completo y $s$ es su punto cerrado. Aplicando [2],
III, 5.5.1, se ve que $X$ es la unión disjunta de dos abiertos $X'$ y $X''$, el
primero propio sobre $S$ y tal que $X'_s=X_s^\circ$. Esto reduce al caso
$X=X'$, es decir, a $X$ propio sobre $S$. Se utiliza entonces un argumento
estándar basado en el criterio valuativo de ser propio para un subconjunto
(omitido en [2], capítulo II).

Demostremos ahora que $G^\tau$ es abierto. Equivalentemente, puesto que la
formación de $G^\tau$ (así como la de $G^\circ$, $G^\sigma$ y $G^\rho$) conmuta
con la extensión de la base, para toda sección $g$ de $G$ sobre $S$, el
conjunto $g^{-1}(G^\tau)$ es abierto. Esto significa dos cosas:

\medskip
\noindent a. Si $y\in S$ y $g(y)\in G_y^\tau$, entonces
$g(y')\in G_{y'}^\tau$ para todo $y'\in\overline{\{y\}}$ suficientemente
próximo a $y$.

\medskip
\noindent b. Si $y'\in S$ y $g(y')\in G_{y'}^\tau$, entonces
$g(y)\in G_y^\tau$ para toda generalización $y$ de $y'$.

Para (a), existe un entero $n>0$ tal que $g^n(y)\in G_y^\circ$. Como
$G^\circ$ es constructible, $(g^n)^{-1}(G^\circ)$ es constructible y, por
tanto, $g^n(y')\in G_{y'}^\circ$ para
$y'\in\overline{\{y\}}$ suficientemente próximo a $y$. Para (b), las
potencias $g(y')^n=g^n(y')$ permanecen en un abierto cuasicompacto de
$G_{y'}$, pues están contenidas en un número finito de clases módulo
$G_{y'}^\circ$. Por consiguiente, existe un abierto cuasicompacto $U$ de $G$
que contiene los $g^n(y')$ y, por tanto, también sus generalizaciones
$g^n(y)$. Así, las potencias de $g(y)$ permanecen en un abierto cuasicompacto
de $G_y$, lo que implica fácilmente que $g(y)\in G_y^\tau$.

\src{226}
Supongamos que $G^\circ$ es cerrado; demostremos lo mismo para $G^\tau$. Como
ya se sabe que $G^\tau$ es abierto y, por tanto, localmente constructible,
basta demostrar que es estable por especialización. Esto se sigue de que es
una unión de conjuntos cerrados, a saber, las imágenes inversas del conjunto
cerrado $G^\circ$ por los homomorfismos de potencia $n$-ésima $\varphi_n$.

El mismo argumento demuestra que $G^\sigma$ y $G^\rho$ son cerrados si
$G^\circ$ es cerrado (en el primer caso, bajo la condición de característica
igual), una vez demostrado que $G^\sigma$ y $G^\rho$ son localmente
constructibles. Para $x\in G^\tau$, sea $\nu(x)$ el menor entero $n>0$ tal que
el homomorfismo de potencia $n$-ésima $\varphi_n$ envía $x$ a $G^\circ$.
Entonces $G^\sigma$ (respectivamente, $G^\rho$) está formado por los
$x\in G$\footnote{El texto impreso dice «$X$», que no está definido en este
pasaje; los cuatro loci considerados son subconjuntos de $G$.} para los cuales
$\nu(x)$ es primo con $p$ (respectivamente, una potencia de $p$). La afirmación
de constructibilidad se deduce ahora del enunciado más preciso:

\result{LEMA}{1.4} La función $\nu$ sobre $G^\tau$ es localmente
constructible.

En efecto, para todo entero $n>0$, el conjunto de los $x\in G^\tau$ tales que
$\nu(x)=n$ es localmente constructible: es la diferencia entre
$\varphi_n^{-1}(G^\circ)$ y la unión de los
$\varphi_d^{-1}(G^\circ)$, donde $d$ recorre los divisores propios de $n$.
Como $G^\circ$ es localmente constructible, también lo son los
$\varphi_d^{-1}(G^\circ)$ y, por consiguiente, también lo es la diferencia
anterior.

Supongamos ahora que $G^\circ$ es cerrado y que $G\longrightarrow S$ es
universalmente abierto en los puntos de $G^\sigma$. Demostremos que
$G^\sigma\longrightarrow S$ es abierto, es decir, que envía un entorno en
$G^\sigma$ de $x\in G^\sigma$ sobre un entorno de $y=f(x)$. Como $G^\sigma$ es
localmente constructible, basta demostrar que para toda generalización $y'$ de
$y$ existe una generalización $x'$ de $x$ en $G^\sigma$ situada sobre $y'$.
Por cambio de base, se reduce al caso en que $S$ es el espectro de un anillo de
valoración discreta y $y,y'$ son, respectivamente, sus puntos cerrado y
genérico. Como $G\longrightarrow S$ es abierto en $x$, existe una
generalización $x_1$ de $x$ en $G$ sobre $y'$. Tras otro cambio de base, se
puede suponer además que existe una sección $g$ de $G$ sobre $S$ con
$x=g(y)$. Si $k(y')$ tiene característica cero, tómese cualquier
generalización $x'$ de $x$ en $G$ sobre $y'$. Pertenece a $G^\tau$ porque
$G^\tau$ es abierto y, por tanto, a $G^\sigma$ porque
$G_{y'}^\sigma=G_{y'}^\tau$. Si $k(y')$ tiene característica $p>0$, escribamos
\[
\nu(g(y'))=p^hm\qquad\text{con }(m,p)=1;
\]

\src{227}
elijamos enteros $a$ y $b$ tales que $ap^h+bm=1$, y pongamos
\[
g_1=g^{ap^h},\qquad g_2=g^{bm},\qquad\text{de modo que}\qquad g=g_1g_2.
\]
Por construcción, $g_1(y')\in G^\sigma$ y $g_2(y')\in G^\rho$. Como
$G^\circ$ es cerrado, se sigue que $g_1(y)\in G^\sigma$ y
$g_2(y)\in G^\rho$. Pero
\[
g(y)=g_1(y)g_2(y)\in G^\sigma,
\]
por lo que $g_2(y)\in G^\sigma$ y, por consiguiente,
\[
g_2(y)\in G_y^\sigma\cap G_y^\rho=G_y^\circ.
\]
La hipótesis, junto con el hecho de que $S$ es el espectro de un anillo de
valoración discreta, implica que $G-(G_y-G_y^\circ)$ es un abierto de $G$ en
el que $G\longrightarrow S$ induce un morfismo abierto. Así, por todo punto de
$G_y^\circ$ pasa una «cuasisección». Después de efectuar una extensión
adicional de la base $S$, se puede suponer, por consiguiente, que existe una
sección $g'_2$ de $G^\circ$ sobre $S$ tal que $g'_2(y)=g_2(y)$. Pongamos
$g'=g_1g'_2$. Entonces, por construcción, $g'(y)=g(y)=x$, y
\[
g'(y')=g_1(y')g'_2(y')\in G^\sigma.
\]
Por tanto, $g'(y')=x'$\footnote{El texto impreso dice «$g(y')=x'$». La sección
construida en la frase anterior, cuya fibra genérica se calcula, es $g'$.} es
una generalización de $x$ en $G^\sigma$ situada sobre $y'$. Esto demuestra
(iii).

Queda por demostrar la última afirmación, (iv). Debemos probar que, si
$x\in G^\rho$, toda generalización $x'$ de $x$ pertenece a $G^\rho$. Tomando
imágenes por $\varphi_{p^h}$ para un $h$ adecuado\footnote{El texto impreso
dice $\varphi_{n^h}$. Para $x\in G^\rho$, la definición de la p.~223 da una
potencia del exponente característico $p$ que envía $x$ a $G^\circ$.}, se
puede suponer que, de hecho, $x\in G^\circ$. Para todo entero $n$ primo con la
característica, $x$ pertenece entonces a la imagen de $\varphi_n$ (pues la
aplicación de potencia $n$-ésima en un grupo conexo de tipo finito sobre un
cuerpo de característica prima con $n$ es sobreyectiva). Como $\varphi_n$ es
abierto, $x'$ pertenece también a la imagen de $\varphi_n$. Más precisamente,
sea $U$ un abierto cuasicompacto de $G^\tau$ que contenga $G_y^\circ$.
Entonces $x'\in\varphi_n(U)$ para todo $n$ primo con $p$. Tomando por $n$ un
múltiplo común de las partes primas con $p$ de los $\nu(z)$, para $z\in U$, se
obtiene $x'\in G^\rho$.

La afirmación (1.1), (iii), se completa mediante el corolario siguiente:

\result{COROLARIO}{1.5} Sea $n>1$ un entero tal que el homomorfismo de potencia
$n$-ésima $\varphi_n:G\longrightarrow G$ sea universalmente abierto (por
ejemplo, étale), pongamos
\[
G^{(n)}=\varphi_n^{-1}(G^\circ),
\]
y supongamos que, en cada fibra geométrica $G_{\bar s}$, el subconjunto
$\bigcup_{h\geq0}|{}_{n^h}G_{\bar s}|$ es denso en
$\bigcup_{h\geq0}G_{\bar s}^{(n^h)}$. Entonces

\src{228}
$G\longrightarrow S$ es universalmente abierto en los puntos de los
$G^{(n^h)}$. En particular, si los $G_s^\circ$ no contienen ninguna componente
aditiva y, para todo entero $n>1$, el homomorfismo $\varphi_n$ es
universalmente abierto en los puntos de característica residual prima con
$n$\footnote{El texto impreso dice «prima con $p$». En este enunciado
cuantificado por $n$, el criterio para que un morfismo sea étale y el argumento que utiliza los
núcleos $ {}_nG$ requieren característica residual prima con $n$.}, entonces
$G\longrightarrow S$ es universalmente abierto en los puntos de $G^\sigma$.
Por tanto, según (1.1), (iii), si $G^\circ$ es cerrado,
$G^\sigma\longrightarrow S$ es universalmente abierto. En estos casos,
$s\longmapsto\dim G_s$ es localmente constante sobre $S$.

En efecto, la hipótesis implica que el núcleo $ {}_nG$ de $\varphi_n$ es
universalmente abierto sobre $S$. Por tanto, $G\longrightarrow S$ es
universalmente abierto en los puntos de $ {}_nG$ y, reemplazando $n$ por
$n^h$, en los puntos de $ {}_{n^h}G$. La hipótesis de densidad dice exactamente que, fibra a fibra, los puntos cuyo
orden es una potencia de $n$ son densos en la unión de los $G^{(n^h)}$. Para
la afirmación particular, la ausencia de componentes aditivas implica esta
densidad para los enteros $n$ invertibles en el campo residual. Se deduce fácilmente que
$G\longrightarrow S$ es universalmente abierto en cada punto de esta unión,
en particular a lo largo de la sección neutra, y, por tanto, que
$s\longmapsto\dim G_s$ es localmente constante.

\result{OBSERVACIÓN}{1.6} Recuérdese que un morfismo $X\longrightarrow S$ se
llama universalmente abierto si envía todo abierto a un abierto y conserva
esta admirable propiedad después de cualquier cambio de base
$S'\longrightarrow S$. Esto también significa (cuando $S$ es localmente
noetheriano y $X\longrightarrow S$ es localmente de tipo finito) que toda
componente irreducible de $X$ domina $S$ y que esta propiedad se conserva
después de cualquier cambio de base $S'\longrightarrow S$. En ambas
afirmaciones basta (bajo las hipótesis de finitud anteriores) comprobar los
cambios de base $S'\longrightarrow S$ en los que $S'$ es el espectro de un
anillo de valoración discreta (completo y con cuerpo residual algebraicamente
cerrado, si se desea~...). La definición se extiende de manera evidente a un
subconjunto $Z$ de $X$ (como el subconjunto $G^\sigma$ de $G$). Es un fenómeno
perfectamente normal, incluso para un morfismo proyectivo liso
$X\longrightarrow S$ con fibras geométricamente conexas (por ejemplo, el
producto fibrado de la familia modular de curvas elípticas consigo misma sobre
$S=\operatorname{Spec}\mathbf C[j]$), que $\operatorname{Pic}_{X/S}$ no sea
universalmente abierto sobre $S$: una componente irreducible de
$\operatorname{Pic}_{X/S}$ puede estar situada enteramente sobre un único
punto de $S$. Esto está relacionado con el hecho de que el rango del grupo de
Néron--Severi de las fibras de $X/S$ puede aumentar bruscamente (fenómenos de
«multiplicación compleja»). Por otra parte, (1.5) nos asegura que, en casos
favorables, $\operatorname{Pic}_{X/S}^\sigma$ (y parece que muy a menudo incluso
$\operatorname{Pic}_{X/S}^\tau$) es universalmente abierto sobre $S$.

\src{229}
Por último, he aquí un caso útil en el que, excepcionalmente, $G^\circ$
consiente en ser abierto:

\result{COROLARIO}{1.7} Supongamos que $G\longrightarrow S$ es universalmente
abierto en los puntos de $G^\circ$ (cf. (1.5)) y que las fibras $G_s$ son
separables y, por tanto, lisas sobre $k(s)$ (esta última condición se satisface
automáticamente en característica residual cero, por un teorema de Cartier).
Entonces $G^\circ$ es abierto en $G$. Si, además, $S$ es reducido, entonces
$G^\circ$ es liso (en particular, plano) sobre $S$.

La primera afirmación puede precisarse del modo siguiente. Si, para un
$s\in S$ dado, $G\longrightarrow S$ es universalmente abierto en los puntos de
$G_s^\circ$ y $G_s^\circ$ es separable sobre $k(s)$, entonces $G^\circ$ es un
entorno de $G_s^\circ$ (y la hipótesis hecha en $s$ seguirá siendo válida en
los puntos próximos). Este enunciado, independiente de toda estructura de
grupo, se demostrará en [2], IV, §7. La última afirmación de (1.7), también
independiente de toda estructura de grupo y de toda cuestión relativa a una
componente conexa, es un caso particular de un criterio de planitud dado en
[2], IV, §5, que implica de manera más general:

\result{COROLARIO}{1.8} Sea $U$ un abierto de una fibra $G_s$ tal que
$G\longrightarrow S$ sea universalmente abierto en los puntos de $U$ (cf.
(1.5)). Si $G_s$ es separable sobre $k(s)$ y $S$ es reducido en $s$, entonces
$G$ es plano sobre $S$ en los puntos de $U$ (y, por tanto, en este caso, liso
sobre $S$ en los puntos de $U$).

\result{OBSERVACIONES}{1.9} No sé si, cuando $G$ es separado sobre $S$,
$G^\circ$ o $G^\tau$ es siempre cerrado en $G$; parece poco probable. En todo
caso, hay contraejemplos evidentes si se suprime la hipótesis de separación.
Por ejemplo, tómese la recta afín con el origen duplicado una infinidad
numerable de veces. Esto da un esquema en grupos $G$ sobre la recta afín cuyas
fibras se reducen todas al grupo neutro salvo una, cuya fibra es $\mathbf Z$.
Este esquema en grupos es un subesquema abierto en grupos, a saber, la clausura
de la sección neutra, del preesquema de Picard del $S$-esquema $X$
correspondiente a una familia de cónicas que degenera en dos rectas que se
cortan.

Además, incluso partiendo de un esquema en grupos finito y plano $G$ sobre el
espectro $S$ de un anillo de valoración discreta $V$, de modo que $G^\circ$ se
reduce a la sección neutra y es, por tanto, cerrado, varias conclusiones pasan
a ser falsas si se suprimen ciertas hipótesis. Supongamos que $V$ tiene
característica igual $p>0$ y sea $G$ el núcleo del homomorfismo
$\mathbf G_a\longrightarrow\mathbf G_a$ definido por el homomorfismo de
funtores
\[
f\longmapsto f^p-tf,
\]
donde $t$ es un uniformizante de $V$.

\src{230}
(Recuérdese que, por definición, el «grupo aditivo» $\mathbf G_a$ sobre $S$
representa el funtor $S'\longmapsto\Gamma(S',\mathcal O_{S'})$.) Así, como
$S$-esquema,
\[
G=\operatorname{Spec}V[x]/(x^p-tx)
\]
es la unión de $p$ «rectas» que se cortan, y representa un grupo
$\mathbf Z/p\mathbf Z$ que degenera en un grupo infinitesimal de tipo aditivo.
En este ejemplo, $G^\circ=G^\sigma$ (y es una de las $p$ rectas), por lo que no
es abierto en $G$. En característica mixta de característica residual $p>0$,
partamos del esquema en grupos $H=\boldsymbol\mu_p$, núcleo de la aplicación de
potencia $p$-ésima en $\mathbf G_m$. Así,
\[
H=\operatorname{Spec}V[x]/(x^p-1);
\]
Este es de nuevo la unión de $p$ rectas que se cortan, y representa un grupo
$\mathbf Z/p\mathbf Z$ (en característica cero) que degenera en un grupo
infinitesimal de tipo multiplicativo en característica $p$. Sea $H'$ el
«esquema en grupos constante» definido por el grupo finito ordinario
$\mathbf Z/p\mathbf Z$, es decir, la unión disjunta de $\mathbf Z/p\mathbf Z$
copias de $S$, o equivalentemente
$\operatorname{Spec}V^{\mathbf Z/p\mathbf Z}$. Entonces $G=H\times_S H'$
describe un grupo $(\mathbf Z/p\mathbf Z)^2$ en característica cero, que
degenera en característica $p$ en el producto de un grupo infinitesimal por
$\mathbf Z/p\mathbf Z$. Aquí $G^\rho$ es la unión de la sección neutra y la
fibra especial, de modo que $G^\rho$ no es abierto, mientras que $G^\sigma$ es
la unión de la sección neutra y la fibra genérica, de modo que no es cerrado,
contrariamente a lo afirmado en los casos de característica igual de (1.1),
(iii) y (iv). Se trata, naturalmente, de fenómenos relacionados con la
característica $p>0$. Los resultados anteriores dan:

\result{COROLARIO}{1.10} Supongamos que todas las características residuales de
$S$ son cero, de modo que $G^\tau=G^\sigma$ y $G^\circ=G^\rho$. Entonces
$G^\tau=G^\sigma$ es abierto; es incluso abierto y cerrado si $G$ es separado
sobre $S$ y los $G_s^\circ$ son propios. Bajo la misma hipótesis,
$G^\circ=G^\rho$ es propio sobre $S$ y, por tanto, cerrado en $G$. Supongamos,
por último, que para todo entero $n>1$ el homomorfismo de potencia $n$-ésima de
$G$ es universalmente abierto. Entonces $G^\circ$ es abierto; y si, además,
los $G_s^\circ$ no tienen componente aditiva (por ejemplo, si los
$G_s^\circ$ son propios, como antes), entonces $G^\tau\longrightarrow S$ es
universalmente abierto, e incluso liso si $S$ es reducido.

Consignemos finalmente el siguiente resultado fácil:

\result{PROPOSICIÓN}{1.11} Existe un abierto $U$ de $S$ tal que el conjunto
$G'$ de los puntos de $G$ en los que $G$ es liso (respectivamente, plano) sobre
$S$ es el conjunto subyacente de un subesquema abierto en grupos inducido por
$G|U$. Además, toda sección de $G$ sobre $U$ es una sección de $G'$ sobre $U$.

\src{231}
\result{COROLARIO}{1.12} Si $G$ es liso (respectivamente, plano) sobre $S$ en
los puntos de la sección neutra, también lo es en los puntos de toda sección de
$G$ sobre $S$ y en todo punto de $G^\circ$. Si, además, para todo entero $n>0$
el homomorfismo de potencia $n$-ésima
$\varphi_n:G\longrightarrow G$ es étale en los puntos de característica prima
con $n$, entonces $G$ es liso (respectivamente, plano) sobre $S$ en todo punto
de $G^\sigma$.

\fsection{2}{Aplicación a las propiedades locales de los esquemas de Picard}

\result{TEOREMA}{2.1}

\noindent(i) Sea $f:X\longrightarrow S$ un morfismo propio y liso, y
supongamos que existe $\operatorname{Pic}_{X/S}$ (por ejemplo, supongamos que
$f$ es proyectivo). Entonces $\operatorname{Pic}_{X/S}$ es separado sobre $S$
y todo subconjunto cerrado $Z$ de $\operatorname{Pic}_{X/S}$ que sea de tipo
finito sobre $S$ es propio sobre $S$.

\noindent(ii) Sea $X$ un preesquema propio geométricamente normal sobre un
cuerpo $k$. Entonces $\operatorname{Pic}_{X/k}^\circ$\footnote{El texto
impreso dice $\operatorname{Pic}_{X/S}^\circ$, aunque $S$ no aparece en este
enunciado sobre el cuerpo $k$.} es propio sobre $k$.

\begin{proof}
(i) Los criterios valuativos ([2], II, §7) nos reducen a demostrar lo
siguiente. Sea $S$ el espectro de un anillo de valoración discreta completo y
sea $U$ el abierto formado por el punto genérico de $S$. Toda sección racional
de $\operatorname{Pic}_{X/S}$ sobre $S$, es decir, toda sección sobre $U$, se
prolonga de manera única a una sección sobre $S$. Por la definición de
$\operatorname{Pic}_{X/S}$, esto equivale al enunciado siguiente: para todo
Módulo invertible $\mathcal L$ sobre $V=f^{-1}(U)$ existe un Módulo invertible sobre $X$
que prolonga $\mathcal L$, único salvo isomorfismo. Esto se deduce fácilmente
de la descripción de los Módulos invertibles sobre $V$ y, respectivamente, sobre
$X$, en términos de clases de divisores de Cartier. En efecto, los anillos
locales de $X$ son regulares ($X$ es liso sobre el esquema regular $S$) y, por
tanto, factoriales por el teorema de Auslander; así, todo divisor sobre
$X$\footnote{El texto impreso dice «sobre $S$». La factorialidad invocada es
la de los anillos locales de $X$.} es un divisor de Cartier. Todo divisor sobre
$V$ se prolonga a un divisor sobre $X$ tomando su clausura.
\end{proof}

\result{OBSERVACIÓN}{2.2} La demostración sigue siendo válida si se supone
solamente que $f$ es plano y que sus fibras $X_s$ son intersecciones completas
locales y lisas sobre $k(s)$ en codimensión $\leq2$. Esto se deduce del teorema
demostrado en [5]: un anillo local noetheriano de intersección completa que sea
regular en codimensión $\leq3$ es factorial («conjetura de Samuel»). Obsérvese
que el

\src{232}
resultado pasa a ser falso si se reemplaza «codimensión $\leq2$» por
«codimensión $\leq1$», es decir, por la hipótesis «normal», como muestra una
familia de cuádricas no singulares que degenera en un cono cuadrático.

(ii) El lema de Chow nos reduce al caso en que $X$ es proyectivo y, por tanto,
está inmerso en $\mathbf P_k^n$; podemos suponer que $X$ es conexo. Si
$\dim X=1$, entonces $X$ es liso sobre $k$ y se aplica (i). Si $\dim X\geq2$,
una variante de los «criterios de equivalencia» conocidos muestra que existe
un número finito de curvas lisas $Y_i$ contenidas en $X$ (obtenidas como
intersecciones de $X$ con subespacios lineales adecuados de $\mathbf P_k^n$),
tales que
\[
\operatorname{Pic}_{X/k}^\tau\longrightarrow
\prod_i\operatorname{Pic}_{Y_i/k}^\tau
\]
es un monomorfismo y, por tanto, induce un monomorfismo sobre las componentes
conexas. La componente conexa del segundo miembro es propia sobre $k$ por el
resultado anterior, y el homomorfismo de esquemas en grupos considerado es
necesariamente una inmersión cerrada. Se sigue que
$\operatorname{Pic}_{X/k}^\circ$ es también propio sobre $k$. Los delicados
criterios de equivalencia pueden evitarse utilizando la estructura de los
grupos algebraicos conmutativos sobre un cuerpo algebraicamente cerrado
(debida a Chevalley--Borel). Se reduce a demostrar que todo morfismo de la
recta afín privada del origen\footnote{El texto impreso dice «de la recta
afín». La fe de erratas vinculada, p.~303, prescribe «recta afín privada del
origen».} a $\operatorname{Pic}_{X/k}^\tau$ es constante. Esto equivale a decir
que todo Módulo invertible sobre $X[t,t^{-1}]$\footnote{El texto impreso dice
$X[t]$; la fe de erratas vinculada, p.~303, prescribe $X[t,t^{-1}]$.} procede
de un Módulo invertible sobre $X$. Este conocido resultado elemental ni siquiera
utiliza que $X$ sea propio sobre $k$ (la normalidad de $X$ reduce
inmediatamente al caso en que $X$ es regular).

\result{COROLARIO}{2.3} Sea $f:X\longrightarrow S$ un morfismo propio normal
(es decir, plano con fibras geométricamente normales), y supongamos que existe
$\operatorname{Pic}_{X/S}$. Entonces $\operatorname{Pic}_{X/S}^\circ$ es propio
sobre $S$ y, por tanto, cerrado en $\operatorname{Pic}_{X/S}$. Además,
$\operatorname{Pic}_{X/S}^\tau$ y $\operatorname{Pic}_{X/S}^\rho$ son también
cerrados, al igual que $\operatorname{Pic}_{X/S}^\sigma$ en «característica
igual».

Aplíquense (1.1) y (2.1), (ii).

\result{COROLARIO}{2.4} Sea $f:X\longrightarrow S$ un morfismo propio y liso
tal que exista $\operatorname{Pic}_{X/S}$ y que este sea una unión disjunta de
esquemas $P^{(i)}$ de tipo finito sobre $S$ (cf. [3], V, 4.1). Entonces cada
$P^{(i)}$ es propio sobre $S$.

Esto se deduce de (2.1), (i). Como se señaló en (2.2), el resultado puede
generalizarse imponiendo hipótesis menos restrictivas sobre las fibras de $f$,
pero pasa a ser falso si se supone meramente que $f$ es normal. En ese caso,
no sé si $\operatorname{Pic}_{X/S}^\tau$ es, no obstante, propio sobre $S$,
incluso suponiendo que sea de tipo finito sobre $S$.

\src{233}
\result{TEOREMA}{2.5} Sea $f:X\longrightarrow S$ un morfismo propio y plano tal
que exista $\operatorname{Pic}_{X/S}$, y para cada entero $n$ sea $\varphi_n$
el homomorfismo de potencia $n$-ésima de este preesquema en grupos. Entonces
$\varphi_n$ es étale en todo punto
$x\in\operatorname{Pic}_{X/S}$\footnote{El texto impreso dice «$x\in X$». La
aplicación $\varphi_n$ es el endomorfismo de potencia del preesquema de
Picard.} cuya característica residual sea prima con $n$.

Por la caracterización infinitesimal de los morfismos étales, esta afirmación
equivale a la siguiente:

\result{LEMA}{2.6} Supongamos que $S$ es el espectro de un anillo artiniano
local $A$ cuyo ideal maximal $\mathfrak m$ tiene potencia $(\nu+1)$-ésima nula,
y pongamos
\[
A_{\nu-1}=A/\mathfrak m^\nu,
\qquad X_{\nu-1}=X\otimes_A A_{\nu-1}.
\]
Sea $\mathcal L$ un Módulo invertible sobre $X$ y sea
$\mathcal L'_{\nu-1}$ un Módulo invertible sobre $X_{\nu-1}$ cuya potencia
tensorial $n$-ésima sea isomorfa a
\[
\mathcal L_{\nu-1}=\mathcal L\otimes_A A_{\nu-1}.
\]
Entonces existe un Módulo invertible $\mathcal L'$ sobre $X$ cuya potencia
tensorial $n$-ésima es isomorfa a $\mathcal L$, siempre que $n$ sea primo con
la característica residual de $k=k(A)$.

Demostremos el lema. Pongamos
\[
V=\mathfrak m^\nu=\mathfrak m^\nu/\mathfrak m^{\nu+1};
\]
este es un espacio vectorial sobre $k=k(A)$. Intentemos primero prolongar
$\mathcal L'_{\nu-1}$ a un Módulo invertible arbitrario $\mathcal L'$ sobre $X$.
La obstrucción pertenece a
\[
H^2(X_0,\mathcal O_{X_0})\otimes_k V.
\]
La hipótesis sobre $\mathcal L'_{\nu-1}$ y el hecho de que
$\mathcal L_{\nu-1}$ pueda prolongarse muestran que $n$ veces esta obstrucción
es cero; la propia obstrucción es, por tanto, nula porque $n$ es primo con la
característica. La ambigüedad en la prolongación elegida pertenece a
\[
H^1(X_0,\mathcal O_{X_0})\otimes_k V,
\]
y la diferencia $\xi$ entre $\mathcal L'^{\otimes n}$ y $\mathcal L$ pertenece
al mismo Módulo. Corregir $\mathcal L'$ para anular esta diferencia equivale a
encontrar un $\eta$ en ese Módulo tal que $n\eta=\xi$. Esto vuelve a ser
posible porque $n$ es primo con la característica.

\result{COROLARIO}{} Bajo las hipótesis de (2.5), supongamos además que los
esquemas de Picard de las fibras $X_s$ no contienen ninguna componente aditiva
(por ejemplo, que los $X_s$ son geométricamente normales; cf. (2.1), (ii)).
Entonces $\operatorname{Pic}_{X/S}\longrightarrow S$ es universalmente abierto
en los puntos de $\operatorname{Pic}_{X/S}^\circ$. Si
$\operatorname{Pic}_{X/S}^\circ$ es cerrado (por ejemplo, si los $X_s$ son
geométricamente normales; cf. (2.3)), entonces
$\operatorname{Pic}_{X/S}^\sigma$ es universalmente abierto sobre $S$. Por
último, en característica igual,
$\operatorname{Pic}_{X/S}^\rho\longrightarrow S$ es universalmente abierto.

Aplíquense (1.5) y (1.1).

\src{234}
\result{COROLARIO}{2.7} Sea $f:X\longrightarrow S$\footnote{El texto impreso
dice «$f:X\longrightarrow Y$», aunque tanto el preesquema de Picard como la
función siguiente tienen base $S$.} un morfismo propio y plano tal que exista
$\operatorname{Pic}_{X/S}$. Entonces la función
\[
s\longmapsto\dim\operatorname{Pic}_{X_s/k(s)}
\]
sobre $S$ es semicontinua superiormente (es decir, puede saltar hacia arriba,
pero no hacia abajo), e incluso continua (es decir, localmente constante) si
los $\operatorname{Pic}_{X_s/k(s)}$ no contienen ninguna componente aditiva.

La primera afirmación es trivial, o casi, para todo preesquema en grupos
localmente de tipo finito sobre una base localmente noetheriana, pues basta
mirar a lo largo de la sección neutra. La segunda afirmación se deduce de
(2.5).

\result{OBSERVACIÓN}{2.8} Sean $s,s'\in S$, con $s$ una especialización de
$s'$. Entonces (2.7) equivale a una desigualdad (respectivamente, una igualdad)
entre las dimensiones de los esquemas de Picard de $X_{s'}$ y de su
«especialización» $X_s$. Incluso antes de que se construyeran los esquemas de
Picard, Serre había observado que, para un morfismo liso
$f:X\longrightarrow S$, la invariancia de la dimensión de los Picard de los
$X_s$ es una consecuencia formal de la teoría de la especialización del grupo
fundamental ([4], X) y de las relaciones clásicas de Kummer entre los puntos de
orden finito de la variedad de Picard y el grupo fundamental abelianizado
([4], XI). En general, sean $\alpha,\mu,\lambda$ las dimensiones de las partes
abeliana, multiplicativa y aditiva de
$\operatorname{Pic}_{X_s/k(s)}^\circ$, y definamos de modo análogo
$\alpha',\mu',\lambda'$. Las relaciones conocidas se expresan mediante las
desigualdades siguientes:
\[
\tag{*}\alpha+\mu+\lambda\geq\alpha'+\mu'+\lambda'
\]
(válida siempre que exista $\operatorname{Pic}_{X/S}$ y, por tanto,
probablemente en todos los casos). Cuando $\lambda=\lambda'=0$, esta
desigualdad se convierte en una igualdad, válida bajo las mismas hipótesis de
existencia:
\[
\alpha+\mu=\alpha'+\mu';
\]
por otra parte,
\[
\tag{**}2\alpha+\mu\leq2\alpha'+\mu'
\]

\src{235}
por el argumento de Serre si los $X_s$ son separables (sin suponer siquiera la
existencia de $\operatorname{Pic}_{X/S}^\tau$), o si los
$ {}_n\operatorname{Pic}_{X/S}$ (los núcleos de $\varphi_n$ en el funtor de
Picard) son separados sobre el abierto de $S$ donde $n$ es invertible (pues son étales sobre ese abierto por (2.5)). Se
inclina uno a conjeturar que (*) es siempre una igualdad, o al menos cuando los
$X_s$ son separables, y que
\[
\tag{***}\alpha\leq\alpha',\qquad\lambda\geq\lambda'
\]
se verifica siempre que exista un preesquema en grupos, localmente de tipo
finito sobre un $S$ localmente noetheriano, cuyas fibras tengan dimensión
constante (véase (1.3) para un resultado positivo en esta dirección).

\result{OBSERVACIÓN}{2.9} En todos los casos conocidos,
$\operatorname{Pic}_{X/S}^\tau$ es universalmente abierto sobre $S$, pero
probablemente no conviene dejarse llevar por ilusiones excesivas, ni siquiera
cuando $f:X\longrightarrow S$ es liso. En todo caso, Mumford ha construido un
ejemplo (ciertamente con $S$ no reducido, de hecho el espectro de un anillo
artiniano) en el que $\operatorname{Pic}_{X/S}^\tau$ no es plano sobre $S$,
haciendo variar infinitesimalmente la superficie de Igusa. El punto considerado
por Mumford pertenece a $\operatorname{Pic}_{X/S}^\rho$, y sigue siendo posible
que, para un $f:X\longrightarrow S$ liso, $\operatorname{Pic}_{X/S}^\tau$ sea
plano sobre $S$ en los puntos de $\operatorname{Pic}_{X/S}^\sigma$. El ponente
duda, sin embargo, que esto sea siempre así, incluso si se restringe a los
puntos de $\operatorname{Pic}_{X/S}^\circ$ y se supone que $S$ es el espectro
de un anillo de valoración discreta. La cuestión está también relacionada con
el estudio de los puntos fijos de un esquema abeliano bajo un grupo finito de
automorfismos, situación para la que parecen faltar ejemplos.

Parecería que, incluso restringiéndose a un morfismo proyectivo liso
$f:X\longrightarrow S$, los resultados de regularidad local para
$\operatorname{Pic}_{X/S}$ enunciados en este parágrafo, junto con las
conjeturas de (2.8), agotan casi todo lo que puede decirse sin hipótesis más
específicas sobre la naturaleza de las fibras de $f$. Recuérdese, sin embargo,
que si las fibras geométricas de $\operatorname{Pic}_{X/S}$ son reducidas y no
tienen componente aditiva, entonces (1.8) y (2.5) implican que
$\operatorname{Pic}_{X/S}$ es liso sobre $S$ en los puntos de
$\operatorname{Pic}_{X/S}^\sigma$ cuando $S$ es reducido. El mismo resultado se
verifica, sin ninguna hipótesis sobre $S$, si $f:X\longrightarrow S$ es normal,
como veremos en (3.5). Consignemos a este respecto:

\result{PROPOSICIÓN}{2.10}

\noindent(i) Si $\operatorname{Pic}_{X/S}$ es liso (respectivamente, plano)
sobre $S$ en los puntos de la sección neutra, también lo es en todo punto de
$\operatorname{Pic}_{X/S}^\sigma$ y en los puntos de toda sección de
$\operatorname{Pic}_{X/S}$ sobre $S$.

\src{236}
\noindent(ii) Sea $s\in S$ tal que
\[
H^2(X_s,\mathcal O_{X_s})=0.
\]
Entonces existe un entorno abierto $U$ de $s$ tal que
$\operatorname{Pic}_{X/S}|U$ es liso sobre $U$.

\noindent(iii) Si $X$ es un esquema propio sobre un cuerpo, entonces
\[
\dim\operatorname{Pic}_{X/k}\leq\dim H^1(X,\mathcal O_X),
\]
con igualdad si y solo si $\operatorname{Pic}_{X/k}$ es liso sobre $k$; este
es siempre el caso cuando $k$ tiene característica cero.

La parte (i) se deduce de (2.5) y (1.12), y (ii) se deduce del criterio
infinitesimal de lisura y de un cálculo conocido de obstrucciones, pues el
«teorema de semicontinuidad» muestra que la hipótesis en $s$ sigue siendo
válida en los puntos próximos. Por último, (iii) se deduce de que
$H^1(X,\mathcal O_X)$ es isomorfo al espacio tangente de Zariski de
$\operatorname{Pic}_{X/k}$ en el elemento neutro. La última afirmación es un
caso particular de un teorema de Cartier que dice que un «grupo formal» en
característica cero es formalmente liso sobre $k$.

\fsection{3}{El subesquema abeliano canónico de
$\operatorname{Pic}_{X/S}$ y el esquema de Albanese}

\result{PROPOSICIÓN}{3.1} Sea $k$ un cuerpo y sea $G$ un esquema conmutativo en
grupos de tipo finito sobre $k$, sin componente aditiva. Entonces
$G_{\mathrm{red}}^\circ$ es separable sobre $k$ y, por consiguiente, es un
subesquema en grupos liso sobre $k$.

Esto es trivial cuando $k$ es perfecto, en particular para $G_{\bar k}$, donde
$\bar k$ es una clausura algebraica de $k$. Basta, por tanto, ver que
$(G_{\bar k}^\circ)_{\mathrm{red}}$ desciende de un subesquema de $G$. De la
hipótesis sobre $G_{\bar k}$ (ausencia de componente aditiva)\footnote{El texto
impreso contiene una construcción mal formada equivalente a «La hipótesis ...
resulta». La fe de erratas vinculada, p.~303, inserta la preposición necesaria
para «De la hipótesis ... resulta».} se deduce fácilmente que existe un entero
$m$ tal que $(G_{\bar k}^\circ)_{\mathrm{red}}$ es la imagen esquemática del
homomorfismo de potencia $m$-ésima de $G_{\bar k}$. Como este homomorfismo
procede del homomorfismo análogo de $G$, la imagen esquemática de este último
da el subesquema buscado.

\result{COROLARIO}{3.2} Sea $X$ un esquema geométricamente normal propio sobre $k$ tal que
exista $\operatorname{Pic}_{X/k}$. Entonces existe un subesquema abeliano $A$
de $\operatorname{Pic}_{X/k}$ cuyo conjunto subyacente es
$\operatorname{Pic}_{X/k}^\circ$.

En efecto, por (2.1), (ii), $\operatorname{Pic}_{X/k}^\circ$ es propio sobre
$k$ y, por tanto, satisface la condición de (3.1).

\src{237}
Así, el resultado anterior muestra que, en ciertos casos, la «variedad de
Picard» clásica (que en la presente teoría se denotaría
$(\operatorname{Pic}_{X/k})_{\mathrm{red}}^\circ$) «está definida sobre $k$»,
sin suponer que $k$ sea perfecto.

Sea ahora $f:X\longrightarrow S$ un esquema relativo propio y plano, con
\[
\mathcal O_S\xrightarrow{\sim}f_*(\mathcal O_X)
\]
para simplificar, tal que exista $\operatorname{Pic}_{X/S}$ y
$\operatorname{Pic}_{X/S}^\circ$ sea propio sobre $S$. Para (3.3), (ii),
suponemos además que existe un abierto de $\operatorname{Pic}_{X/S}$ que
contiene $\operatorname{Pic}_{X/S}^\circ$ y es cuasiproyectivo sobre $S$. Como
se vio anteriormente, esta condición se verifica cuando $f$ es proyectivo con
fibras geométricamente separables e irreducibles. Recuérdese que un esquema
abeliano sobre $S$ es un esquema en grupos sobre $S$ propio y liso sobre $S$,
con fibras geométricamente conexas. Preguntamos si existe un subesquema en
grupos $A$ de $\operatorname{Pic}_{X/S}$ que sea un esquema abeliano y cuyo
conjunto subyacente sea $\operatorname{Pic}_{X/S}^\circ$. Acabamos de ver que
siempre existe uno cuando $S$ es el espectro de un cuerpo. He aquí lo que puede
decirse en el caso general considerado:

\result{TEOREMA}{3.3} Bajo las hipótesis anteriores:

\noindent(i) Si existe un subesquema abeliano de
$\operatorname{Pic}_{X/S}$ cuyo conjunto subyacente sea
$\operatorname{Pic}_{X/S}^\circ$, es único. Por consiguiente, su formación es
compatible con el cambio de base.

\noindent(ii) Para que exista tal subesquema abeliano, basta que exista después
de todo cambio de base $S'\longrightarrow S$ con $S'$ artiniano local. Si $S$
es el espectro de un anillo local, basta comprobar los
$S'=\operatorname{Spec}(A_n)$, donde $A_n=A/\mathfrak m^{n+1}$. Si $S$ es
reducido, basta asimismo comprobar los cambios de base
$S'\longrightarrow S$ en los que $S'$ es el espectro de un anillo de
valoración discreta (completo y con cuerpo residual algebraicamente cerrado,
si se desea).

\noindent(iii) Supongamos que existe $A$, y sea
$B=\operatorname{Alb}^0(X/S)$ el esquema abeliano dual (es decir,
$B=\operatorname{Pic}_{A/S}^\circ$ [7]). Se puede construir entonces,
canónicamente, un espacio homogéneo principal
$P=\operatorname{Alb}^1(X/S)$ bajo $B$, junto con un $S$-morfismo
$X\longrightarrow P$ universal para los $S$-morfismos de $X$ a esquemas
para-abelianos (es decir, espacios homogéneos principales bajo esquemas
abelianos). La formación de $\operatorname{Alb}^0(X/S)$,
$\operatorname{Alb}^1(X/S)$ y
$X\longrightarrow\operatorname{Alb}^1(X/S)$ conmuta con el cambio de base.

Esbocemos la demostración.

La parte (i) es una propiedad general de rigidez de los subesquemas abelianos
de los esquemas conmutativos en grupos: si dos de esos subesquemas coinciden
como conjuntos en un punto $s\in S$, coinciden sobre toda la componente conexa

\src{238}
de $s$ [7]. (Este resultado generaliza un teorema clásico de Chow.)

(ii) Los esquemas de Hilbert muestran que el funtor que asocia a cada $S'$ sobre
$S$ el conjunto (con uno o cero elementos) de los subesquemas abelianos
canónicos de $(\operatorname{Pic}_{X/S})\times_S S'$ está representado por un
esquema $T$ de tipo finito sobre $S$. Por (i), $T\longrightarrow S$ es un
monomorfismo y, por (3.2), es sobreyectivo. La existencia de un subesquema
abeliano canónico de $\operatorname{Pic}_{X/S}$ significa que $T$ posee una
sección sobre $S$, o equivalentemente que $T\longrightarrow S$ es un
isomorfismo. Esto equivale también a que $T\longrightarrow S$ sea étale o,
cuando $S$ es reducido, a que $T\longrightarrow S$ sea propio. Esto demuestra
(ii).

(iii) Directamente de la definición de $\operatorname{Pic}_{X/S}$, para todo
esquema abeliano $C$ sobre $S$, dar un $S$-morfismo de $X$ a un espacio
homogéneo principal bajo $C$ equivale a dar un homomorfismo de grupos
$C'\longrightarrow\operatorname{Pic}_{X/S}$, donde $C'$ es el esquema abeliano
dual de $C$. Si existe el subesquema abeliano canónico $A$ de
$\operatorname{Pic}_{X/S}$, estos homomorfismos se factorizan necesariamente a
través de $A$ (como puede verse también utilizando los puntos de orden finito).
Esto demuestra (iii).

\result{OBSERVACIONES}{3.4} Denotaremos por
$\operatorname{Pic}_{X/S}^{\circ\circ}$ el subesquema abeliano canónico de
$\operatorname{Pic}_{X/S}$, cuando exista. Por desgracia, no siempre existe,
como se ve haciendo variar infinitesimalmente la superficie de Igusa (mediante
la deformación modular de primer orden). Es posible, no obstante, que
$\operatorname{Pic}_{X/S}^{\circ\circ}$ exista al menos cuando $S$ sea reducido
o, equivalentemente por (ii), cuando $S$ sea el espectro de un anillo de
valoración discreta. Sean $X_0,X_1$ las fibras especial y genérica de $X$, y
pongamos
\[
A_1=\operatorname{Pic}_{X_1/K}^\circ,
\]
donde $K$ es el cuerpo de fracciones del anillo de valoración $V$. Por el
teorema de Koizumi, existe un esquema abeliano $A$ sobre $S$, esencialmente
único, cuya fibra genérica es $A_1$. Como en (2.1), (i), se ve fácilmente
(suponiendo ahora que $X$ es liso sobre $S$) que el morfismo identidad de $A_1$
se prolonga a un morfismo
\[
A\longrightarrow\operatorname{Pic}_{X/S}.
\]
De ahí se obtiene un homomorfismo
\[
\tag{*}A_0\longrightarrow\operatorname{Pic}_{X_0/k}^{\circ\circ},
\]
que se demuestra fácilmente que es sobreyectivo y tiene por núcleo un
$p$-grupo finito, donde $p$ es el exponente característico del cuerpo residual
$k$

\src{239}
(de nuevo utilizando los puntos de orden finito). Sentado esto, las condiciones
siguientes sobre $X/S$ son equivalentes:

\medskip
\noindent a. El homomorfismo anterior (*) es un isomorfismo (lo que Shimura
expresaría diciendo que la formación de la «variedad de Picard» es «compatible
con la especialización»).

\noindent b. Existe $\operatorname{Pic}_{X/S}^{\circ\circ}$ (en cuyo caso no es
otro que $A$).

\noindent c. (Para referencia) Existen los
$\operatorname{Pic}_{X_n/S}^{\circ\circ}$.

Por la observación hecha sobre el núcleo de (*), la condición (a) se verifica
cuando la característica residual es cero, pero este resultado se generalizará
considerablemente en (3.5).

Por supuesto, si $\operatorname{Pic}_{X/S}$ es liso sobre $S$ en los puntos de
$\operatorname{Pic}_{X/S}^\circ$, este último es abierto en
$\operatorname{Pic}_{X/S}$ (cf. (1.7)). Provisto de la estructura inducida, es
por tanto un subesquema abeliano de $\operatorname{Pic}_{X/S}$, luego igual a
$\operatorname{Pic}_{X/S}^{\circ\circ}$, que existe en este caso. Pero se tiene
algo mucho más fuerte:

\result{TEOREMA}{3.5} Bajo las hipótesis de (3.3), sea $s\in S$ tal que
$\operatorname{Pic}_{X_s/k(s)}$ sea liso sobre $k(s)$ (o, equivalentemente,
\[
\dim\operatorname{Pic}_{X_s/k(s)}=\dim H^1(X_s,\mathcal O_{X_s})).
\]
Entonces existe un entorno abierto $U$ de $s$ tal que
$\operatorname{Pic}_{X/S}$ es liso sobre $S$ en los puntos de
$\operatorname{Pic}_{X/S}^\circ|U$, que es, por consiguiente, un subesquema
abeliano abierto de $\operatorname{Pic}_{X/S}|U$. A fortiori, existe
$\operatorname{Pic}_{X|U/U}^{\circ\circ}$.

Indiquemos el principio de la demostración. La discusión anterior nos reduce
al caso en que $S$ es el espectro de un anillo artiniano local $A$, y se razona
por inducción sobre el orden infinitesimal de $A$. Podemos suponer, por tanto,
que $\operatorname{Pic}_{X_n/A_n}^\circ$ es liso sobre $A_n$ y basta demostrar
que $\operatorname{Pic}_{X_{n+1}/A_{n+1}}^\circ$ es liso sobre $A_{n+1}$. Para
ello basta construir un esquema abeliano $P_{n+1}$ sobre $A_{n+1}$ que prolongue
\[
P_n=\operatorname{Pic}_{X_n/A_n}^\circ,
\]
junto con un Módulo invertible $\mathcal L_{n+1}$ sobre
$X_{n+1}\times_{A_{n+1}}P_{n+1}$ que prolongue el Módulo invertible
$\mathcal L_n$ sobre $X_n\times_{A_n}P_n$ que interviene en la definición del
esquema de Picard $\operatorname{Pic}_{X_n/A_n}$ como solución de un problema
universal. (N.B. --- Podemos suponer que $X$ posee una sección sobre $S$~...)
La construcción utiliza el resultado clave siguiente: todo esquema abeliano
definido sobre un cociente de un anillo artiniano local puede prolongarse (en
otros términos, el «esquema formal de Módulos» absoluto ([3], II) de un esquema
abeliano sobre un cuerpo algebraicamente cerrado es liso sobre el anillo de

\src{240}
vectores de Witt de $k$). Este resultado se deduce, a su vez, de las propiedades
formales generales de la obstrucción al levantamiento y de las operaciones de
grupo. Admitiéndolo, se prolonga primero $P_n$ arbitrariamente a $P_{n+1}$. Se
encuentra entonces una obstrucción al levantamiento de $\mathcal L_n$, que
pertenece a
\[
H^2(X_0\times P_0,\mathcal O_{X_0\times P_0})\otimes_k V
\qquad(\text{donde }V=\mathfrak m^{n+1}/\mathfrak m^{n+2}),
\]
y, de hecho, al subespacio
\[
H^1(X_0,\mathcal O_{X_0})\otimes
H^1(P_0,\mathcal O_{P_0})\otimes_k V
\]
(porque la restricción de $\mathcal L_n$ a cada uno de los dos factores $X_n$
y $P_n$ es trivial). Este último espacio es precisamente
\[
H^1(P_0,\mathcal T_{P_0/k})\otimes V,
\]
donde $\mathcal T_{P_0/k}$ es el haz tangente de $P_0/k$, y es también, por
tanto, el espacio que mide la indeterminación en el levantamiento de $P_n$ a
$P_{n+1}$ ([3], II). Por consiguiente, este levantamiento puede corregirse (de
una y solo una manera, como debe ser) para anular la obstrucción al
levantamiento de $\mathcal L_n$.

\result{COROLARIO}{3.6} Bajo las hipótesis de (3.5),
$R^1f_*(\mathcal O_X)$ es un Módulo localmente libre sobre $S$ en un entorno de
$s$, y su formación conmuta con el cambio de base.

En efecto, este Módulo es el Módulo tangente de $\operatorname{Pic}_{X/S}$ a lo largo
de la sección neutra.

\result{OBSERVACIÓN}{3.7} El mismo argumento que en (3.5) muestra que, si $S'$
es un subesquema de $S$ definido por un Ideal coherente nilpotente y se supone
solamente que existe $\operatorname{Pic}_{X'/S'}^\circ$ y que tiene fibras
lisas, entonces existe necesariamente $\operatorname{Pic}_{X/S}^\circ$ y es un
esquema abeliano sobre $S$. Esto permite construir esquemas de Picard en
algunos casos pese a la ausencia de una hipótesis de proyectividad; por
ejemplo, siempre existe el esquema abeliano dual de un esquema abeliano sobre
un anillo artiniano. Por otra parte, aplicando (3.6) cuando $X$ es un esquema
abeliano sobre $S$ y utilizando la estructura conocida de álgebra exterior de
$H^*(X_s,\mathcal O_{X_s})$ sobre $H^1(X_s,\mathcal O_{X_s})$
(Rosenlicht--Serre), se obtiene que, para todo $i$, $R^if_*(\mathcal O_X)$ es
localmente libre y, más precisamente, isomorfo a la potencia exterior $i$-ésima
de $R^1f_*(\mathcal O_X)$.

\result{OBSERVACIÓN}{3.8} Para un morfismo proyectivo liso
$f:X\longrightarrow S$, con $S$ reducido y todas las características
residuales nulas, (3.6) era conocido por métodos trascendentes como consecuencia
de la teoría de Hodge. De hecho, todos los $R^pf_*(\Omega^q_{X/S})$ son
entonces localmente libres. Por otra parte, existen contraejemplos

\src{241}
en característica mixta a la libertad local de $R^1f_*(\mathcal O_X)$,
obtenidos a partir de variedades de Serre ([8], p.~50). No parece haber
contraejemplos en característica igual.

\result{COROLARIO}{3.8} Bajo las hipótesis de (3.5),
$\operatorname{Pic}_{X/S}|U$ es liso sobre $U$ en todo punto de
$\operatorname{Pic}_{X/S}^\sigma|U$.

Aplíquese (2.10), (i). En particular, teniendo en cuenta (2.10), (iii):

\result{COROLARIO}{3.9} Supongamos que todas las características residuales de
$S$ son nulas. Entonces $\operatorname{Pic}_{X/S}^\tau$ es liso sobre $S$.

\result{OBSERVACIÓN}{3.10} Si $\operatorname{Pic}_{X/S}^\tau$ es además propio sobre
$S$, se sigue, por ejemplo, que el subgrupo de torsión de los grupos de
Néron--Severi de las fibras geométricas de $f$ es constante sobre cada
componente conexa de $S$ (como también resulta claro por métodos trascendentes
cuando $f$ es liso y proyectivo). Una aplicación directa de (2.5) muestra, más
generalmente, que si $\operatorname{Pic}_{X/S}^\tau$ es propio sobre $S$ (por
ejemplo, si $f:X\longrightarrow S$ es liso y proyectivo) y $q$ es un primo
distinto de todas las características residuales de $S$, entonces la
componente $q$-primaria de los subgrupos de torsión de los grupos de
Néron--Severi de las fibras geométricas de $X/S$ es constante sobre cada
componente conexa de $S$. En característica $p>0$, esto deja de ser cierto para
la componente $p$-primaria de la torsión. Sigue siendo posible, sin embargo,
que el rango sobre $k(s)$ de
\[
\operatorname{Pic}_{X_s/k(s)}^\tau/
\operatorname{Pic}_{X_s/k(s)}^{\circ\circ}=T_{X_s/k(s)}
\]
sea localmente constante. Cuando $S$ es reducido, se puede demostrar que esto
equivale a la existencia de $\operatorname{Pic}_{X/S}^{\circ\circ}$ junto con
la planitud de $\operatorname{Pic}_{X/S}^\tau$ sobre $S$, y basta comprobar la
cuestión cuando $S$ es el espectro de un anillo de valoración discreta. Esto es
lo que he comprobado en los pocos ejemplos que he examinado; pero como el
enunciado correspondiente con $S$ artiniano es falso (cf. Observaciones (2.9)
y (3.4)), probablemente no convenga entregarse a un optimismo excesivo.

\fsection{4}{El teorema de finitud para el esquema de Picard}

Sea $f:X\longrightarrow S$ un morfismo proyectivo y plano tal que exista
$\operatorname{Pic}_{X/S}$. La consideración de los «polinomios de Hilbert»
$Q\in\mathbf Q[t]$ descompone $\operatorname{Pic}_{X/S}$ en una unión disjunta
de abiertos $\operatorname{Pic}_{X/S}^Q$. Sin hipótesis adicionales, por
ejemplo que aseguren que $\operatorname{Pic}_{X/S}$ sea separado sobre $S$,
estos abiertos no son necesariamente de tipo finito sobre $S$; se obtiene un
contraejemplo cuando $X$ es «una cónica que degenera en dos rectas». Es

\src{242}
posible, sin embargo, que sean de tipo finito cuando $f$ sea separable y tenga
fibras geométricamente irreducibles. La cuestión está relacionada con saber si
$\operatorname{Pic}_{X/S}^\tau$ es de tipo finito sobre $S$, lo que podría ser
cierto sin ninguna hipótesis sobre las fibras de $X/S$. Cuando
$f:X\longrightarrow S$ es liso, obsérvese que
$\operatorname{Pic}_{X/S}^\tau$ está contenido en uno de los
$\operatorname{Pic}_{X/S}^Q$ (pues sobre una variedad proyectiva no singular la
equivalencia por torsión es más fina que la equivalencia numérica para los
divisores y, de hecho, por Matsusaka, coincide con ella), y es, por tanto, de
tipo finito sobre $S$ si los $\operatorname{Pic}_{X/S}^Q$ lo son. Las
cuestiones de finitud que acabamos de plantear conservan un sentido evidente
incluso sin suponer la existencia de $\operatorname{Pic}_{X/S}$: dicen que
ciertas familias de Módulos invertibles son «acotadas» en el sentido de [3], IV.
Para todo cuerpo algebraicamente cerrado $k$, considérense los subesquemas
íntegros de $\mathbf P_k^r$ de dimensión $n$ y grado $d$, y los Módulos
invertibles sobre esos preesquemas con polinomio de Hilbert $Q$, donde
$r,n,d,Q$ están fijados. El problema consiste en demostrar que la familia de
estos Módulos (considerados como Módulos coherentes sobre los $\mathbf P_k^r$) es
acotada, es decir, que puede parametrizarse mediante un esquema de tipo finito
sobre $\mathbf Z$.

Utilizando el método de Matsusaka [6], una demostración bastante técnica (que
emplea los «criterios de equivalencia» en una forma adecuada) da una respuesta
afirmativa cuando uno se restringe a las subvariedades no singulares de
$\mathbf P_k^r$. Más precisamente:

\result{TEOREMA}{4.1} Sea $f:X\longrightarrow S$\footnote{El texto impreso
dice «Sea: $X\longrightarrow S$»; el morfismo utilizado más adelante es $f$.}
un morfismo proyectivo liso con fibras geométricamente conexas, siendo $S$
noetheriano. Sea $\mathcal O_X(1)$ un Módulo relativamente muy amplio sobre $X$,
y sea $E$ un subconjunto de $\operatorname{Pic}_{X/S}$ correspondiente a una
familia $(\mathcal L_i)$ de Módulos invertibles sobre las fibras geométricas
$\overline X_{s_i}$ de $X/S$. Sea $D_i$ un divisor (no necesariamente
efectivo) sobre $\overline X_{s_i}$ que defina $\mathcal L_i$, sea
\[
a_n^{(i)}t^n+\cdots+a_0^{(i)}
\]
el polinomio de Hilbert de $\mathcal L_i$, y sea $\xi=\xi_i$ un divisor que
defina $\mathcal O_X(1)$, es decir, una sección hiperplana. Las condiciones
siguientes son equivalentes:

\noindent a. $E$\footnote{El texto impreso dice «$Q$». El subconjunto
introducido en el enunciado, cuya cuasicompacidad expresa la acotación de la
familia, es $E$.} es cuasicompacto, es decir, está contenido en un abierto de
tipo finito sobre $S$; equivalentemente, la familia $(\mathcal L_i)$ es
acotada.

\noindent b. $E$ está contenido en la unión de un número finito de los
conjuntos $\operatorname{Pic}_{X/S}^Q$, $Q\in\mathbf Q[t]$; equivalentemente,
solo aparece un número finito de polinomios de Hilbert entre los
$\mathcal L_i$.

\noindent b'. (Si todas las fibras de $X/S$ tienen la misma dimensión $n$.)
Los dos coeficientes $a_{n-1}^{(i)}$ y $a_{n-2}^{(i)}$ de los polinomios de
Hilbert de los $\mathcal L_i$ recorren conjuntos finitos.

\src{243}
\noindent b''. (Si todas las fibras de $X/S$ tienen la misma dimensión $n$.)
Los enteros
\[
\xi^{n-1}D_i\qquad\text{y}\qquad\xi^{n-2}D_i^2
\]
(los grados proyectivos de $D_i$ y $D_i^2$) permanecen acotados superiormente.

\result{COROLARIO}{4.2} Sea $f:X\longrightarrow S$ un morfismo proyectivo liso
con fibras geométricamente conexas. Entonces los esquemas
$\operatorname{Pic}_{X/S}^Q$ y $\operatorname{Pic}_{X/S}^\tau$
($Q\in\mathbf Q[t]$) son proyectivos sobre $S$.

Son de tipo finito sobre $S$ por (4.1), por lo que se aplica (2.4).

\section*{BIBLIOGRAFÍA}
\addcontentsline{toc}{section}{Bibliografía}

\begin{description}
\item[{[1]}] CHEVALLEY (Claude). -- Sur la théorie de Picard,
\emph{Amer. J. of Math.}, vol. 82, 1960, pp. 435--490.

\item[{[2]}] GROTHENDIECK (A.) y DIEUDONNÉ (J.). -- Éléments de géométrie
algébrique, capítulo I y siguientes. -- París, Presses universitaires de
France, 1960--1961 (Institut des Hautes Études Scientifiques, 4, 8, 11,
\ldots).

\item[{[3]}] GROTHENDIECK (Alexander). -- Technique de descente et théorèmes
d'existence en géométrie algébrique, I--V, \emph{Séminaire Bourbaki}, t. 12,
1959/60, n.º 190, 29 pp., y n.º 195, 22 pp.; t. 13, 1960/61, n.º 212,
20 pp., y n.º 221, 28 pp.; t. 14, 1961/62, n.º 232, 19 pp.

\item[{[4]}] GROTHENDIECK (Alexander). -- Séminaire de géométrie algébrique,
1960/61. -- París, Institut des Hautes Études Scientifiques, 1961.

\item[{[5]}] GROTHENDIECK (Alexander). -- Séminaire de géométrie algébrique,
1961/62, redactado por un grupo de oyentes. -- París, Institut des Hautes Études
Scientifiques (en preparación).

\item[{[6]}] MATSUSAKA (T.). -- The criteria for algebraic equivalence and the
torsion group, \emph{Amer. J. Math.}, vol. 79, 1957, pp. 53--66.

\item[{[7]}] MUMFORD (D.) y TATE (J.). -- Séminaire de géométrie algébrique,
Harvard University, semestre de primavera de 1962 (en preparación).

\item[{[8]}] SERRE (Jean-Pierre). -- Sur la topologie des variétés algébriques
en caractéristique $p$, \emph{Symposium internacional de topología algebraica}
[1956, México], pp. 24--53. -- México, Universidad Nacional Autónoma, 1958.
\end{description}

\def\FgaUnitKey{e236}\clearpage
\chapter*{Fe de erratas y complemento a la exposición n.º 236}
\addcontentsline{toc}{chapter}{Fe de erratas y complemento a la exposición n.º 236}
\src{303}

\noindent\textbf{8. TDTE VI: Los esquemas de Picard. Propiedades generales.}

\begin{center}
[Seminario Bourbaki, t. 14, 1961/62, n.º 236, 23 pp.]
\end{center}

\noindent\textbf{Página 12, línea 16.} --- En lugar de «recta afín», léase
«recta afín privada del origen».

\noindent\textbf{Línea 18.} --- En lugar de «$X[t]$», léase
«$X[t,t^{-1}]$».

\noindent\textbf{Página 16, proposición 3.1, 6.ª línea.} --- En lugar de «Ahora
la hipótesis ...», léase «Ahora bien, de la hipótesis ...».

\noindent\textbf{Página 21, § 4: El teorema de finitud para el esquema de
Picard.} --- Las cuestiones de finitud del género de las planteadas en este
parágrafo han quedado resueltas casi por completo desde que se redactó la
presente exposición. Indiquemos los principales hechos conocidos actualmente
en esta dirección. Para simplificar los enunciados, suponemos implícitamente
que existen todos los preesquemas de Picard que intervienen en ellos, aunque
una modificación evidente de estos enunciados permite prescindir de toda
hipótesis de existencia. En lo sucesivo, $S$ designa un esquema noetheriano y
$X,Y$ esquemas propios sobre $S$.

\medskip
\noindent(i) Sea $f:X\longrightarrow Y$ un morfismo sobreyectivo. Entonces
\[
f^*: \operatorname{Pic}_{Y/S}\longrightarrow\operatorname{Pic}_{X/S}
\]
es un morfismo de tipo finito.

\src{304}
\noindent(ii) Supongamos que $Y$ es proyectivo sobre $S$ y está provisto de un
Módulo invertible relativamente amplio sobre $S$, y sea $X$ el esquema de ceros
de una sección cualquiera de este Módulo, con $f:X\longrightarrow Y$ la inmersión
canónica. Supongamos, por último, que las componentes irreducibles de las
fibras de $Y/S$ tienen dimensión $\geq3$. Entonces
\[
f^*: \operatorname{Pic}_{Y/S}\longrightarrow\operatorname{Pic}_{X/S}
\]
es de tipo finito.

\noindent(iii) Supongamos que $X$ es proyectivo sobre $S$ y que todas sus
fibras geométricas son íntegras y de dimensión $n$. Sea $\mathcal O_X(1)$ un
Módulo invertible amplio sobre $X$ que permita definir polinomios de Hilbert. Sea
$M$ una parte de $\operatorname{Pic}_{X/S}$; entonces $M$ es cuasicompacta si
y solo si, en los polinomios de Hilbert
\[
a_0x^n+a_1x^{n-1}+a_2x^{n-2}+\cdots
\]
de los elementos de $M$, los coeficientes $a_1$ y $a_2$ permanecen acotados, omitiendo todo coeficiente cuyo monomio tendría grado negativo; así, solo interviene $a_1$ si $n=1$, y ninguno de los dos si $n=0$.

\noindent(iv) Para todo entero $n\neq0$, el homomorfismo de potencia $n$-ésima
del preesquema en grupos $\operatorname{Pic}_{X/S}$ es un morfismo de tipo
finito.

Obsérvese que (i) y (ii) significan también que, bajo las hipótesis
consideradas respectivamente en esos enunciados, una parte $M$ de
$\operatorname{Pic}_{Y/S}$ es cuasicompacta si y solo si lo es su imagen en
$\operatorname{Pic}_{X/S}$. Se deduce que un Módulo invertible $\mathcal L$ sobre
$Y$ es $\tau$-equivalente a cero si y solo si lo es su imagen inversa sobre
$X$; en otros términos, $\operatorname{Pic}_{Y/S}^\tau$ es la imagen inversa
de $\operatorname{Pic}_{X/S}^\tau$. En particular, para demostrar que el
primer preesquema es de tipo finito sobre $S$, basta demostrarlo para el
segundo, puesto que $\operatorname{Pic}_{Y/S}\longrightarrow
\operatorname{Pic}_{X/S}$ es de tipo finito. Utilizando entonces (i), el lema
de Chow y (iii), se obtiene:

\medskip
\noindent(v) $\operatorname{Pic}_{X/S}^\tau$ es de tipo finito sobre $S$.

En general, la conjunción de (i) para un morfismo finito y de (ii) permite
reducir la mayoría de las cuestiones de finitud al caso en que $X/S$ tiene
fibras geométricas íntegras normales de dimensión $\leq2$; a menudo, aplicando
(i) a un morfismo sobreyectivo no necesariamente finito y la resolución de
singularidades de las superficies algebraicas (debida, en característica
arbitraria, a Abhyankar), se reduce incluso al caso en que $X/S$ es liso y,
por tanto, tiene fibras geométricas no singulares de dimensión 2. Esto permite,
por ejemplo, teniendo en cuenta (v) y la desigualdad de Picard--Igusa que
acota el rango del grupo de Néron--Severi de una superficie proyectiva no
singular, demostrar la siguiente generalización del teorema de finitud de
Néron:

\medskip
\noindent(vi) Sea $X/S$ propio sobre $S$ y por lo demás arbitrario. Entonces
los grupos de Néron--Severi
\[
\operatorname{Pic}_{X_i/k_i}/\operatorname{Pic}_{X_i/k_i}^\circ
\]
de las fibras geométricas $X_i/k_i$ de $X/S$

\src{305}
son de tipo finito, y sus rangos y los órdenes de sus subgrupos de torsión
permanecen acotados.

El mismo método de reducción al caso de superficies no singulares, junto con
los teoremas conocidos para este caso (a saber, el teorema de finitud de Néron
y el hecho de que la forma de intersección sobre el grupo de Néron--Severi es
no degenerada), implica:

\medskip
\noindent(vii) Sea $X$ un esquema propio sobre un cuerpo algebraicamente
cerrado. Existe un número finito de curvas cerradas íntegras $C_i$
($1\leq i\leq r$) en $X$ con la propiedad siguiente: para toda parte $M$ de
$\operatorname{Pic}_{X/k}$, $M$ es cuasicompacta si y solo si los enteros
\[
\deg_{C_i'}\mathcal L\qquad(\mathcal L\in M)
\]
permanecen acotados (donde $C_i'$ designa la normalización de $C_i$).

Aquí se puede tomar como $r$ el rango del grupo de Néron--Severi. Una vez
sabido que este último es de tipo finito, (vii) se reduce al hecho de que las
formas lineales sobre el grupo de Néron--Severi definidas por curvas $C$ en
$X$ solo se anulan simultáneamente sobre los elementos de torsión del grupo de
Néron--Severi. En el caso en que $X$ es proyectivo y no singular, este
resultado, así como (v), se debía a Matsusaka. Utilizando (vii), se obtiene
fácilmente la siguiente caracterización de los Módulos invertibles
$\tau$-equivalentes a cero sobre $X$:

\medskip
\noindent(viii) Sean $X/k$ un esquema propio sobre un cuerpo y $\mathcal L$ un
Módulo invertible sobre $X$. Las condiciones siguientes son equivalentes:

\noindent a. $\mathcal L$ es $\tau$-equivalente a cero;

\noindent b. Para todo Módulo coherente $F$, se tiene
\[
\chi(F\otimes\mathcal L)=\chi(F),
\]
donde $\chi$ designa la característica de Euler--Poincaré;

\noindent b'. Como (b), pero con $F=\mathcal O_Y$, donde $Y$ es un
subesquema cerrado íntegro de dimensión 1 en $X$;

\noindent c. Para todo $Y$ como el anterior, si $Y'$ designa la curva
normalizada, se tiene
\[
\deg\mathcal L_{Y'}=0.
\]

Si $X/k$ es proyectivo y está provisto de un Módulo invertible amplio
$\mathcal O_X(1)$, las condiciones precedentes son también equivalentes a las
siguientes:

\noindent d. Para todo entero $m$, $\mathcal L^{\otimes m}(1)$ es amplio.

\noindent e. (Si $X$ es íntegro.) Para todo par de enteros $m,n$, se tiene
\[
\chi(\mathcal L^{\otimes m}(n))=\chi(\mathcal O_X(n))
\]

\src{306}
(es decir, (b) es cierto tomando $F=\mathcal L^{\otimes m}(n)$).

Para la suficiencia de esta última condición, obsérvese que significa que los
polinomios de Hilbert de los $\mathcal L^{\otimes m}$ son todos iguales; por
consiguiente, en virtud del criterio de Mumford (iii), los
$\mathcal L^{\otimes m}$ permanecen en una parte cuasicompacta de
$\operatorname{Pic}_{X/k}$, es decir, se tiene (a). Las condiciones (b), (b')
y (c) deben considerarse variantes, sobre un esquema propio arbitrario, de la
noción de equivalencia numérica definida habitualmente sobre variedades
proyectivas no singulares. Para estas últimas, la equivalencia de (a) y (c)
era, naturalmente, bien conocida (Matsusaka).

El criterio (e) implica asimismo el resultado siguiente:

\medskip
\noindent(ix) Sea $f:X\longrightarrow S$ un morfismo proyectivo y plano cuyas
fibras geométricas son íntegras. Entonces $\operatorname{Pic}_{X/S}^\tau$ es
abierto y cerrado en $\operatorname{Pic}_{X/S}$.

Nos limitaremos a algunos comentarios sobre la demostración de los resultados
clave (i), (ii) y (iii) (el resultado (iv) se aparta un poco de los demás y se
demuestra utilizando únicamente (i) para los morfismos finitos sobreyectivos
radiciales, concretamente para un morfismo de Frobenius). Para (i), se
utilizan de manera esencial las ideas del descenso no plano (véase en
particular TDTE I, página 9). Resulta que, dado que solo se buscan resultados
de finitud, la falta de criterios de efectividad para datos de descenso es
inofensiva. Por otra parte, Mumford ha demostrado recientemente una forma un
poco más débil de (iii), a saber, el criterio que hace intervenir todos los
coeficientes del polinomio de Hilbert. Su argumento es extremadamente simple
y está inspirado en la demostración de un criterio de amplitud de Nakai
(enunciado por este último para las superficies no singulares y generalizado
por Mumford a morfismos proyectivos arbitrarios). Me parece, sin embargo, que
este argumento solo es válido bajo una ligera restricción adicional sobre las
fibras de $X/S$, la propiedad $(S_2)$ de Serre, que se verifica cuando las
fibras geométricas son normales. Se utiliza entonces este criterio restringido
en la demostración de (ii): el criterio (i) permite, en efecto, reducirse al
caso en que $Y/S$ es plano con fibras geométricas íntegras y normales y,
aplicando el criterio de Mumford, se reduce fácilmente al caso en que $X/S$
satisface las mismas condiciones. La hipótesis de dimensión implica entonces
que las fibras geométricas de $Y$ y de $X$ tienen profundidad $\geq2$ en sus
puntos cerrados, lo que permite aplicar los «criterios de equivalencia» en la
forma que se les da en SGA 1962 (exposición XII) y concluir la demostración de
(ii). Una vez que se dispone de (i) y (ii), no es difícil eliminar del criterio
de Mumford toda hipótesis de normalidad sobre las fibras y demostrarlo también
en la forma más fuerte dada en (iii).

\src{307}
Obsérvese, por último, que la demostración de (i) y (ii) muestra asimismo que,
cuando $S$ es el espectro de un cuerpo $k$, el morfismo
\[
\operatorname{Pic}_{Y/k}\longrightarrow\operatorname{Pic}_{X/k}
\]
es afín (y no solamente de tipo finito).

\def\FgaUnitKey{com}\unit{Comentarios de 1962}
\sid{com}
\src{297}
\seminares{14.º año, 1961/62, Suplemento\\Mayo de 1962}
\paperhead{FUNDAMENTOS DE LA GEOMETRÍA ALGEBRAICA\\COMENTARIOS}%
  {por Alexander GROTHENDIECK}

\section*{Advertencia}
\addcontentsline{toc}{section}{Advertencia}

Entre 1957 y 1962 presentamos en el Seminario Bourbaki ocho exposiciones
relativas a los Fundamentos de la geometría algebraica. Salvo la primera,
estas exposiciones están redactadas en el marco de los esquemas. Todos los
resultados enunciados encontrarán su lugar en los «Éléments de géométrie
algébrique» de Jean DIEUDONNÉ y Alexander GROTHENDIECK. Sin embargo, el
contenido de ninguna de estas exposiciones está cubierto actualmente ni por
alguno de los capítulos ya publicados o en preparación, ni por ningún otro
libro o artículo, y sin duda no lo estará todavía durante algunos años. Esta
es la razón principal que nos ha llevado a reunir estas exposiciones, que
pondrán a disposición de los lectores cierto número de nociones y resultados
clave de la teoría de esquemas mientras se espera una redacción sistemática.
Por otra parte, la lectura de estos textos permitirá familiarizarse
rápidamente con dichos resultados y nociones, sin verse entorpecido por el
detalle, necesariamente prolijo, de un tratado sistemático, y proporcionará
las motivaciones indispensables para el estudio de tal tratado.

Para comodidad del lector, hemos reunido aquí cierto número de comentarios y
erratas, agrupados por exposiciones, que señalarán en particular, de paso, los
progresos realizados desde la redacción del texto e indicarán algunas
referencias adicionales.

Diversos resultados que figuran en estas exposiciones han sido tratados en
detalle en el Seminario de geometría algebraica del IHES, así como en dos
seminarios consecutivos celebrados en Harvard en 1961/62, impartidos el
primero por mí mismo y el segundo por MUMFORD--TATE, para los cuales se están
preparando actualmente notas mimeografiadas (redactadas por el Sr.
LICHTENBAUM).

\section*{Siglas empleadas}
\addcontentsline{toc}{section}{Siglas empleadas}

\begin{description}
\item[SGA :] Seminario de geometría algebraica, Institut des Hautes Études
Scientifiques, 1960--1962.

\src{298}
\item[SHG :] Seminario GROTHENDIECK, Harvard University, 1961/62.

\item[SHMT :] Seminario MUMFORD--TATE, Harvard University, 1961/62.

\item[EGA :] GROTHENDIECK (A.) y DIEUDONNÉ (J.). --- Éléments de géométrie
algébrique. --- París, Presses universitaires de France, 1960, 1961, …
(Institut des Hautes Études Scientifiques, Publications mathématiques, 4, 8,
11, …).

\item[TDTE :] GROTHENDIECK (Alexander). --- Technique de descente et théorèmes
d'existence en géométrie algébrique, I--VI, \emph{Séminaire Bourbaki}, t. 12,
1959/60, n.os 190 y 195; t. 13, 1960/61, n.os 212 y 221; t. 14, 1961/62,
n.os 232 y 236.
\end{description}

\end{document}
